\documentclass[11pt]{article}

\usepackage[margin=1.05in]{geometry}
\usepackage{amsmath,amssymb,amsthm,mathtools,mathrsfs}
\usepackage{enumitem}
\usepackage{booktabs,array,longtable,tabularx}
\usepackage[colorlinks=true,linkcolor=blue,citecolor=blue,urlcolor=blue]{hyperref}
\usepackage[nameinlink,capitalise]{cleveref}
\setlength{\emergencystretch}{3em}
\hbadness=10000
\raggedbottom

% ---------------------------------------------------------------------
% theorem environments
% ---------------------------------------------------------------------
\newtheorem{theorem}{Theorem}[section]
\newtheorem{proposition}[theorem]{Proposition}
\newtheorem{lemma}[theorem]{Lemma}
\newtheorem{corollary}[theorem]{Corollary}
\newtheorem{openproblem}[theorem]{Open Problem}
\theoremstyle{definition}
\newtheorem{definition}[theorem]{Definition}
\newtheorem{assumption}{Assumption}[section]
\newtheorem{construction}[theorem]{Construction}
\newtheorem{example}[theorem]{Example}
\theoremstyle{remark}
\newtheorem{remark}[theorem]{Remark}
\crefname{assumption}{Assumption}{Assumptions}
\Crefname{assumption}{Assumption}{Assumptions}
\crefname{remark}{Remark}{Remarks}
\Crefname{remark}{Remark}{Remarks}
\crefname{example}{Example}{Examples}
\Crefname{example}{Example}{Examples}
\crefname{proposition}{Proposition}{Propositions}
\Crefname{proposition}{Proposition}{Propositions}
\crefname{corollary}{Corollary}{Corollaries}
\Crefname{corollary}{Corollary}{Corollaries}
\crefname{definition}{Definition}{Definitions}
\Crefname{definition}{Definition}{Definitions}
\crefname{lemma}{Lemma}{Lemmas}
\Crefname{lemma}{Lemma}{Lemmas}
\crefname{theorem}{Theorem}{Theorems}
\Crefname{theorem}{Theorem}{Theorems}

% ---------------------------------------------------------------------
% macros
% ---------------------------------------------------------------------
\newcommand{\N}{\mathbb N}
\newcommand{\Z}{\mathbb Z}
\newcommand{\R}{\mathbb R}
\newcommand{\C}{\mathbb C}
\newcommand{\F}{\mathbb F}
\newcommand{\T}{\mathbb T}
\newcommand{\one}{\mathbf 1}
\newcommand{\id}{\mathrm{id}}
\newcommand{\Aut}{\operatorname{Aut}}
\newcommand{\Ad}{\operatorname{Ad}}
\newcommand{\Tr}{\operatorname{Tr}}
\newcommand{\tr}{\operatorname{tr}}
\newcommand{\rank}{\operatorname{rank}}
\newcommand{\rad}{\operatorname{rad}}
\newcommand{\spec}{\operatorname{spec}}
\newcommand{\sgn}{\operatorname{sgn}}
\newcommand{\Span}{\operatorname{span}}
\newcommand{\End}{\operatorname{End}}
\newcommand{\Hom}{\operatorname{Hom}}
\newcommand{\Cl}{\mathrm{Cl}}
\newcommand{\CP}{\mathrm{CP}}
\newcommand{\UCP}{\mathrm{UCP}}
\newcommand{\Alt}{\operatorname{Alt}}
\newcommand{\Sym}{\operatorname{Sym}}
\newcommand{\SO}{\mathrm{SO}}
\newcommand{\norm}[1]{\lVert#1\rVert}
\newcommand{\A}{\mathcal A}
\newcommand{\E}{\mathcal E}
\newcommand{\Hh}{\mathcal H}
\newcommand{\K}{\mathcal K}
\newcommand{\M}{\mathcal M}
\newcommand{\Pp}{\mathfrak P}
\newcommand{\Qq}{\mathsf Q}
\newcommand{\Rr}{\mathcal R}
\newcommand{\Ww}{\mathcal W}
\newcommand{\Gg}{\mathcal G}
\newcommand{\D}{\mathsf D}
\newcommand{\Ham}{\mathsf H}
\newcommand{\J}{\mathsf J}
\newcommand{\GammaR}{\Gamma_{\rho}}
\newcommand{\sharpP}{\sharp_{\rho}}
\newcommand{\Ztwo}{\mathbb Z/2}
\newcommand{\Zfour}{\mathbb Z/4}
\newcommand{\eps}{\varepsilon}
\newcommand{\dd}{\mathrm d}
\newcommand{\op}{\mathrm{op}}
\newcommand{\rev}{\mathrm{rev}}
\newcommand{\pred}{\mathrm{pred}}
\newcommand{\rec}{\mathrm{rec}}
\newcommand{\spat}{\mathrm{sp}}
\newcommand{\HS}{\mathrm{HS}}
\newcommand{\Pred}{\operatorname{Pred}}

\title{\textbf{From Operational Prediction to Signed Renewal Memory}}
\author{A.~P\'elissier}
\date{}

% ATMP strengthened revision V6: retains the V5 stability and no-go results and
% additionally proves that no natural history-valued selector is determined by
% the predictive quotient, gives an explicit rooted-orbit tomography
% counterexample, derives a canonical resolved instrument from a specified
% interaction-complete stable-pointer phase, and adds a necessity/no-go ledger.

\begin{document}
\maketitle

\begin{abstract}
We investigate whether the mathematical structure used to describe emergent Lorentzian spacetime can itself arise from a more basic principle: prediction. We begin with composable operations and identify histories whenever no admissible future experiment can distinguish them. This removes redundant microscopic information and produces a minimal memory containing exactly what is needed to predict future outcomes. Under explicit operational assumptions, this reduced memory acquires the structure of finite-dimensional complex quantum dynamics and remains stable under small errors in predictive compression.
%
We then show how resolved outcomes generate well-defined successor states and multiplicative path weights. Forward and retrodictive dynamics define a real measure of irreversibility, while stable-record transport produces an independent binary orientation class and a two-sheeted signed cover. Reversible transformations generate a Clifford algebra, pressure selects an intrinsic renewal scale, and spatial record statistics provide the remaining finite ingredients for a companion construction of Lorentzian geometry.
%
The result is conditional: it identifies clear assumptions and phase conditions under which spacetime-like structure emerges, while allowing unoriented, subcritical, and torsional alternatives.
%
A \texttt{Lean} formalization is available at
\url{https://github.com/AI-MathPhys/renewal-geometry}.
\end{abstract}



\setcounter{tocdepth}{2}
\tableofcontents

% =====================================================================
\section{Introduction}
\label{sec:introduction}

A companion work, \emph{Renewal Spectral Geometry and the Emergence of
Lorentzian Spacetime} \cite{PelissierLorentzian2026}, starts from a finite
family of unital completely positive maps and studies the geometry encoded by
histories of their composition.  Under explicit regularity and phase
conditions, that construction produces a positive spectral geometry, a
minimal signed cover with a Krein fundamental symmetry, a primitive Clifford
factor generated by reversible predictive transformations, and controlled
flat and curved Lorentzian Dirac limits.  Spacetime is therefore not assumed
there, but one substantial object still is: the signed renewal memory from
which the geometry is built.

The purpose of the present paper is to reconstruct that object from a more
primitive operational level and to identify, with sharp logical boundaries,
which later ingredients are canonical and which remain phase selections.  The
central question is not simply whether quantum channels can be used to model a
pregeometric dynamics.  It is why completely positive renewal maps, a
predictive quotient, a retrodictive comparison, a signed orientation line, a
modular scale, and a torsion-free translational sector should arise together
at all.  These structures are often packaged as if they had a common origin,
but they do not.  Positivity is weaker than complete positivity; a
retrodictive channel is not an inverse; nonequilibrium affinity is real-valued
whereas orientation is mod two; a nontrivial double cover does not select one
of its sheets; stationarity need not imply exchange symmetry; and detailed
balance need not remove torsion.  A continuum limit, moreover, requires
analytic regularity that should not be disguised as a foundational axiom.

Our strategy is reconstructive.  The primitive language contains composable
operations, parallel composition, terminal tests, and real weights assigned
to closed experiments.  Histories are generated composites rather than
trajectories in a pre-existing time.  Two rooted histories are identified
when every admissible continuation and terminal test gives the same result.
This future-equivalence relation produces a universal reduced predictive
memory: every other future-sufficient representation factors through it.  The
construction belongs to the broad operational and process-theoretic programme
for reconstructing quantum theory
\cite{Hardy2001,ChiribellaDAPS2009,ChiribellaDArianoPerinotti2011,MasanesMuller2011,DakicBrukner2011,DArianoChiribellaPerinotti2017,SelbyScandoloCoecke2021,vandeWetering2019},
but the organising object here is a quotient of histories rather than a
single-shot state space.  In particular, ordinary operational equivalence is
compared with its complete ancillary extension.  Their equality is isolated
as a genuine no-holism condition, rather than being hidden inside a local-tomography axiom; the obstruction
exhibited by real-vector-space quantum theory is correspondingly retained
\cite{BarnumWilce2014,HardyWootters2012,delaTorre2012}.

The first reconstruction theorem passes from this predictive quotient to
finite quantum dynamics.  A reduced sharp-reversible realisation supplies the
purification and sharpness structure needed for operational diagonalisation
and the Euclidean Jordan route
\cite{ChiribellaScandolo2015,BarnumLeeScandoloSelby2017,BarnumUdudecVdW2023}.
Minimal continuous two-path interference then selects complex matrix type,
while a compatible phase response fixes the associative orientation; this is
the point at which the Jordan structure becomes a complex $C^*$-algebra
\cite{HancheOlsen1985}.  Stability under arbitrary finite ancillas upgrades
positive deterministic transformations to UCP maps.  The paper also proves a
quantitative perturbation theorem: a nearly idempotent, nearly
future-sufficient compression with the correct rank and a transversality bound
lies near an exact minimal predictive projection, and the induced compressed
renewal dynamics is stable.  Thus the finite UCP renewal alphabet is derived
from operational prediction rather than postulated as primitive geometry.

A second reconstruction concerns resolved histories.  Computational mechanics
identifies stochastic pasts that induce the same conditional future law, and
the resulting causal states form a minimal unifilar predictive model
\cite{CrutchfieldYoung1989,ShaliziCrutchfield2001}; quantum memory advantages
for such models have also been studied
\cite{GuWiesnerRieperVedral2012,BinderThompsonGu2018}.  Here the same principle
is lifted to resolved completely positive branches.  Equality of conditional
future profiles is proved to be stable under appending an observed outcome,
so the conditional predictive quotient has a deterministic outcome-labelled
successor map.  In every reduced future-tomographic operator realisation, the
normalised output of a branch is therefore forced to equal the state
representing that successor class.  This yields the branchwise
reference-renewal identity and multiplicative scalar path weights as theorems,
not hidden-state assumptions.  A canonical successor refinement repairs
under-resolved target graphs; Kraus changes within a fixed resolved outcome
leave the weights invariant; and approximate future tomography gives an
explicit error bound along every finite history.

Resolved outcomes themselves cannot, however, be recovered from an averaged
channel without additional physical information.  To make this boundary
precise, we give an optional microscopic bridge.  A minimal dilation together
with a specified family of system--environment interactions determines a
basis-independent environment coefficient algebra.  In an
interaction-complete monitoring phase, the associated dephasing semigroup
converges to the trace-preserving conditional expectation onto its commutant,
and the centre of the stable nondemolition algebra yields a canonical pointer
instrument.  This construction is related to einselection and decoherence
\cite{Zurek1981,Zurek2003,Schlosshauer2004}, conditional expectations and
fixed-point algebras
\cite{BlanchardOlkiewicz2003,Evans1977,FrigerioVerri1982,FagnolaSassoUmanita2019,TicozziViola2008},
and the redundancy viewpoint of quantum Darwinism
\cite{BlumeKohoutZurek2006}.  Its additional conclusion is a canonicity
statement relative to the specified interaction family: the stable pointer
centre is fixed before the conditional predictive quotient is taken, while an
averaged channel alone remains insufficient to select an instrument.

The stationary and retrodictive layers are reconstructed next.  A primitive
weighted transfer selects a unique faithful stationary state.  When the
elementary channels preserve that state, the symmetric KMS pairing determines
their common completely positive retrodictive adjoints, in the lineage of the
Accardi--Cecchini and Petz transpose maps
\cite{AccardiCecchini1982,Petz1986,Petz2003} and KMS formulations of quantum
detailed balance
\cite{Alicki1976,KFGV1977,GoldsteinLindsay1995,FagnolaUmanita2010}.  In the
commutative reduction, a stationary forward path has exactly the same weight
as the reversed path under the canonical Petz chain.  Consequently a genuine
irreversibility observable must compare forward and retrodictive branch laws
in the same path orientation, consistently with time-reversal and Bayesian
readings of fluctuation relations \cite{Crooks2008,BuscemiScarani2021}.  The
resulting cycle affinity belongs to the established network-thermodynamic
lineage of Schnakenberg and its extensions
\cite{Schnakenberg1976,LebowitzSpohn1999,AndrieuxGaspard2007}.

The Lorentzian sign has a different origin.  Conditional future prediction
reconstructs the minimal algebra of stable resolved records, and determinant
transport on its contrast module defines a real line local system and a class
in $H^1(G;\Ztwo)$; the determinant-line viewpoint is standard in its natural
geometric setting \cite{ZingerDeterminant2016}.  We extend the construction
from reversible record permutations to positive noisy correspondences whose
action on future-distinct contrasts remains invertible.  The key result is a
strict separation theorem: the real affinity class in $H^1(G;\R)$ and the
record-orientation class in $H^1(G;\Ztwo)$ are logically and cohomologically
independent.  There is no coefficient homomorphism that reduces the former to
the latter, and all four combinations of vanishing or nonvanishing affinity
and orientation occur.  To our knowledge, this theorem-level separation of
network affinity from stable-record parity has not previously been formulated.
It shows in particular that entropy production does not generate the signed
cover and that orientability does not imply detailed balance.

The remaining steps are phase-theoretic.  On a selected recurrent component,
reversible predictive units act projectively on a finite factor.  Their
commutator bicharacter and radical quotient produce a primitive complex
Clifford algebra, while the signed cover fixes its temporal row.  Reciprocal
branch capacities admit a unique decomposition into symmetric conductance and
half-affinity.  Perron--Frobenius pressure then selects a dimensionless renewal
exponent, in the finite-graph thermodynamic-formalism tradition
\cite{Bowen1975,Ruelle2004,ParryPollicott1990}.  Its stationary law obeys the
exact identity
\[
  \beta\mu_\ell
  =h_{\mathrm{rate}}+\mathbb E[\log c_e]+\frac12\sigma,
\]
which ties modular depth, entropy rate, symmetric capacity, and entropy
production in one formula.  A graded path cover realises the corresponding
affine modular clock, giving a renewal implementation of the thermal-time
relation between modular flow and physical time \cite{ConnesRovelli1994} while
avoiding the obstruction to a nonzero common exponent on a finite return
cycle.

Spatial and Cartan data are obtained from record statistics rather than added
as independent geometry.  Covariance of positive spatial second moments
derives orthogonal comparison between neighbouring frames.  At the
plaquette level, an exchange-covariant primitive law on microscopic two-step
diamonds forces the exchange-odd translational defect to have zero stationary
mean while leaving the exchange-even rotational defect unconstrained.  The
same condition is the zero-source Euler--Lagrange equation of a convex
exchange pressure.  Since torsion is precisely the translational obstruction
separating Levi--Civita from general metric-compatible transport
\cite{Hehl1976}, this provides the finite input needed by the downstream
Cartan reconstruction.  An explicit Ising-type countermodel shows at the same
time that one-step detailed balance and stationarity can support both
exchange-symmetric torsion-free states and symmetry-broken torsional states.
Torsion freedom is therefore a genuine phase selection rather than an
automatic consequence of reversibility.

The main contributions may be summarised as follows.
\begin{enumerate}[label=\textup{(\roman*)},leftmargin=2.5em]
\item We construct the universal rooted and complete predictive quotients,
      reconstruct a finite complex $C^*$-algebra and UCP renewal dynamics from
      explicit operational assumptions, and prove quantitative stability of
      approximate predictive compression.
\item We derive conditional unifilarity, branchwise reference renewal, and
      multiplicative path weights from the predictive quotient and future
      tomography, including canonical refinement and finite-error control.
\item Relative to a specified interaction family, we extract a canonical
      stable pointer instrument from the centre of the nondemolition
      commutant, while proving why no averaged channel can determine such an
      instrument by itself.
\item We separate the real affinity class from the mod-two record-orientation
      class, extend determinant orientation to noisy regular record transport,
      and classify the resulting signed cover and its amplitude lifts.
\item We derive the entropy--affinity--depth identity, the pressure-selected
      affine clock, orthogonal frame comparison, and exchange-selected
      translational closure, and we provide explicit countermodels showing
      which conclusions are phase dependent.
\end{enumerate}

These statements are conditional reconstruction results, not a claim that
every operational theory becomes Lorentzian spacetime.  The predictive
quotient, stationary reference, KMS--Petz dual, record algebra, determinant
line, projective multiplier, and pressure root are canonical once their
respective domains and hypotheses are present.  By contrast, nontrivial sign
holonomy, sufficient reversible rank, selection of one deck orientation,
pressure supercriticality, spatial ellipticity, exchange-symmetric
translation, and continuum regularity are genuine phase or scaling
conditions.  Their failure produces mathematically consistent holistic,
unsigned, subcritical, unoriented, spatially degenerate, or torsional regimes.
The aim is therefore not to replace one large geometric axiom by an equally
opaque operational list, but to expose a short auditable dependency chain in
which every indispensable premise has a distinct role and a distinct failure
mode.


% =====================================================================
\section{Predictive reduction and the emergence of UCP renewal memory}
\label{sec:predictive}

The first task is to construct the object whose geometry is studied
downstream.  We begin before vector spaces, probabilities, or operator
algebras have been chosen.  The primitive relation is operational
indistinguishability under all admissible futures.

The argument proceeds in three conceptual stages.  First, future equivalence is
shown to define a canonical minimal memory, both for rooted histories and for
processes placed in arbitrary ancillary contexts.  Second, the exact quotient is
related to practical finite-dimensional compressions and is proved stable under
controlled approximation.  Third, additional operational hypotheses reconstruct
the algebraic form of the memory and identify its deterministic updates as UCP
maps.  The final subsection then supplies the tomographic and stationary
interfaces needed for the retrodictive constructions of \Cref{sec:retrodictive}.

\subsection{Rooted predictive memory and process equivalence}
\label{subsec:rooted-and-process}

\begin{definition}[Finite operational process system]
\label{def:process-system}
A finite operational process system consists of:
\begin{enumerate}[label=\textup{(\roman*)},leftmargin=2.4em]
\item a small strict symmetric monoidal category $\mathsf C$ of typed processes,
      with tensor unit $I$;
\item for every object $A$, a deterministic discard effect
      $u_A:A\to I$, satisfying $u_{A\otimes B}=u_A\otimes u_B$;
\item a finite generating family $\mathsf G$ of processes, together with the
      identities and symmetries needed for typed sequential and parallel
      composition;
\item convex mixing of parallel processes and an evaluation
      $\Pr(\omega)\in[0,1]$ for every closed experiment $\omega:I\to I$;
\item specified admissible continuations and terminal effects, closed under
      composition, classical postprocessing, and tensoring with identities.
\end{enumerate}
A \emph{rooted history ending at $B$} is a generated process $h:I\to B$.
The word ``future'' means admissible postcomposition, not a pre-existing
physical time.
\end{definition}

\begin{definition}[Rooted future equivalence]
\label{def:future-equivalence}
Let $h,h':I\to B$ be rooted histories.  Write $h\sim_+h'$ when
\[
 \Pr(e\circ r\circ h)=\Pr(e\circ r\circ h')
\]
for every admissible continuation $r:B\to C$ and terminal effect $e:C\to I$.
Write
\[
 \Qq_+(B):=\{h:I\to B\text{ generated}\}/\!\sim_+.
\]
\end{definition}

\begin{proposition}[Rooted future congruence]
\label{prop:largest-congruence}
For each endpoint $B$, the relation $\sim_+$ is an equivalence relation.  If
$h\sim_+h'$ and $k:B\to C$ is an admissible generated continuation, then
\[
 k\circ h\sim_+k\circ h'.
\]
Consequently every continuation induces a well-defined map
\[
 \Qq_+(k):\Qq_+(B)\longrightarrow\Qq_+(C),\qquad
 \Qq_+(k)[h]=[k\circ h],
\]
and $\Qq_+$ is a covariant functor from the generated continuation category to
$\mathbf{Set}$ on the reachable component.
\end{proposition}

\begin{proof}
Associate to $h:I\to B$ its future-profile function
\[
 \Pi_B(h)(r,e):=\Pr(e\circ r\circ h).
\]
Then $h\sim_+h'$ exactly when $\Pi_B(h)=\Pi_B(h')$, proving reflexivity,
symmetry, and transitivity.  If $h\sim_+h'$ and $k:B\to C$, then for every
$d:C\to D$ and $e:D\to I$,
\[
 \Pr(e\circ d\circ k\circ h)
 =\Pr(e\circ(d\circ k)\circ h)
 =\Pr(e\circ(d\circ k)\circ h')
 =\Pr(e\circ d\circ k\circ h'),
\]
because $d\circ k$ is an admissible future continuation from $B$.  Thus
postcomposition descends.  The identity and composition laws follow from those
of $\mathsf C$.
\end{proof}

Rooted future equivalence is intentionally one-sided.  It is the correct object
for predictive memory, but it is not by itself a congruence on all hom-sets.
General process equality must also quantify over the admissible past.

The preceding proposition establishes the essential predictive fact: once two
rooted histories have the same future profile, no later admissible operation can
separate them.  This gives a well-defined state-like memory.  To compare general
processes, however, one must also allow different admissible preparations before
the process.  The next definition therefore enlarges the comparison from rooted
histories to arbitrary input--output transformations.

\begin{definition}[Ordinary two-sided operational equivalence]
\label{def:two-sided-equivalence}
For generated processes $p,q:A\to B$, write $p\equiv_1q$ when
\[
 \Pr(e\circ r\circ p\circ a)
 =\Pr(e\circ r\circ q\circ a)
\]
for every generated preparation $a:I\to A$, every admissible continuation
$r:B\to C$, and every terminal effect $e:C\to I$.
\end{definition}

\begin{proposition}[Two-sided equivalence is a sequential congruence]
\label{prop:two-sided-congruence}
The relation $\equiv_1$ is stable under arbitrary typed precomposition and
postcomposition.  Hence the ordinary process quotients
\[
 \Qq_1(A,B):=\mathsf C_{\mathsf G}(A,B)/{\equiv_1}
\]
form a category under sequential composition.
\end{proposition}

\begin{proof}
Let $p\equiv_1q:A\to B$, $u:A'\to A$, and $v:B\to B'$.  Every rooted
preparation before $p\circ u$ has the form $u\circ a'$ for an admissible
$a':I\to A'$, and every future continuation after $v\circ p$ is absorbed into
$r\circ v$.  The defining equality for $p\equiv_1q$ therefore gives
$p\circ u\equiv_1q\circ u$ and $v\circ p\equiv_1v\circ q$.  Composition of
classes is consequently well defined.
\end{proof}

\begin{definition}[Complete contextual equivalence]
\label{def:complete-equivalence}
For $p,q:A\to B$, write $p\equiv_{\parallel}q$ when, for every ancillary object
$R$, every joint preparation $\omega:I\to A\otimes R$, every continuation
$c:B\otimes R\to C$, and every terminal effect $e:C\to I$,
\[
 \Pr\!\left(e\circ c\circ(p\otimes\id_R)\circ\omega\right)
 =
 \Pr\!\left(e\circ c\circ(q\otimes\id_R)\circ\omega\right).
\]
Set
\[
 \Qq_{\parallel}(A,B)
 :=\mathsf C_{\mathsf G}(A,B)/{\equiv_{\parallel}}.
\]
\end{definition}

\begin{theorem}[Canonical parallel completion]
\label{thm:parallel-completion}
Complete contextual equivalence is a monoidal congruence and implies ordinary
operational equivalence.  Moreover it is the largest monoidal congruence
contained in $\equiv_1$: if $\mathcal R$ is an equivalence relation on every
process set such that
\begin{enumerate}[label=\textup{(\alph*)},leftmargin=2.4em]
\item $p\,\mathcal R\,q$ implies $p\equiv_1q$;
\item $\mathcal R$ is stable under sequential composition;
\item $\mathcal R$ is stable under tensoring with arbitrary transformations,
\end{enumerate}
then $p\,\mathcal R\,q$ implies $p\equiv_{\parallel}q$.
\end{theorem}

\begin{proof}
Equivalence and sequential stability follow by absorbing pre- and
postprocessings into the joint preparation and terminal continuation.  Tensor
stability follows by regarding any further system as part of the ancillary
object.  Taking the trivial ancilla proves
$\equiv_{\parallel}\subseteq\equiv_1$.

Now suppose $p\,\mathcal R\,q$.  Tensor stability gives
$p\otimes\id_R\,\mathcal R\,q\otimes\id_R$ for every $R$.  By (a), the two
extended processes are ordinarily equivalent on the composite input and output,
so they have identical values on every joint preparation and every joint
terminal test.  This is exactly $p\equiv_{\parallel}q$.  Since
$\equiv_{\parallel}$ itself has properties (a)--(c), it is the largest such
monoidal congruence.
\end{proof}

The theorem identifies complete contextual equivalence as the correct process
notion whenever parallel composition is physically meaningful.  Ordinary
experiments may miss distinctions that become visible only through an ancilla;
the comparison map below measures exactly that possible loss of information.

Because $\equiv_{\parallel}$ is finer than $\equiv_1$, there is a canonical
surjection
\[
 \kappa_{A,B}:\Qq_{\parallel}(A,B)\longrightarrow\Qq_1(A,B),
 \qquad [p]_{\parallel}\longmapsto[p]_1.
\]
Its fibres are the process distinctions invisible to all unassisted
preparation--test experiments but visible in some ancillary context.

\begin{definition}[Parallel completeness]
\label{def:parallel-completeness}
The operational theory is \emph{parallel complete} when every $\kappa_{A,B}$ is
bijective, equivalently when
\[
 \equiv_1=\equiv_{\parallel}
\]
on every process set.  This is the precise no-holism condition.
\end{definition}

\begin{theorem}[Universal complete predictive quotient]
\label{thm:universal-quotient}
The quotient
\[
 \Qq_{\pred}:=\mathsf C_{\mathsf G}/\!\equiv_{\parallel}
\]
is a symmetric monoidal category.  Its quotient functor is universal among
symmetric monoidal functors that are constant on complete contextual
classes.  Under parallel completeness it is canonically isomorphic to the
ordinary two-sided operational quotient $\mathsf C_{\mathsf G}/\!\equiv_1$.
\end{theorem}

\begin{proof}
The first assertion follows because $\equiv_{\parallel}$ is a monoidal
congruence.  For the universal property, if
$F:\mathsf C_{\mathsf G}\to\mathsf D$ is symmetric monoidal and
$p\equiv_{\parallel}q$ implies $F(p)=F(q)$, define
$\overline F([p]_{\parallel})=F(p)$.  This is well defined, symmetric monoidal,
and unique because the quotient functor is surjective on morphism classes.  If
parallel completeness holds, the two kernel relations coincide, giving the
canonical quotient isomorphism.
\end{proof}

\begin{proposition}[Real quantum theory has a strict parallel completion]
\label{prop:rebit-parallel-counterexample}
In finite-dimensional real quantum theory there are channels $T_+,T_-$ such
that $T_+\equiv_1T_-$ but $T_+\not\equiv_{\parallel}T_-$.  Thus ordinary
process reduction need not be monoidal even in a theory with purification and
sharp state--effect duality.
\end{proposition}

\begin{proof}
Let
\[
 Y_{\R}:=\begin{pmatrix}0&-1\\1&0\end{pmatrix},
 \qquad Y_{\R}^{\mathsf T}=-Y_{\R},
 \qquad Y_{\R}^2=-I_2,
\]
and define two real two-rebit density matrices
\[
 \rho_\pm:=\frac14\bigl(I_2\otimes I_2
              \pm Y_{\R}\otimes Y_{\R}\bigr).
\]
The matrix $Y_{\R}\otimes Y_{\R}$ is real symmetric and squares to the
identity, so $\rho_\pm\ge0$ and $\Tr\rho_\pm=1$.  Their first marginal is
$I_2/2$ because $\Tr Y_{\R}=0$.  Hence the real Choi correspondence associates
to them trace-preserving completely positive maps $T_\pm$.

For real symmetric local effects $A,B$,
\[
 \Tr\!\left[(A\otimes B)(\rho_+-\rho_-)\right]
 =\frac12\Tr(AY_{\R})\Tr(BY_{\R})=0,
\]
because the trace pairing of a symmetric and an antisymmetric real matrix
vanishes.  Thus the two Choi states have identical product-effect statistics,
and the corresponding channels are indistinguishable by every unassisted
preparation--effect experiment.  On the other hand
\[
 P_+:=\frac12\bigl(I_4+Y_{\R}\otimes Y_{\R}\bigr)
\]
is a valid real-symmetric projection and
$\Tr(P_+\rho_+)=1$, $\Tr(P_+\rho_-)=0$.  A joint ancillary experiment therefore
distinguishes the channels.  This proves the strictness of the comparison map
$\kappa$; compare the general real-quantum obstruction in
\cite{BarnumWilce2014}.
\end{proof}

\begin{remark}[Why no-holism is load-bearing]
\label{rem:no-holism}
Parallel completeness does not follow from sequential prediction or
purification.  Finite-dimensional real quantum theory contains distinct
composite states with identical product-effect statistics, and the associated
faithful Choi processes are ordinarily indistinguishable but ancilla-distinct;
see \cite{BarnumWilce2014,ChiribellaProcessTomography2021}.  Thus a faithful
local-channel representation of all process distinctions requires either the
explicit equality $\Qq_{\parallel}\cong\Qq_1$ or a richer carrier retaining the
holistic fibres of $\kappa$.
\end{remark}

Taken together, the quotient theorems and the real-quantum counterexample fix the
logical starting point for renewal.  Prediction canonically determines an
equivalence class, but monoidal consistency may require the complete rather than
the ordinary quotient.  Having identified the information that must be retained,
we next ask how that information can be represented as a minimal repeatedly
updated memory.

\subsection{Renewal as exact predictive compression}

Fix an endpoint $B$ and write $\mathsf H_B$ for the generated rooted histories
$I\to B$.  The canonical quotient map is
\[
 q_B:\mathsf H_B\twoheadrightarrow\Qq_+(B).
\]

The purpose of this subsection is not to choose a preferred microscopic history.
It is to show that every future-sufficient representation contains the same
reduced predictive content, and that any concrete renewal rule is merely a choice
of representatives for that canonical content.

\begin{definition}[Future-sufficient realization]
\label{def:future-sufficient}
A future-sufficient realization of the rooted histories consists of sets
$X(B)$, surjective encodings $\eta_B:\mathsf H_B\to X(B)$, transition maps
$\delta_k:X(B)\to X(C)$ for generated continuations $k:B\to C$, and output
functions reproducing every terminal-test weight, such that
\[
 \delta_k\eta_B(h)=\eta_C(k\circ h).
\]
It is \emph{reduced} when distinct elements of $X(B)$ are separated by some
future continuation and terminal effect.
\end{definition}

\begin{theorem}[Minimality and uniqueness of predictive memory]
\label{thm:minimal-predictive-memory}
The rooted quotient $\Qq_+$ with the transition maps of
\Cref{prop:largest-congruence} is future sufficient and reduced.  For every
future-sufficient realization $X$ there is a unique surjection
\[
 \pi_B^X:X(B)\longrightarrow\Qq_+(B),
 \qquad
 \pi_B^X(\eta_B(h))=[h],
\]
intertwining all transitions and outputs.  The realization $X$ is reduced if
and only if every $\pi_B^X$ is bijective.  Hence $\Qq_+$ is the unique reduced
future-sufficient memory, up to the unique isomorphism fixing all encoded
histories.
\end{theorem}

\begin{proof}
If $\eta_B(h)=\eta_B(h')$, transition compatibility and equality of all output
functions imply that every future experiment has the same value on $h$ and
$h'$.  Thus $h\sim_+h'$, proving that the displayed formula for $\pi_B^X$ is
well defined.  It is surjective because every class has a representative.
Compatibility with continuations follows from
\[
 \pi_C^X(\delta_k\eta_B(h))
 =\pi_C^X(\eta_C(k\circ h))
 =[k\circ h]
 =\Qq_+(k)[h].
\]
The quotient itself is sufficient by definition and is reduced because
inequivalent classes differ on some future profile.

If $X$ is reduced and $\pi_B^X(x)=\pi_B^X(y)$, choose histories representing
$x,y$.  Their quotient classes agree, so all future outputs agree; reducedness
gives $x=y$.  Conversely, if $\pi_B^X$ is injective and two realized states have
identical future outputs, their representing histories are future equivalent,
so their images under $\pi_B^X$ agree and injectivity gives equality.  Thus
reducedness is equivalent to bijectivity.  Uniqueness follows from surjectivity
of the encodings.
\end{proof}

The theorem is the minimal-memory statement on which the rest of the paper rests.
It says that predictive states are not introduced as hidden variables: they are
forced by the equivalence relation, and every other sufficient memory maps onto
them.  The corollary translates this abstract quotient into the familiar language
of an idempotent renewal operation on raw histories.

\begin{corollary}[Exact renewal retract]
\label{cor:renewal-retract}
Suppose a type-preserving section
$s_B:\Qq_+(B)\to\mathsf H_B$ is chosen for every reachable endpoint.  Then
\[
 \mathsf R_B:=s_Bq_B
\]
is idempotent, satisfies $q_B\mathsf R_B=q_B$, and replaces every history by a
future-equivalent representative.  Its image contains exactly one representative
of every predictive class.  Conversely every idempotent map with these
properties is obtained from a unique section.  The induced update
\[
 \Phi_k^{(s)}:=\mathsf R_C\circ(k\circ-)\big|_{\operatorname{im}\mathsf R_B}
\]
forms an exact action of the continuation category and is conjugate to the
canonical action on $\Qq_+$.
\end{corollary}

\begin{proof}
Since $q_Bs_B=\id$,
$\mathsf R_B^2=s_Bq_Bs_Bq_B=\mathsf R_B$ and
$q_B\mathsf R_B=q_B$.  Conversely, an idempotent with one fixed point in each
class defines the inverse of $q_B$ restricted to its image.  For updates,
\[
 q_C\Phi_k^{(s)}=\Qq_+(k)q_B.
\]
The restricted quotient maps are bijective, so identity and composition of the
canonical transition system imply the same laws for $\Phi^{(s)}$.
\end{proof}

The quotient is canonical; a section is not.  Renewal therefore means repeated
return to minimal sufficient predictive content, not the selection of a
preferred microscopic past.


\begin{proposition}[No natural history-valued selector]
\label{prop:no-natural-selector}
Let \(q_B:\mathsf H_B\twoheadrightarrow\Qq_+(B)\) be a rooted predictive
quotient.  Suppose one quotient class contains two distinct histories
\(h_1,h_2\), and suppose there is a bijection
\[
 \alpha_B:\mathsf H_B\longrightarrow\mathsf H_B
\]
which preserves endpoint type and every future profile, satisfies
\(q_B\alpha_B=q_B\), and exchanges \(h_1\) and \(h_2\).  Then no section
\(s_B:\Qq_+(B)\to\mathsf H_B\) can be invariant under all
future-profile-preserving automorphisms of the presentation.  Consequently
the quotient data alone do not determine a natural exact renewal retract on
raw histories.
\end{proposition}

\begin{proof}
Let \(x=q_B(h_1)=q_B(h_2)\).  If an invariant section existed, then
\(\alpha_Bs_B=s_B\).  Since \(s_B(x)\) is one element of the fibre
\(q_B^{-1}(x)\), invariance would require it to be fixed by every
profile-preserving automorphism.  In particular, if \(s_B(x)=h_1\), then
\[
 s_B(x)=\alpha_Bs_B(x)=\alpha_B(h_1)=h_2,
\]
contradicting \(h_1\ne h_2\); the same contradiction occurs if
\(s_B(x)=h_2\).  More generally, any transitive automorphism orbit in a fibre
contains no invariant choice of one representative.  Since an exact minimal
history-valued retract is equivalent to a section by
\Cref{cor:renewal-retract}, no such retract is naturally determined by the
quotient alone.
\end{proof}

\begin{remark}[Representative gauge]
The obstruction is a naturality statement, not a failure of existence.
Finite quotients admit sections, and arbitrary set-sized quotients admit them
under the axiom of choice.  What prediction fixes canonically is the quotient
and its induced transition action; choosing a representative history requires
extra structure, such as a normal form, metric projection, conditional
expectation, or entropy criterion.
\end{remark}

The exact theory therefore separates two levels.  The quotient and its transition
action are canonical, whereas a history-valued section is a gauge choice and in
general cannot be selected naturally.  Real implementations will moreover only
approximate the exact quotient.  The next subsection shows that this is not fatal:
under explicit nondegeneracy conditions, an approximate compression determines a
nearby exact minimal predictive memory.

\subsection{Finite-dimensional stability of predictive compression}
\label{subsec:predictive-compression-stability}

The quotient construction above is exact and set theoretic.  A finite
implementation may instead provide a linear compression that is only nearly
idempotent and only nearly preserves all predictive profiles.  We now show
that, below the sharp spectral threshold and under explicit rank and
transversality conditions, such a compression lies near an exact minimal
predictive retract.  No positivity is required in this subsection.

The strategy is deliberately modular.  We first classify exact linear retracts,
then turn a nearly idempotent map into a genuine projection, and finally correct
that projection so that it preserves the predictive quotient exactly.  This
separates algebraic idempotence from the independent requirements of sufficiency
and minimal rank.

\begin{definition}[Finite linear predictive presentation]
\label{def:linear-predictive-presentation}
A finite linear predictive presentation is a surjective linear map
\[
 q:X\twoheadrightarrow Q
\]
between finite-dimensional normed spaces.  The quotient \(Q\) is the linear
span of the retained predictive profiles, and
\[
 K_{\pred}:=\ker q,\qquad r:=\dim Q
\]
is the space of observationally null directions.  A linear map \(P:X\to X\)
is an \emph{exact predictive projection} when
\[
 P^2=P,\qquad qP=q.
\]
It is \emph{minimal} when \(\rank P=r\).
\end{definition}

\begin{proposition}[Classification of exact linear predictive retracts]
\label{prop:linear-predictive-retracts}
Let \(P\) be an exact predictive projection for
\(q:X\twoheadrightarrow Q\).  Then \(\rank P\ge r\) and
\(\ker P\subseteq K_{\pred}\).  Moreover the following are equivalent:
\[
 \rank P=r,\qquad
 \ker P=K_{\pred},\qquad
 q|_{\operatorname{im}P}:\operatorname{im}P\longrightarrow Q
 \text{ is an isomorphism}.
\]
Consequently minimal exact predictive projections are in bijection with
linear complements \(M\) of \(K_{\pred}\):
\[
 X=M\oplus K_{\pred},\qquad
 P_M=(q|_M)^{-1}q.
\]
Any two such carriers are related by the unique quotient-fixed linear
isomorphism.
\end{proposition}

\begin{proof}
From \(qP=q\) one obtains \(\rank q\le\rank P\).  If \(Px=0\), then
\(qx=qPx=0\), proving \(\ker P\subseteq K_{\pred}\).  Since \(P\) is a
projection,
\[
 \dim\ker P=\dim X-\rank P,\qquad
 \dim K_{\pred}=\dim X-r.
\]
Thus \(\rank P=r\) is equivalent to \(\ker P=K_{\pred}\).  The restriction
\(q|_{\operatorname{im}P}\) is surjective because \(qP=q\); it is injective
exactly when \(\operatorname{im}P\cap K_{\pred}=0\), which is equivalent to
the same dimension condition.

If \(M\) complements \(K_{\pred}\), then \(q|_M\) is an isomorphism and
\(P_M=(q|_M)^{-1}q\) is an idempotent with \(qP_M=q\), image \(M\), and
kernel \(K_{\pred}\).  Conversely the image of every minimal exact
projection is such a complement.  For two complements \(M,N\), the map
\[
 U_{M,N}:=(q|_N)^{-1}q|_M:M\longrightarrow N
\]
is the unique isomorphism satisfying \(qU_{M,N}=q|_M\).
\end{proof}

\begin{lemma}[Idempotentization below the quarter threshold]
\label{lem:predictive-idempotentization}
Let \(E\in\End(X)\) and assume
\[
 \delta:=\norm{E^2-E}<\frac14.
\]
Set
\[
 D:=E^2-E,\qquad A:=2E-\one,
\]
and define, by the norm-convergent binomial series,
\[
 P_E:=\frac12\left[\one+A(\one+4D)^{-1/2}\right].
\]
Then
\[
 P_E^2=P_E,\qquad P_EE=EP_E,
\]
and
\[
 \norm{P_E-E}
 \le
 \alpha_B(E,\delta)
 :=
 \frac{\norm{2E-\one}}2
 \left((1-4\delta)^{-1/2}-1\right).
\]
If \(E=E^*\) on a finite-dimensional Hilbert space, then
\[
 P_E=\mathbf 1_{[1/2,\infty)}(E),\qquad
 \norm{P_E-E}
 \le
 \alpha_H(\delta):=
 \frac{1-\sqrt{1-4\delta}}2,
\]
and the constant \(\alpha_H\) is sharp.
\end{lemma}

\begin{proof}
The defect \(D\) is a polynomial in \(E\), hence commutes with \(A\), and
\[
 A^2=(2E-\one)^2=\one+4D.
\]
Because \(4\norm D<1\), the inverse square root
\((\one+4D)^{-1/2}\) is defined by its absolutely norm-convergent binomial
series.  Therefore
\[
 S:=A(\one+4D)^{-1/2}
\]
satisfies \(S^2=\one\), and \(P_E=(\one+S)/2\) is an idempotent commuting
with \(E\).  Furthermore,
\[
 P_E-E
 =\frac12A\left[(\one+4D)^{-1/2}-\one\right].
\]
Majorizing the absolute binomial coefficients gives
\[
 \norm{(\one+4D)^{-1/2}-\one}
 \le (1-4\norm D)^{-1/2}-1,
\]
which proves the Banach-space estimate.

In the self-adjoint case, the scalar inequality
\(|\lambda^2-\lambda|\le\delta<1/4\) separates the spectrum into two
components lying near \(0\) and \(1\).  Functional calculus therefore gives
\(P_E=\mathbf 1_{[1/2,\infty)}(E)\).  Solving
\(t(1-t)=\delta\) on \(0\le t\le1/2\) gives the maximal spectral distance
\[
 t=\frac{1-\sqrt{1-4\delta}}2
\]
to \(\{0,1\}\).  Equality is attained by a scalar spectral value at an
endpoint of the allowed component, proving sharpness.
\end{proof}

\begin{theorem}[Quantitative stability of minimal predictive memory]
\label{thm:predictive-compression-stability}
Let \(q:X\twoheadrightarrow Q\) be a finite linear predictive presentation.
Choose a bounded right inverse \(s:Q\to X\), \(qs=\id_Q\), and put
\(\kappa:=\norm{s}\).  Suppose \(E\in\End(X)\) satisfies
\[
 \norm{E^2-E}\le\delta<\frac14,\qquad
 \norm{qE-q}\le\varepsilon.
\]
Let \(P=P_E\) be the projection of
\Cref{lem:predictive-idempotentization}, and set
\[
 \alpha:=\norm{P-E},\qquad
 \eta:=\varepsilon+\norm q\,\alpha.
\]
Assume
\[
 \rank P=r=\dim Q,\qquad \kappa\eta<1.
\]
Then \(B:=qPs\in\End(Q)\) is invertible and
\[
 \boxed{R:=PsB^{-1}q}
\]
is an exact minimal predictive projection with
\[
 R^2=R,\qquad qR=q,\qquad
 \rank R=r,\qquad \operatorname{im}R=\operatorname{im}P.
\]
Moreover
\[
 \boxed{
 \norm{R-E}
 \le
 \alpha+
 \frac{\norm P\,\kappa\,\eta}{1-\kappa\eta}.}
\]
Thus approximate idempotence and approximate future sufficiency determine a
nearby exact minimal predictive memory whenever the near-one spectral
subspace has the correct dimension and remains transverse to the
observationally null directions.
\end{theorem}

\begin{proof}
First,
\[
 \norm{qP-q}
 \le \norm{q(P-E)}+\norm{qE-q}
 \le\eta.
\]
Hence
\[
 \norm{B-\id_Q}
 =\norm{(qP-q)s}
 \le\kappa\eta<1.
\]
The Neumann series makes \(B\) invertible and gives
\(\norm{B^{-1}}\le(1-\kappa\eta)^{-1}\).  The definition of \(R\) yields
\[
 qR=qPsB^{-1}q=BB^{-1}q=q
\]
and
\[
 R^2
 =PsB^{-1}qPsB^{-1}q
 =PsB^{-1}BB^{-1}q
 =R.
\]
The range of \(R\) lies in \(\operatorname{im}P\).  Since \(qR=q\),
\(\rank R\ge r\), whereas \(\rank R\le\rank P=r\); consequently
\(\rank R=r\) and the two ranges agree.

The map \(q|_{\operatorname{im}P}\) is therefore an isomorphism with inverse
\(PsB^{-1}\).  For every \(x\in X\),
\[
 Px=PsB^{-1}qPx,
\]
and hence
\[
 R-P=PsB^{-1}q(\one-P).
\]
Because \(q(\one-P)=q-qP\),
\[
 \norm{R-P}
 \le
 \norm P\,\kappa\,\norm{B^{-1}}\,\norm{q(\one-P)}
 \le
 \frac{\norm P\,\kappa\,\eta}{1-\kappa\eta}.
\]
Adding \(\norm{P-E}=\alpha\) proves the estimate.
\end{proof}

The theorem gives the central robustness conclusion: approximate predictive
compression is close to an exact quotient carrier provided the retained spectral
subspace has the correct dimension and is transverse to observationally null
directions.  The next corollary propagates this static estimate to the renewal
updates themselves.

\begin{corollary}[Stability of compressed renewal dynamics]
\label{cor:compressed-dynamics-stability}
Let \(U_\sigma:X\to X\) be bounded raw updates and let
\(M:=\operatorname{im}P=\operatorname{im}R\).  Define
\[
 \widetilde\Phi_\sigma:=PU_\sigma P|_M,\qquad
 \Phi_\sigma:=RU_\sigma R|_M.
\]
Then
\[
 \boxed{
 \norm{\widetilde\Phi_\sigma-\Phi_\sigma}
 \le\norm{P-R}\,\norm{U_\sigma}.}
\]
For every fixed word length, the corresponding products differ by the
standard telescoping sum and are therefore norm stable.
\end{corollary}

\begin{proof}
Both \(P\) and \(R\) restrict to the identity on their common range \(M\).
Thus, for \(x\in M\),
\[
 (\widetilde\Phi_\sigma-\Phi_\sigma)x
 =(P-R)U_\sigma x,
\]
which gives the displayed bound.  For products, insert and subtract one
factor at a time:
\[
 A_n\cdots A_1-B_n\cdots B_1
 =
 \sum_{j=1}^n
 A_n\cdots A_{j+1}(A_j-B_j)B_{j-1}\cdots B_1,
\]
and take norms.
\end{proof}

\begin{remark}[Why the nondegeneracy hypotheses are necessary]
\label{rem:predictive-stability-boundaries}
Approximate idempotence and exact sufficiency do not imply minimality:
if \(X=Q\oplus K_{\pred}\), \(q\) is the first projection, and \(E=\one_X\),
then \(E^2=E\) and \(qE=q\), but \(E\) retains all redundant memory.
Conversely, an exact rank-\(r\) projection may erase the predictive quotient.
The rank and transversality clauses in
\Cref{thm:predictive-compression-stability} therefore control genuinely
independent failure modes.
\end{remark}

Thus the predictive quotient is stable as a finite memory and its compressed
dynamics is stable along every fixed finite history.  Up to this point, however,
the carrier is only a finite real linear space.  We now add the operational
principles that determine its order structure, its scalar field, and ultimately
its associative operator algebra.

\subsection{Finite sharp-complex reconstruction}

To pass from predictive sets to operator algebras, additional operational
structure is required.  We isolate the hypotheses that produce the Jordan
stage, the complex field, and the associative orientation.

The order of the reconstruction is important.  Sharp purification first yields
spectral decomposition and a Euclidean Jordan algebra.  Two-path interference
then eliminates the non-complex Jordan alternatives, while a compatible phase
response chooses one of the two opposite associative products.  Each step solves
a distinct ambiguity and is therefore stated separately.

\begin{assumption}[Finite operational reconstruction package]
\label{ass:reconstruction-package}
On each operationally irreducible reduced predictive sector the following hold.
\begin{enumerate}[label=\textup{(O\arabic*)},leftmargin=3em]
\item normalized predictive states form a finite-dimensional compact convex set,
      admissible effects separate states, and the resolved process theory is
      closed under sequential and parallel composition and convex mixing;
\item \emph{reduced sharp-reversible realisation}: there is a unique deterministic
      discard; every physical process admits a pure reversible dilation with a
      pure environment, unique up to reversible changes of a minimal environment;
      operationally pure processes with no retained environmental record are
      closed under sequential and parallel composition; and every normalized
      pure preparation has a normalized pure reverse effect certain on it;
\item ordinary process reduction is parallel complete in the sense of
      \Cref{def:parallel-completeness};
\item for every nonclassical pair of primitive outcomes, the balanced two-path
      pure states form one effective transitive orbit of a compact connected
      one-dimensional phase group;
\item pair phases have primitive period $2\pi$ and obey the linear, pair-local,
      permutation-covariant and triangle-compatible spectral response stated in
      \Cref{ass:pair-phase-response} below.
\end{enumerate}
\end{assumption}

\begin{definition}[Finite sharp theory with purification]
\label{def:sharp-purification}
A finite reduced operational theory satisfies the \emph{sharp-purification
package} when:
\begin{enumerate}[label=\textup{(SP\arabic*)},leftmargin=3em]
\item it is causal, meaning that every system has a unique deterministic effect;
\item sequential and parallel composites of pure processes are pure;
\item every system possesses a pure effect occurring with unit probability on
      some normalized pure state;
\item every normalized state is the marginal of a pure bipartite state, and two
      purifications on the same purifying system differ by a reversible
      transformation of that system.
\end{enumerate}
\end{definition}

\begin{proposition}[Sharp purification from reduced reversible realisation]
\label{prop:reversible-realisation-sharp-purification}
Item~\textup{(O2)} of \Cref{ass:reconstruction-package} implies the
sharp-purification package of \Cref{def:sharp-purification}.
\end{proposition}

\begin{proof}
The unique discard in \textup{(O2)} is precisely causality, giving
\textup{(SP1)}.  Apply the reversible-dilation clause to a preparation
$\rho:I\to A$.  Its dilation is a pure preparation
$\Psi:I\to A\otimes E$ followed by discarding $E$, so every state has a
purification.  Essential uniqueness of minimal reversible dilations implies the
usual uniqueness of purifications up to a reversible transformation on the
purifying system; this is the reversible-realisation form of purification proved
in \cite{ChiribellaDAPS2009}.  Hence \textup{(SP4)} holds.

The no-retained-record closure of pure dilations under sequential and parallel
composition is exactly purity preservation, giving \textup{(SP2)}.  Finally, if
$\alpha:I\to A$ is normalized and pure, its normalized pure reverse effect
$\alpha^\sharp:A\to I$ satisfies
$\alpha^\sharp\circ\alpha=1$.  Thus a pure effect is certain on a pure state,
which is \textup{(SP3)}.
\end{proof}

\begin{proposition}[Reversible dilation alone does not imply pure sharpness]
\label{prop:dilation-does-not-imply-sharpness}
Unique discard, pure-environment reversible dilation, and essential uniqueness
of minimal dilations do not imply pure sharpness.  In particular, the
pure-process reversal clause in item~\textup{(O2)} of
\Cref{ass:reconstruction-package} is logically independent of the
transformation-level dilation clauses.
\end{proposition}

\begin{proof}
Begin with finite-dimensional complex quantum theory, retaining its state
spaces, channels, reversible dilations, tensor product, and deterministic
discard.  These data satisfy Stinespring dilation and essential uniqueness of
minimal dilations.  Modify only the admitted non-deterministic effects.

Fix \(0<\epsilon<1/2\).  On a qubit, for every rank-one projection \(P\), let
\[
 e_{P,\epsilon}(\rho)
 :=
 \epsilon+(1-2\epsilon)\Tr(P\rho),
\]
and admit the convex effect system generated by
\[
 0,\quad u,\quad e_{P,\epsilon},\quad u-e_{P,\epsilon}.
\]
The family separates states: equality of all \(e_{P,\epsilon}\)-values implies
equality of all rank-one projection probabilities because \(1-2\epsilon\ne0\),
and rank-one projections separate density matrices.  Thus the predictive
state theory remains reduced.

Every non-deterministic generator takes values in
\([\epsilon,1-\epsilon]\), so no such effect occurs with unit probability on
any state.  Convex combinations cannot increase this maximum unless the
deterministic effect \(u\) is used.  But \(u\) has a nontrivial refinement
\[
 u=e_{P,\epsilon}+(u-e_{P,\epsilon})
\]
into two nonzero admitted effects and is therefore not pure in the
operational, non-refinability sense.  Hence the modified theory has no pure
effect certain on a normalized pure state: pure sharpness fails.

The state spaces, channels, reversible dilations, minimal environments, and
discard have not changed.  Therefore the dilation assumptions alone cannot
determine pure sharpness; the admitted effect theory is independent
operational data.
\end{proof}

\begin{remark}[Exact boundary of the reconstruction package]
The reversible-dilation clauses derive causality and purification, while the
no-retained-record clause supplies purity preservation.  Pure sharpness
requires a separate preparation--observation reversal principle.  The last
clause of item~\textup{(O2)} is therefore not a hidden consequence of
Stinespring-type realization.
\end{remark}

\begin{theorem}[Operational diagonalisation]
\label{thm:operational-diagonalization}
Under \Cref{def:sharp-purification}, every normalized predictive state has a
finite spectral decomposition
\[
 \rho=\sum_{j=1}^{r}p_j\alpha_j,
 \qquad p_j>0,
 \qquad \sum_jp_j=1,
\]
where the $\alpha_j$ are perfectly distinguishable normalized pure states.
Reversible transformations act transitively on normalized pure states, and each
system has a unique invariant completely mixed state.
\end{theorem}

\begin{proof}
The diagonalisation statement is the finite operational diagonalisation theorem
of Chiribella and Scandolo \cite{ChiribellaScandolo2015}.  Its hypotheses are
exactly causality, purity preservation, pure sharpness, and purification, which
are \textup{(SP1)}--\textup{(SP4)}.  Pure-state transitivity and existence of the
invariant completely mixed state are standard consequences of purification with
essential uniqueness in the finite operational-probabilistic framework
\cite{ChiribellaDAPS2009}.
\end{proof}

\begin{theorem}[Jordan reconstruction from sharp purification]
\label{thm:operational-jordan}
Under \Cref{def:sharp-purification}, the ordered state space of every
operationally irreducible predictive sector is order-isomorphic to the positive
cone of a finite-dimensional Euclidean Jordan algebra.  The cone is self-dual
for a canonical invariant inner product, and the decompositions of
\Cref{thm:operational-diagonalization} are its Jordan spectral decompositions.
\end{theorem}

\begin{proof}
Sharp theories with purification possess a canonical state--effect dagger,
self-duality, pure face projectors, and the symmetry-of-transition-probability
property.  Barnum, Lee, Scandolo, and Selby prove that these structures identify
each finite system with the state space of a Euclidean Jordan algebra
\cite{BarnumLeeScandoloSelby2017}.  The present theory is already reduced, so its
states and effects satisfy the separation hypotheses of that theorem.  The
spectral decompositions agree with those of
\Cref{thm:operational-diagonalization} by uniqueness of the operational dagger
and the associated face projectors.
\end{proof}

\begin{remark}[What is now derived]
\label{rem:jordan-derived-not-assumed}
Homogeneity, self-duality, and spectral decomposition are no longer primitive
clauses of \Cref{ass:reconstruction-package}.  They are consequences of the
reduced sharp-reversible realisation principle through
\Cref{prop:reversible-realisation-sharp-purification,thm:operational-diagonalization,thm:operational-jordan}.
The remaining complex-field and associative-orientation inputs are genuinely
composite and interferometric: parallel completeness and the two-path phase
laws.
\end{remark}

The preceding results explain what is already fixed before complex quantum
structure is invoked: the predictive state cone is homogeneous, self-dual, and
spectral.  They do not yet distinguish real, complex, quaternionic, spin-factor,
or exceptional Jordan systems.  The next results use the geometry of balanced
two-path interference to make precisely that distinction.

Let $J$ be the Euclidean Jordan algebra supplied by
\Cref{thm:operational-jordan}, and let $(p_1,\ldots,p_r)$ be a Jordan frame.  We
now show that the two-path phase condition selects complex matrix type rather
than merely postulating a qubit.  For $i\ne j$, write $J_{ij}$ for the
Peirce coherence space and
\[
 J[i,j]:=\R p_i\oplus\R p_j\oplus J_{ij}
\]
for the rank-two face generated by $p_i,p_j$.  With the invariant Euclidean
norm, the balanced pure states in this face are canonically a sphere in
$J_{ij}$; see the rank-two classification in
\cite{JordanVonNeumannWigner1934,BarnumHilgert2020}.

\begin{theorem}[Two-path complex face]
\label{thm:two-path-complex-face}
Under item \textup{(O4)}, every nonclassical Peirce space has
\[
 \dim_{\R}J_{ij}=2,
\]
and every nonclassical rank-two face is Jordan-isomorphic to
$M_2(\C)_{\mathrm{sa}}$.  The connected phase group is isomorphic to $\T$ and
acts on $J_{ij}\cong\R^2$ by ordinary planar rotations, up to orientation.
\end{theorem}

\begin{proof}
Let $m=\dim_{\R}J_{ij}$.  After rescaling the invariant norm, the balanced pure
states form the sphere $S^{m-1}$.  A transitive orbit of a one-dimensional Lie
group has manifold dimension at most one, so $m-1\le1$.  The case $m=1$ gives
$S^0$, on which a connected group cannot act transitively; nonclassicality
excludes $m=0$.  Hence $m=2$.  The unique rank-two Euclidean Jordan algebra with
a two-dimensional Peirce coherence space is $M_2(\C)_{\mathrm{sa}}$.  A compact
connected one-dimensional Lie group is a circle, and every effective transitive
orthogonal circle action on $S^1$ is conjugate to the standard rotation action,
with the two conjugacy choices corresponding to orientation.
\end{proof}

\begin{corollary}[Pairwise complex matrix type]
\label{cor:pairwise-complex-type}
Every nonclassical irreducible predictive Jordan sector is
\[
 J\cong M_r(\C)_{\mathrm{sa}}
\]
for some $r\ge2$; the remaining irreducible sectors are one-dimensional
classical sectors.
\end{corollary}

\begin{proof}
For rank two this is \Cref{thm:two-path-complex-face}.  For rank at least three,
the simple Euclidean Jordan classification gives real, complex, quaternionic
Hermitian matrix algebras and the exceptional Albert algebra; their pairwise
Peirce dimensions are respectively $1,2,4,8$.  Spin factors have rank two.
Since every nonclassical pair has Peirce dimension two, only the complex
Hermitian matrix algebra remains
\cite{JordanVonNeumannWigner1934}.  Operational irreducibility excludes a
nontrivial direct sum inside one sector.
\end{proof}

Pairwise two-path interference has now selected complex Hermitian matrix type, but
the Jordan product alone still forgets the sign of the commutator and hence the
orientation of the associative multiplication.  The following phase-response
assumption and rigidity lemma recover this remaining binary choice consistently
across all pairs.

\begin{assumption}[Compatible spectral phase response]
\label{ass:pair-phase-response}
For $a=\sum_k a_kp_k$ diagonal in a Jordan frame, let
$\ell_{ij}(a)$ be the oriented angular velocity induced on $J_{ij}$.  Assume:
\begin{enumerate}[label=\textup{(P\arabic*)},leftmargin=2.8em]
\item $a\mapsto\ell_{ij}(a)$ is real linear and is unchanged by adding a scalar
      multiple of $\one$;
\item $\ell_{ij}(a)$ depends only on $a_i,a_j$ and is covariant under frame
      permutations;
\item the contrast $p_i-p_j$ has primitive unit angular speed;
\item for every triple of distinct indices,
      $\ell_{ij}(a)+\ell_{jk}(a)+\ell_{ki}(a)=0$.
\end{enumerate}
\end{assumption}

\begin{lemma}[Rigidity of pairwise phase velocity]
\label{lem:phase-velocity}
Under \Cref{ass:pair-phase-response}, there is one sign $s\in\{\pm1\}$ on each
connected nonclassical block such that
\[
 \ell_{ij}(a)=\frac{s}{2}(a_i-a_j)
\]
for every frame and every pair.
\end{lemma}

\begin{proof}
Linearity and pair locality give
$\ell_{ij}(a)=\alpha_{ij}a_i+\beta_{ij}a_j$.  Scalar invariance gives
$\alpha_{ij}+\beta_{ij}=0$, so
$\ell_{ij}(a)=c_{ij}(a_i-a_j)$.  Since the contrast $p_i-p_j$ has eigenvalues
$1,-1$ on the pair, primitive unit speed gives $2|c_{ij}|=1$.  Write
$c_{ij}=s_{ij}/2$.  Substituting the three formulas into the triangle identity
and comparing the independent coefficients of $a_i,a_j,a_k$ gives
$s_{ij}=s_{jk}=s_{ik}$.  Connectivity of a Jordan frame propagates one sign
through the irreducible block, and frame covariance makes it independent of the
chosen frame.
\end{proof}

\begin{theorem}[Finite complex predictive algebra]
\label{thm:complex-algebra}
Under \Cref{ass:reconstruction-package}, the predictive observable algebra is
\[
 \boxed{
 \A_{\pred}\cong\bigoplus_{r=1}^{s}M_{n_r}(\C).}
\]
On each nonclassical connected block there is one sign $s=\pm1$ such that the
operational phase flow generated by $a=a^*$ is
\[
 \tau_a^t(b)=e^{sita/2}b\,e^{-sita/2},
 \qquad
 \frac{\dd}{\dd t}\tau_a^t(b)\bigg|_{t=0}
 =s\frac{i}{2}[a,b].
\]
Thus the ordered Jordan structure determines complex matrix type, while the
composed phase response selects one of the two opposite associative products.
The remaining freedom is blockwise reversal of this associative orientation.
\end{theorem}

\begin{proof}
By
\Cref{prop:reversible-realisation-sharp-purification,thm:operational-jordan},
items \textup{(O1)}--\textup{(O2)} supply the finite Euclidean Jordan
decomposition without separately assuming homogeneity, self-duality, or
spectrality.  Item~\textup{(O4)} and
\Cref{cor:pairwise-complex-type} then force every nonclassical irreducible
Jordan block to be complex Hermitian matrix type.

In a frame diagonalizing $a$, item~\textup{(O5)} and
\Cref{lem:phase-velocity} give the angular velocity
$s(a_i-a_j)/2$ on every Peirce coherence plane.  The unitary
$e^{sita/2}$ fixes each $p_i$ and has exactly those velocities, so the
operational and matrix phase flows agree on every Peirce summand and hence on
the whole algebra.  Degenerate spectral subspaces cause no ambiguity because
the exponential is scalar there.  Differentiation yields
$s\frac{i}{2}[a,\,\cdot\,]$.  Replacing $s$ by $-s$ reverses the
associative multiplication while leaving the Jordan product, order, and all
single-system probabilities unchanged; compare
\cite{AlfsenShultz1998,AlfsenShultz2003}.

Item~\textup{(O3)} is not used to select the single-system complex field.  It
ensures that independently prepared systems have no hidden ancillary fibres and
supplies the faithful composite matrix ordering used in the UCP step below.
Finite-dimensionality gives the stated direct sum of full matrix algebras.
\end{proof}

\begin{remark}[Sharp logical boundary]
Real or quaternionic quantum systems, higher spin factors, and the exceptional
Jordan system remain possible if the one-dimensional transitive two-path phase
condition is removed.  Conversely, single-system complex interference does not
imply parallel completeness: a composite may contain globally visible degrees
of freedom invisible to all product tests.  The complex selector and the
no-holism axiom are therefore independent.
\end{remark}

The complex-algebra theorem completes the single-system reconstruction.  Its
conclusion is stronger than a choice of notation: it shows that complex matrix
algebras and their commutator orientation arise from operational phase data.  A
separate question remains for transformations.  Positivity on one system does not
guarantee positivity when the system is embedded in a larger experiment, so the
next subsection identifies ancillary stability with complete positivity.

\subsection{Ancillary stability and complete positivity}

Let $V_A=\A_A^{\mathrm{sa}}$ be the observable order-unit space associated with
a predictive object $A$.  A deterministic process $T:A\to B$ acts in the
Heisenberg picture by a linear map
\[
 \Phi_T:\A_B\longrightarrow\A_A.
\]
Closed-experiment positivity makes $\Phi_T$ positive, and preservation of
discard makes it unital.  Complete positivity is the additional statement
that this remains true in every matrix ancillary context.

This is the transformation-level counterpart of parallel completeness.  The
reconstructed algebra tells us what the predictive observables are; complete
ancillary stability tells us which positive updates remain admissible in every
composite context.

\begin{definition}[Complete ancillary stability]
\label{def:ancillary-stability}
A deterministic predictive transformation $T$ is completely ancillary stable
when, for every $k\ge1$, its extension by the identity on a $k$-level complex
ancilla induces a positive map
\[
 \id_{M_k}\otimes\Phi_T:
 M_k(\C)\otimes\A_B\longrightarrow M_k(\C)\otimes\A_A.
\]
\end{definition}

\begin{theorem}[Operational-to-UCP bridge]
\label{thm:operational-ucp}
Assume \Cref{ass:reconstruction-package}.  Every deterministic predictive
transformation that is completely ancillary stable induces a unital completely
positive map on $\A_{\pred}$.  Conversely every UCP map defines a deterministic
transformation of the finite predictive operator system and remains positive
under arbitrary finite ancillas.  Thus the downstream renewal alphabet may be
taken to be a finite family
\[
 \Phi_\sigma:\A_{\pred}\to\A_{\pred},
 \qquad \sigma\in E,
\]
of UCP maps, derived from operational prediction rather than imposed as
primitive geometry.
\end{theorem}

\begin{proof}
Let $T$ be deterministic.  If $b\ge0$ is an effect or positive observable in
$\A_B$, then every state $\omega$ on $\A_A$ gives a nonnegative operational
probability $\omega(\Phi_T(b))$; state separation implies
$\Phi_T(b)\ge0$.  Determinism gives $\Phi_T(\one)=\one$, so $\Phi_T$ is positive
and unital.

By complete ancillary stability, $\id_{M_k}\otimes\Phi_T$ is positive for all
$k$.  This is the definition of complete positivity.  In finite dimensions it
is enough to test $k$ equal to the largest matrix-block size, by Choi's theorem
\cite{Choi1975}.  Hence $\Phi_T$ is UCP.

Conversely, a UCP map is unital and positive, so it preserves the order unit
and sends effects to effects.  Tensoring with an identity map preserves
complete positivity, hence positivity, on every finite ancilla.  Pairing with
states therefore defines consistent probabilities in all parallel contexts,
which is precisely a deterministic transformation of the reconstructed
operator system.
\end{proof}

\begin{corollary}[Predictive units are automorphisms]
\label{cor:units-automorphisms}
If a UCP predictive transformation $\Phi$ has a UCP inverse, then $\Phi$ is a
$*$-automorphism.
\end{corollary}

\begin{proof}
Let $\Psi$ be the UCP inverse.  Kadison--Schwarz gives
$\Phi(a)^*\Phi(a)\le\Phi(a^*a)$.  Applying $\Psi$ and then Kadison--Schwarz to
$\Psi$ at $\Phi(a)$ yields
\[
 a^*a=\Psi(\Phi(a))^*\Psi(\Phi(a))
 \le\Psi(\Phi(a)^*\Phi(a))\le a^*a.
\]
Thus equality holds and every $a$ lies in the multiplicative domain of
$\Phi$.  Hence $\Phi$ is a unital $*$-homomorphism; bijectivity makes it a
$*$-automorphism.
\end{proof}

\begin{remark}[The first pillar]
The output of this section is not yet spacetime, causality, or a signed
geometry.  It is the canonical finite object on which those structures may be
built: a reduced predictive $C^*$-algebra and a monoid of UCP renewal updates.
The construction has separated three issues that are often conflated:
predictive quotienting is categorical, complex algebra reconstruction is
operational, and complete positivity is ancillary stability.
\end{remark}

We have therefore reached the first finite dynamical object promised in the
introduction: a reduced complex predictive algebra acted on by UCP renewal maps.
The remaining issue in this section is one of identification and reference.  We
must know when operationally indistinguishable histories determine the same full
channel, and we must select a faithful stationary state relative to which forward
dynamics can later be compared with retrodiction.

\subsection{Tomographic identification and a canonical stationary reference}

The predictive quotient is defined operationally.  To identify a class of
histories with a full accumulated channel, rather than only with its action on
one prepared orbit, one needs a process-tomography interface.

The distinction is subtle but essential for the later branchwise arguments.  A
rooted history only probes the orbit generated from the preparations actually
available, whereas a channel is a transformation on the entire state space.  The
following assumption makes the passage from the first notion to the second
explicit rather than tacit.

\begin{assumption}[Preparation--effect process tomography]
\label{ass:process-tomography}
For every reconstructed predictive object, the admissible preparations span
its real ordered state space and the admissible effects separate that space.
Two transformations are operationally equal when all pairings
\[
 a(Tv)=a(T'v)
\]
agree for admissible spanning preparations $v$ and separating effects $a$.
\end{assumption}

\begin{proposition}[Operational profiles determine accumulated channels]
\label{prop:tomographic-channel-equality}
Under \Cref{ass:process-tomography}, two generated histories have identical
transformation-level operational profiles if and only if their accumulated
linear state maps agree.  Equivalently, after the reconstruction of
\Cref{thm:complex-algebra,thm:operational-ucp}, their accumulated Heisenberg
UCP maps agree.
\end{proposition}

\begin{proof}
Equality of the accumulated maps plainly implies equality of every
preparation--effect pairing.  Conversely, suppose all such pairings agree.
For each admissible spanning preparation $v$, separation by admissible effects
gives $Tv=T'v$.  Linearity and spanning then imply $T=T'$ on the entire state
space.  Taking duals gives equality of the Heisenberg maps.
\end{proof}

\begin{example}[A rooted orbit does not determine a channel]
\label{ex:rooted-orbit-not-tomographic}
Let \(V=\R^2\), suppose the only admissible initial preparation is
\(v_0=e_1\), and let the terminal effects separate vectors in \(V\).  Define
\[
 T=\begin{pmatrix}1&0\\0&1\end{pmatrix},
 \qquad
 T'=\begin{pmatrix}1&1\\0&0\end{pmatrix}.
\]
Then \(Tv_0=T'v_0=e_1\).  If every admissible future continuation is generated
from the actual rooted output and never supplies an independent preparation
with a component along \(e_2\), all rooted future statistics of \(T\) and
\(T'\) agree.  Nevertheless
\[
 Te_2=e_2,\qquad T'e_2=e_1,
\]
so \(T\ne T'\).  Thus equality on one reachable rooted orbit does not imply
equality of accumulated transformations.  The spanning-preparation clause in
\Cref{ass:process-tomography} is therefore an independent interface
condition, not a consequence of arbitrarily many future tests on the given
orbit.
\end{example}

The proposition and counterexample isolate the exact role of tomography: future
profiles determine accumulated channels only when preparations span and effects
separate.  Once this interface is present, the UCP alphabet may be treated as a
genuine family of channels.  We can then ask which state provides the common
reference needed for stationarity and retrodictive duality.

Let a selected simple recurrent block be $M_n(\C)$ and let
$\{\Phi_\sigma\}_{\sigma\in E}$ be its UCP renewal alphabet.  A stationary
reference cannot be selected from the unweighted alphabet alone: one must also
specify how elementary renewals are sampled.

\begin{definition}[Weighted averaged renewal transfer]
\label{def:averaged-transfer}
Choose $p_\sigma>0$ with $\sum_\sigma p_\sigma=1$ and set
\[
 \overline\Phi:=\sum_{\sigma\in E}p_\sigma\Phi_\sigma,
 \qquad T:=\overline\Phi_*.
\]
The trace-preserving CP map $T$ is \emph{primitive} when some $m\ge1$ satisfies
\[
 X\ge0,\ X\ne0 \quad\Longrightarrow\quad T^m(X)>0.
\]
\end{definition}

\begin{theorem}[Primitive transfer selects a unique faithful state]
\label{thm:primitive-stationary-state}
If the weighted averaged renewal transfer of \Cref{def:averaged-transfer} is
primitive, then there is a unique density matrix $\rho_*>0$ satisfying
\[
 T(\rho_*)=\rho_*.
\]
Moreover, for some $m\ge1$ and $0<c<1$,
\[
 \norm{T^{km}(\rho)-\rho_*}_1
 \le c^k\norm{\rho-\rho_*}_1
\]
for every density matrix $\rho$ and every $k\ge0$.
\end{theorem}

\begin{proof}
The compact convex set of density matrices is invariant under the continuous
affine map $T$, so it has a fixed point $\rho_*$ by Brouwer's theorem.
Primitivity gives $\rho_*=T^m(\rho_*)>0$.

Let
\[
 \K_0:=\{X=X^*: \Tr X=0,\ \norm X_1=1\}.
\]
This set is compact.  For $X\in\K_0$, write its Jordan decomposition
$X=X_+-X_-$; then $\Tr X_+=\Tr X_-=1/2$.  Primitivity makes
$A=T^m(X_+)$ and $B=T^m(X_-)$ positive definite.  Hence for some
$\delta_X>0$ one has $A,B\ge\delta_X\one$, and therefore
\[
 \norm{T^m(X)}_1=\norm{A-B}_1
 \le \Tr(A-\delta_X\one)+\Tr(B-\delta_X\one)
 =1-2n\delta_X<1.
\]
The continuous function $X\mapsto\norm{T^m(X)}_1$ has a maximum $c<1$ on
$\K_0$.  Homogeneity gives
$\norm{T^m(X)}_1\le c\norm X_1$ for every traceless Hermitian $X$.
Applying this to $\rho-\rho_*$ and iterating proves the estimate.  If two
stationary densities existed, the same strict contraction applied to their
difference would force them to coincide.
\end{proof}

\begin{corollary}[Common KMS--Petz family]
\label{cor:common-petz-family}
Under \Cref{thm:primitive-stationary-state}, the averaged channel
$\overline\Phi$ has a canonical UCP KMS--Petz dual relative to $\rho_*$.  If,
in addition,
\[
 (\Phi_\sigma)_*(\rho_*)=\rho_*
 \qquad(\sigma\in E),
\]
then every elementary channel has a UCP KMS--Petz dual relative to the same
faithful state.  Without channelwise stationarity the elementary Petz
transposes remain CP, but need not be unital.
\end{corollary}

\begin{proof}
Apply \Cref{thm:petz-properties} to $\overline\Phi$ and to each channel that
preserves $\rho_*$.  From the defining formula,
$\Phi_{\rho_*}^\sharp(\one)=\one$ is equivalent to
$\Phi_*(\rho_*)=\rho_*$.
\end{proof}

The weighted primitive transfer completes the predictive stage by selecting a
unique faithful reference state and, under channelwise stationarity, a common
family of KMS--Petz duals.  Section~
ef{sec:predictive} has therefore supplied
three ingredients: the minimal predictive memory, its UCP renewal dynamics, and
the stationary reference against which opposite temporal comparisons can be
made.  The next section shows that two such comparisons---real affinity and record
orientation---must nevertheless be kept distinct.

% =====================================================================
\section{Retrodictive duality and the signed orientation line}
\label{sec:retrodictive}

The positive predictive algebra distinguishes what future tests can retain.  A
Lorentzian construction additionally needs an oppositely oriented comparison.
There are two different comparisons available: a real forward/backward
weight, measured by affinity, and a parity of record transport, measured by a
determinant line.  This section proves that they are mathematically
independent and that only the second classifies the signed cover.

The section follows the operational order in which these structures become
available.  We first construct the canonical KMS--Petz dual.  We then resolve
individual outcomes and prove that conditional predictive successors enforce
multiplicative branch weights.  Those weights define real affinity, while stable
future-distinct records define an independent determinant parity.  The final
results compare the two classes and classify the additional amplitude lifts of
the signed cover.

\subsection{The canonical KMS--Petz retrodictive dual}

Let $\A=M_n(\C)$ for the moment; direct sums are handled blockwise.  Let
$\rho>0$ be a faithful density matrix and let $\Phi:\A\to\A$ be UCP.  Write
$\Phi_*$ for the trace adjoint, determined by
\[
 \Tr(\Phi_*(x)a)=\Tr(x\Phi(a)).
\]
Assume $\rho$ is stationary, $\Phi_*(\rho)=\rho$.

\begin{definition}[Symmetric KMS pairing]
\label{def:kms-pairing}
Define
\[
 \langle a,b\rangle_{\rho,1/2}
 :=\Tr\!\left(\rho^{1/2}a^*\rho^{1/2}b\right),
 \qquad
 \GammaR(a):=\rho^{1/2}a\rho^{1/2}.
\]
\end{definition}

\begin{definition}[KMS--Petz anti-renewal dual]
\label{def:petz-dual}
The retrodictive or anti-renewal dual of $\Phi$ relative to $\rho$ is
\[
 \boxed{
 \Phi_\rho^\sharp
 :=\GammaR^{-1}\circ\Phi_*\circ\GammaR,
 \qquad
 \Phi_\rho^\sharp(a)
 =\rho^{-1/2}\Phi_*(\rho^{1/2}a\rho^{1/2})\rho^{-1/2}.}
\]
\end{definition}

\begin{theorem}[Canonical retrodictive adjoint]
\label{thm:petz-properties}
Let $\Phi$ be UCP and let $\rho>0$ satisfy $\Phi_*(\rho)=\rho$.  Then:
\begin{enumerate}[label=\textup{(\roman*)},leftmargin=2.4em]
\item $\Phi_\rho^\sharp$ is UCP and preserves the same stationary state;
\item it is the adjoint of $\Phi$ for the symmetric KMS pairing,
      \[
      \langle a,\Phi(b)\rangle_{\rho,1/2}
      =\langle\Phi_\rho^\sharp(a),b\rangle_{\rho,1/2};
      \]
\item $(\Phi_\rho^\sharp)_\rho^\sharp=\Phi$;
\item if $\Phi$ and $\Psi$ preserve $\rho$, then
      \[
      (\Phi\circ\Psi)_\rho^\sharp
      =\Psi_\rho^\sharp\circ\Phi_\rho^\sharp.
      \]
\end{enumerate}
Thus retrodiction reverses compositional order but is not an inverse channel in
general.
\end{theorem}

This theorem gives the unique backward comparison compatible with the stationary
KMS pairing.  Its importance is conceptual: retrodiction is an adjoint operation,
not an undoing of dissipation.  The next proposition compares this choice with
more familiar adjoints and explains why the symmetric modular placement is needed
to retain complete positivity.

\begin{proof}
The maps $a\mapsto\rho^{1/2}a\rho^{1/2}$ and
$x\mapsto\rho^{-1/2}x\rho^{-1/2}$ are completely positive congruences.
The trace adjoint $\Phi_*$ of a CP map is CP.  Hence the composition defining
$\Phi_\rho^\sharp$ is CP.  Stationarity gives
\[
 \Phi_\rho^\sharp(\one)
 =\rho^{-1/2}\Phi_*(\rho)\rho^{-1/2}=\one,
\]
so the dual is UCP.  To verify that it preserves the same state, for every
$a\in\A$ compute
\[
 \Tr\!\left(\rho\,\Phi_\rho^\sharp(a)\right)
 =\Tr\!\left(\Phi_*(\rho^{1/2}a\rho^{1/2})\right)
 =\Tr\!\left(\rho^{1/2}a\rho^{1/2}\Phi(\one)\right)
 =\Tr(\rho a).
\]
Thus $(\Phi_\rho^\sharp)_*(\rho)=\rho$.

For the KMS adjoint identity,
\begin{align*}
 \langle a,\Phi(b)\rangle_{\rho,1/2}
 &=\Tr\!\left(\rho^{1/2}a^*\rho^{1/2}\Phi(b)\right)\\
 &=\Tr\!\left(\Phi_*(\rho^{1/2}a^*\rho^{1/2})b\right)\\
 &=\Tr\!\left(\rho^{1/2}(\Phi_\rho^\sharp(a))^*
                \rho^{1/2}b\right),
\end{align*}
using $*$-preservation of positive maps.  This proves (ii).  The adjoint of an
adjoint for a nondegenerate inner product is the original map, giving (iii).
Alternatively it follows by direct substitution and the relation between the
trace adjoints.  Finally,
\[
 (\Phi\Psi)_*=\Psi_*\Phi_*,
\]
so conjugation by $\GammaR$ gives (iv).
\end{proof}

\begin{proposition}[Why the symmetric KMS adjoint is selected]
\label{prop:reverse-comparison}
Besides $\Phi_\rho^\sharp$, two common alternatives are the Hilbert--Schmidt
adjoint $\Phi^\dagger:=\Phi_*$ and the right-GNS adjoint
\[
 \Phi_\rho^{\mathrm G}(a):=\Phi_*(a\rho)\rho^{-1}.
\]
The Hilbert--Schmidt adjoint is CP and trace preserving but is unital only when
$\Phi$ is bistochastic.  The right-GNS adjoint satisfies
\[
 \Phi_\rho^{\mathrm G}
 =\mathcal M_\rho^{-1}\circ\Phi_\rho^\sharp\circ\mathcal M_\rho,
 \qquad
 \mathcal M_\rho(a)=\rho^{-1/2}a\rho^{1/2},
\]
and need not preserve Hermiticity because $\mathcal M_\rho$ is generally not
$*$-preserving.  If $\Phi$ commutes with the modular group of $\rho$, then
$\Phi_\rho^{\mathrm G}=\Phi_\rho^\sharp$.  Thus the symmetric placement of the
two factors $\rho^{1/2}$ is precisely what gives an unconditional CP adjoint
relative to a faithful reference state.
\end{proposition}

\begin{proof}
The Hilbert--Schmidt assertions follow from trace duality, with unitality
$\Phi_*(\one)=\one$ equivalent to trace preservation of the Heisenberg map.
Let $R_\rho(a)=a\rho$.  By definition
$\Phi_\rho^{\mathrm G}=R_\rho^{-1}\Phi_*R_\rho$, while
$R_\rho=\Gamma_\rho\mathcal M_\rho$.  Substitution gives the displayed
similarity relation.  If $\Phi$ is modular covariant, then so are $\Phi_*$ and
$\Phi_\rho^\sharp$; finite-dimensional analytic continuation to the imaginary
half-step $\sigma_{i/2}^\rho=\mathcal M_\rho$ makes the similarity trivial.
\end{proof}

\begin{example}[The right-GNS adjoint can fail positivity]
\label{ex:gns-not-positive}
Let
\[
 H=\frac1{\sqrt2}\begin{pmatrix}1&1\\1&-1\end{pmatrix},
 \qquad
 \tau=\begin{pmatrix}4/5&0\\0&1/5\end{pmatrix},
\]
and define the Schr\"odinger channel
\[
 T(x)=\frac12HxH^*+\frac12\Tr(x)\tau.
\]
It has faithful stationary state
\[
 \rho=\begin{pmatrix}7/10&1/10\\1/10&3/10\end{pmatrix}.
\]
For the Heisenberg channel $\Phi=T^*$,
\[
 \Phi_\rho^{\mathrm G}(E_{00})
 =T(E_{00}\rho)\rho^{-1}
 =\begin{pmatrix}129/200&57/200\\19/100&67/100\end{pmatrix},
\]
which is not Hermitian.  Hence the right-GNS adjoint is not positive, whereas
$\Phi_\rho^\sharp$ is UCP by \Cref{thm:petz-properties}.
\end{example}

\begin{definition}[KMS detailed balance and cancellation current]
\label{def:detailed-balance}
The channel satisfies symmetric KMS detailed balance when
$\Phi=\Phi_\rho^\sharp$.  Its oriented cancellation current is
\[
 \mathcal J_\rho(\Phi):=\Phi-\Phi_\rho^\sharp.
\]
\end{definition}

\begin{proposition}[Exact cancellation decomposition]
\label{prop:cancellation}
Every stationary UCP channel has the unique decomposition
\[
 \Phi=\Phi_{\mathrm{sym}}+\Phi_{\mathrm{asym}},
 \qquad
 \Phi_{\mathrm{sym}}:=\frac12(\Phi+\Phi_\rho^\sharp),
 \qquad
 \Phi_{\mathrm{asym}}:=\frac12\mathcal J_\rho(\Phi),
\]
where $\Phi_{\mathrm{sym}}$ is KMS self-adjoint and
$\Phi_{\mathrm{asym}}$ is KMS skew-adjoint.  Detailed balance is equivalent to
$\mathcal J_\rho(\Phi)=0$.
\end{proposition}

\begin{proof}
Apply the involution $\sharp$ of \Cref{thm:petz-properties} to the two terms.
One obtains $\Phi_{\mathrm{sym}}^\sharp=\Phi_{\mathrm{sym}}$ and
$\Phi_{\mathrm{asym}}^\sharp=-\Phi_{\mathrm{asym}}$.  Their sum is $\Phi$ and
the decomposition into the $\pm1$ eigenspaces of an involution is unique.
The last assertion is immediate.
\end{proof}

\begin{theorem}[Composition law for the cancellation current]
\label{thm:current-composition}
Let $\Phi$ and $\Psi$ preserve the same faithful state $\rho$.  Then
\[
\boxed{
\begin{aligned}
 \mathcal J_\rho(\Phi\circ\Psi)
 &=\mathcal J_\rho(\Phi)\circ\Psi
   +\Phi_\rho^\sharp\circ\mathcal J_\rho(\Psi)
   +[\Phi_\rho^\sharp,\Psi_\rho^\sharp]\\
 &=\Phi\circ\mathcal J_\rho(\Psi)
   +\mathcal J_\rho(\Phi)\circ\Psi_\rho^\sharp
   +[\Phi_\rho^\sharp,\Psi_\rho^\sharp].
\end{aligned}}
\]
Consequently the current is not an ordinary derivation on a noncommutative
renewal monoid: reversal of compositional order produces the commutator term.
\end{theorem}

\begin{proof}
Since $(\Phi\Psi)_\rho^\sharp=\Psi_\rho^\sharp\Phi_\rho^\sharp$,
\[
 \mathcal J_\rho(\Phi\Psi)
 =\Phi\Psi-\Psi_\rho^\sharp\Phi_\rho^\sharp.
\]
For the first identity, insert and subtract
$\Phi_\rho^\sharp\Psi$ and
$\Phi_\rho^\sharp\Psi_\rho^\sharp$; for the second, insert and subtract
$\Phi\Psi_\rho^\sharp$ and
$\Phi_\rho^\sharp\Psi_\rho^\sharp$.
\end{proof}

\begin{corollary}[Balanced factors can have an unbalanced product]
\label{cor:balanced-product}
If $\Phi=\Phi_\rho^\sharp$ and $\Psi=\Psi_\rho^\sharp$, then
\[
 \mathcal J_\rho(\Phi\circ\Psi)=[\Phi,\Psi].
\]
Thus the composite of two individually KMS-detailed-balanced channels is
KMS-detailed-balanced if and only if they commute.
\end{corollary}

\begin{proof}
Substitute the two vanishing individual currents into
\Cref{thm:current-composition}.
\end{proof}

\begin{remark}[Anti-renewal terminology]
In this paper \emph{anti-renewal} means the KMS--Petz retrodictive dual.  It is
not the opposite sheet of the later signed cover, not microscopic time
reversal, and not the inverse of a dissipative channel.  The opposite sheet is
called the \emph{counter-oriented renewal realisation}.
\end{remark}


At the channel level, the KMS--Petz dual is now fixed and its failure of detailed
balance is measured by the cancellation current.  These statements concern
averaged transformations.  To obtain path weights, cycle affinities, and later a
renewal graph, we must descend to resolved outcomes and identify the predictive
state reached after each successful branch.

\subsection{Conditional predictive unifilarity and derived reference renewal}
\label{subsec:conditional-reference-renewal}

The scalar edge law needed below is not fundamentally a consequence of
stationarity.  It is forced when resolved outcomes are represented at the same
operational resolution as their conditional predictive successors.

The key point is that scalar renewal is not assumed at the microscopic level.  It
will emerge because conditioning removes the probability of the past, and future
equivalence makes the observed outcome determine a unique successor class.  A
future-tomographic realization then has no freedom to represent that successor by
a different normalized state.

\begin{definition}[Conditional future profile]
\label{def:conditional-future-profile}
Let $h:I\to B$ be a successful resolved history and put
\[
 p(h):=\Pr(u_B\circ h)>0.
\]
A future experiment after $h$ is an admissible resolved continuation followed
by a terminal effect, written as a single process $F:B\to I$.  Define
\[
 \Pred^{\mathrm{cond}}(h)(F)
 :=\frac{\Pr(F\circ h)}{p(h)}.
\]
Two successful histories are \emph{conditionally predictively equivalent},
written $h\sim_{\mathrm{cond}}h'$, when their conditional future profiles agree
for every admissible future experiment.
\end{definition}

The normalization removes the probability that the past itself occurred.
Probabilities of later resolved outcomes remain part of the profile because
such outcomes may be included at the beginning of $F$.

\begin{lemma}[Resolved-outcome successor congruence]
\label{lem:conditional-successor}
Let $h\sim_{\mathrm{cond}}h'$ and let $e:B\to C$ be a resolved continuation
admissible after both histories.  If
\[
 p(e\circ h)>0,\qquad p(e\circ h')>0,
\]
then
\[
 e\circ h\sim_{\mathrm{cond}}e\circ h'.
\]
\end{lemma}

\begin{proof}
Let $F:C\to I$ be any future experiment after $e$.  Conditional equivalence of
$h$ and $h'$ applied first to $F\circ e$ and then to the discard experiment
$u_C\circ e$ gives
\[
 \frac{\Pr(F\circ e\circ h)}{p(h)}
 =
 \frac{\Pr(F\circ e\circ h')}{p(h')},
 \qquad
 \frac{p(e\circ h)}{p(h)}
 =
 \frac{p(e\circ h')}{p(h')}.
\]
The second common value is positive.  Dividing the first equality by the
second yields
\[
 \frac{\Pr(F\circ e\circ h)}{p(e\circ h)}
 =
 \frac{\Pr(F\circ e\circ h')}{p(e\circ h')},
\]
which is precisely conditional predictive equivalence of the extended
histories.
\end{proof}

\begin{definition}[Conditional predictive quotient and successor map]
\label{def:conditional-predictive-quotient}
Let
\[
 \Qq_{\mathrm{cond}}
 :=
 \{\text{successful resolved histories}\}/\!\sim_{\mathrm{cond}}.
\]
For a resolved outcome $e$ and a class $x=[h]$ for which
$p(e\circ h)>0$, define
\[
 \delta_e(x):=[e\circ h].
\]
By \Cref{lem:conditional-successor}, this is a well-defined partial map.  The
conditional predictive presentation is therefore \emph{outcome unifilar}: the
current predictive class and the observed resolved outcome determine one
successor class.
\end{definition}

\begin{remark}[Predictive, not microscopic, unifilarity]
No unifilar hidden-state model is assumed.  Many microscopic states may realize
one conditional predictive class, and their hidden dynamics may be
nonunifilar.  Deterministic successors appear only after quotienting pasts by
all conditional future predictions.
\end{remark}

Let $X$ be a finite reachable component of
$\Qq_{\mathrm{cond}}$.  Represent every $x\in X$ by a normalized density
matrix $\rho_x$ on its finite-dimensional support algebra $\A_x$, so
$\rho_x$ is faithful on $\A_x$.  A resolved
edge $e:x\to y=\delta_e(x)$ is represented in the Schr\"odinger picture by a
CP trace-nonincreasing map
\[
 \E_e:\A_x\longrightarrow\A_y.
\]

\begin{definition}[Reduced resolved operator realization]
\label{def:reduced-resolved-realization}
The family $(\rho_x,\E_e)$ is a \emph{reduced resolved realization} when:
\begin{enumerate}[label=\textup{(F\arabic*)},leftmargin=2.8em]
\item the future statistics of $\rho_x$ equal the conditional future profile of
      the class $x$;
\item for every positive-weight edge
      \[
       w_e:=\Tr\E_e(\rho_x)>0,
      \]
      the normalized branch output
      \[
       \widehat\rho_e:=\frac{\E_e(\rho_x)}{w_e}
      \]
      represents the operational history obtained by appending the resolved
      outcome $e$;
\item the represented admissible future effects separate the reachable
      normalized states in every target algebra.
\end{enumerate}
The second clause is implementation consistency; the third is the
future-tomographic reducedness condition.
\end{definition}

\begin{theorem}[Predictive derivation of reference renewal]
\label{thm:predictive-reference-renewal}
In every reduced resolved operator realization and for every positive-weight
edge $e:x\to y=\delta_e(x)$,
\[
 \boxed{
 \E_e(\rho_x)=w_e\rho_y,
 \qquad
 w_e=\Tr\E_e(\rho_x).}
\]
Thus branchwise renewal on the canonical predictive orbit is derived without
assuming stationarity, detailed balance, target atomicity, or invertibility of
the branch.
\end{theorem}

\begin{proof}
By implementation consistency, $\widehat\rho_e$ represents the successful
extended history $e\circ h$ for any $h$ in the class $x$.  By
\Cref{def:conditional-predictive-quotient}, this history belongs to the unique
successor class $y=\delta_e(x)$.  Hence $\widehat\rho_e$ and $\rho_y$ have the
same value on every admissible future effect.  Reducedness separates reachable
states, so $\widehat\rho_e=\rho_y$.  Multiplication by $w_e$ proves the
identity.
\end{proof}

The theorem converts quotient-level unifilarity into an operator identity.  Every
resolved branch carries the canonical reference state of its source to a scalar
multiple of the canonical reference state of its successor.  Iteration therefore
reduces the full branch dynamics on this orbit to a product of one-edge
probabilities, as stated next.

\begin{corollary}[Exact scalar path law]
\label{cor:predictive-reference-path-law}
Let $\gamma=e_n\cdots e_1$ be a successful resolved path with
$x_k=\delta_{e_k}(x_{k-1})$ and
$W_k:=\prod_{j=1}^k w_{e_j}$.  Then
\[
 \boxed{
 \E_{e_n}\cdots\E_{e_1}(\rho_{x_0})
 =W_n\rho_{x_n}.}
\]
In particular, the probability of the resolved path on the canonical
predictive orbit is the product of its one-edge probabilities.
\end{corollary}

\begin{proof}
Induct on the path length using
$\E_{e_k}(\rho_{x_{k-1}})=w_{e_k}\rho_{x_k}$.
\end{proof}

\begin{proposition}[Branchwise Petz anti-renewal on the predictive orbit]
\label{prop:branchwise-petz-reference-renewal}
Let $e:x\to y$ satisfy the reference-renewal identity with faithful support
states,
\[
 \E_e(\rho_x)=w_e\rho_y,
 \qquad 0<w_e\le1.
\]
Let $\E_e^*$ be the trace adjoint and define
\[
 \E_e^{\leftarrow}
 :=
 \Gamma_{\rho_x}\circ\E_e^*\circ\Gamma_{\rho_y}^{-1},
 \qquad
 \Gamma_\rho(a)=\rho^{1/2}a\rho^{1/2}.
\]
Then $\E_e^{\leftarrow}:\A_y\to\A_x$ is a CP trace-nonincreasing branch,
\[
 \Tr\E_e^{\leftarrow}(\rho_y)=w_e,
 \qquad
 (\E_e^{\leftarrow})^*(\one_x)=w_e\one_y,
\]
and it renews the reverse reference state,
\[
 \E_e^{\leftarrow}(\rho_y)=w_e\rho_x
\]
if and only if $\E_e^*(\one_y)=w_e\one_x$.  Reversal is involutive:
\[
 \left(\E_e^{\leftarrow}\right)^{\leftarrow}=\E_e.
\]
For composable branches $\E_e:x\to y$ and $\E_f:y\to z$,
\[
 (\E_f\E_e)^{\leftarrow}
 =\E_e^{\leftarrow}\E_f^{\leftarrow}.
\]
\end{proposition}

\begin{proof}
Complete positivity follows because the three factors in the definition are
CP.  Taking the trace adjoint gives
\[
 (\E_e^{\leftarrow})^*
 =\Gamma_{\rho_y}^{-1}\circ\E_e\circ\Gamma_{\rho_x}.
\]
Hence
\[
 (\E_e^{\leftarrow})^*(\one_x)
 =\rho_y^{-1/2}\E_e(\rho_x)\rho_y^{-1/2}
 =w_e\one_y.
\]
Since $w_e\le1$, this is the success effect of a trace-nonincreasing CP map.
Moreover,
\[
 \Tr\E_e^{\leftarrow}(\rho_y)
 =\Tr\!\left(\rho_x\E_e^*(\one_y)\right)
 =\Tr\E_e(\rho_x)=w_e.
\]
Directly from the definition,
\[
 \E_e^{\leftarrow}(\rho_y)
 =\Gamma_{\rho_x}\E_e^*(\one_y)
 =\rho_x^{1/2}\E_e^*(\one_y)\rho_x^{1/2},
\]
which equals $w_e\rho_x$ exactly when
$\E_e^*(\one_y)=w_e\one_x$, because $\rho_x$ is faithful.  Finally,
inserting
$\Gamma_{\rho_y}^{-1}\Gamma_{\rho_y}=\id$ proves contravariant composition,
and applying the construction twice cancels the two modular congruences and
returns $\E_e$.
\end{proof}

\begin{corollary}[Reciprocal weights from paired Petz branches]
\label{cor:petz-reciprocal-weights}
For every pair of opposite resolved branches
$e:x\to y$ and $\bar e:y\to x$, define the same-orientation retrodictive weight
of $e$ to be the reference-state weight of the Petz reverse of the opposite
branch:
\[
 w_e^-:=
 \Tr\!\left(\E_{\bar e}^{\leftarrow}(\rho_x)\right).
\]
Then
\[
 \boxed{w_e^-=w_{\bar e}^+.}
\]
If the forward branches are reference renewing, their path weights are
multiplicative by \Cref{cor:predictive-reference-path-law}.  If, in addition,
each opposite branch satisfies
\[
 \E_{\bar e}^*(\one_x)=w_{\bar e}^+\one_y,
\]
then
$\E_{\bar e}^{\leftarrow}(\rho_x)=w_{\bar e}^+\rho_y$, so the forward branches
and the reindexed Petz branches form a ray-homogeneous bi-instrument with
$p_e^-=p_{\bar e}^+$.  Without this additional reverse-renewal condition, the
Petz flux still supplies the reciprocal scalar weight, but not a
state-independent retrodictive path law.
\end{corollary}

\begin{proof}
The equality is the reference-state weight identity of
\Cref{prop:branchwise-petz-reference-renewal} applied to the opposite branch.
\end{proof}

\begin{remark}[The exact operational boundary]
A primitive stationary averaged channel does not force branchwise reference
renewal.  The decisive datum is that the target of a resolved branch is the
actual conditional predictive successor represented in a reduced
future-tomographic realization.  If outcomes are forgotten, distinct
successors are replaced by their mixture; if target labels are coarser than
conditional predictive classes, the graph must first be refined.
\end{remark}

The preceding results also locate the exact boundary of the scalar description.
It is valid when target labels already coincide with conditional predictive
successors; an averaged or under-resolved presentation can merge distinct future
states.  The following construction repairs only this labelling defect, without
inventing any unobserved outcome decomposition.

\begin{construction}[Canonical predictive successor refinement]
\label{constr:successor-refinement}
Suppose a resolved instrument is presented on a graph whose target labels may be
coarser than the conditional predictive quotient.  For every reachable pair
$(x,e)$ with
\[
 w_e(x):=\Tr\E_e(\rho_x)>0,
\]
form the normalized branch output
\[
 \widehat\rho_{x,e}:=
 \frac{\E_e(\rho_x)}{\Tr\E_e(\rho_x)}.
\]
Declare
\[
 (x,e)\approx_{\mathrm{succ}}(x',e')
\]
when the two normalized outputs have identical conditional future profiles.
The refined target vertices are the $\approx_{\mathrm{succ}}$-classes of
reachable pairs, and the branch $(x,e)$ is redirected to its own class.
\end{construction}

\begin{proposition}[Minimality of successor refinement]
\label{prop:successor-refinement}
In a reduced future-tomographic realization,
\Cref{constr:successor-refinement} is the coarsest relabelling of the already
resolved branches for which every positive-weight branch is reference renewing.
It is finite whenever the reachable conditional predictive rank is finite.
\end{proposition}

\begin{proof}
All normalized outputs assigned to one refined vertex have the same conditional
future profile.  Reduced future tomography therefore represents them by one and
the same normalized target state.  Consequently every redirected branch
satisfies
\[
 \E_e(\rho_x)=w_e(x)\rho_{[(x,e)]},
\]
so the refined graph is reference renewing.

Conversely, let $r$ be any relabelling for which each positive-weight branch is
reference renewing in a reduced realization.  If
$r(x,e)=r(x',e')$, then the two normalized outputs equal the same target
reference state and hence have identical conditional future profiles.  Thus
$(x,e)\approx_{\mathrm{succ}}(x',e')$.  Every such relabelling therefore factors
uniquely through the quotient of
\Cref{constr:successor-refinement}, proving coarseness and minimality.  Finiteness
is immediate when only finitely many conditional predictive successor classes
are reachable.
\end{proof}

\begin{remark}[Relabelling is not hidden outcome splitting]
\label{rem:refinement-not-splitting}
The construction refines labels of outcomes that are already operationally
resolved.  If only an averaged CP map is available, decomposing it into finer CP
branches is additional instrument data and is not canonically determined by the
channel.  The refinement corrects under-resolved targets; it does not infer an
unobserved Kraus decomposition.
\end{remark}

Canonical refinement gives an exact reference-renewing graph whenever the
conditional predictive quotient is available.  For empirical or approximate
models, exact equality of all future statistics is too strong.  The next theorem
replaces it with a finite separating family and quantifies how local tomographic
errors accumulate along a path.

Exact future tomography may be weakened to a finite quantitative separating
family.

\begin{definition}[Tomographic frame constant]
\label{def:tomographic-frame-constant}
Let $\mathcal Q_y$ be a finite family of effects on the target support algebra.
A number $\kappa_y<\infty$ is a \emph{tomographic frame constant} when
\[
 \norm{\tau-\sigma}_1
 \le
 \kappa_y\max_{Q\in\mathcal Q_y}
 \left|\Tr((\tau-\sigma)Q)\right|
\]
for all reachable normalized states $\tau,\sigma$ at $y$.  Such a constant
exists whenever $\mathcal Q_y$ separates the finite-dimensional affine span of
the reachable state set.
\end{definition}

\begin{theorem}[Quantitative approximate reference renewal]
\label{thm:approximate-reference-renewal}
Let $e:x\to y$ have $w_e>0$, and suppose its normalized output
$\widehat\rho_e=\E_e(\rho_x)/w_e$ obeys
\[
 \max_{Q\in\mathcal Q_y}
 \left|\Tr((\widehat\rho_e-\rho_y)Q)\right|
 \le\varepsilon_e.
\]
Then
\[
 \boxed{
 \norm{\E_e(\rho_x)-w_e\rho_y}_1
 \le w_e\kappa_y\varepsilon_e.}
\]
For a path $\gamma=e_n\cdots e_1$, with
$W_k=\prod_{j=1}^k w_{e_j}$ and $W_0=1$,
\[
 \boxed{
 \norm{
 \E_{e_n}\cdots\E_{e_1}(\rho_{x_0})
 -W_n\rho_{x_n}}_1
 \le
 \sum_{k=1}^n W_k\kappa_{x_k}\varepsilon_{e_k}.}
\]
Consequently approximate conditional predictive sufficiency is stable along
every fixed finite history.
\end{theorem}

\begin{proof}
The frame inequality gives
$\norm{\widehat\rho_e-\rho_y}_1\le\kappa_y\varepsilon_e$; multiplication by
$w_e$ proves the one-edge estimate.

For the path estimate set
\[
 S_k:=\E_{e_k}\cdots\E_{e_1}(\rho_{x_0}),
 \qquad
 D_k:=\norm{S_k-W_k\rho_{x_k}}_1.
\]
A positive trace-nonincreasing map is contractive on Hermitian inputs in trace
norm.  Indeed, if $a=a_+-a_-$ is the Jordan decomposition, then
\[
 \norm{\E(a)}_1
 \le \Tr\E(a_+)+\Tr\E(a_-)
 \le \Tr a_++\Tr a_-=\norm a_1.
\]
Therefore
\[
\begin{aligned}
 D_k
 &\le
 \norm{\E_{e_k}(S_{k-1}-W_{k-1}\rho_{x_{k-1}})}_1\\
 &\quad+
 W_{k-1}
 \norm{\E_{e_k}(\rho_{x_{k-1}})-w_{e_k}\rho_{x_k}}_1\\
 &\le D_{k-1}+W_k\kappa_{x_k}\varepsilon_{e_k}.
\end{aligned}
\]
Iteration from $D_0=0$ gives the result.
\end{proof}

\begin{remark}[Rare resolved outcomes]
The normalized conditional error is independent of $w_e$, whereas the
subnormalized branch error is multiplied by $w_e$.  A rare outcome can
therefore have a poorly estimated conditional state without strongly changing
the unconditioned path law.  Relative statements conditioned on that rare
outcome require separate statistical control.
\end{remark}

We have thus derived both the exact scalar path law and its finite-error version
from conditional prediction.  The matrix formulation is the operational source
of the construction, but later cohomological arguments use only the resulting
one-dimensional action on reference rays.  The next subsection records that
minimal abstract structure.

\subsection{Abstract ray-homogeneous formulation}
\label{subsec:ray-homogeneous}

\Cref{thm:predictive-reference-renewal} derives ray homogeneity on the canonical
predictive reference orbit.  The following ordered-space formulation is useful
when the forward and retrodictive branch families are specified abstractly.
In matrix realizations the retrodictive branches may be taken to be the
branchwise Petz maps of
\Cref{prop:branchwise-petz-reference-renewal}, reindexed in the same path
orientation.

This abstraction deliberately discards all information not needed for path
weights.  It preserves the source and target rays, the positive edge scalars, and
the distinction between forward and retrodictive instruments.  These are exactly
the data from which affinity will be built.

A generic resolved quantum instrument does not assign state-independent
numbers to arbitrary microscopic inputs.  The reference-renewal theorem states
that the necessary one-dimensional closure nevertheless holds on the reduced
conditional predictive orbit.  Abstractly, let $G$ be a finite bidirected graph.  At each vertex $x$, let $X_x$ be an
ordered real vector space with a fixed nonzero positive reference vector
$\eta_x$.

\begin{definition}[Ray-homogeneous bi-instrument]
\label{def:ray-bi-instrument}
A ray-homogeneous bi-instrument consists of positive resolved branch maps
\[
 F_e,R_e:X_{s(e)}\longrightarrow X_{t(e)}
\]
and scalars $p_e^+,p_e^->0$ such that
\[
 F_e\eta_{s(e)}=p_e^+\eta_{t(e)},
 \qquad
 R_e\eta_{s(e)}=p_e^-\eta_{t(e)}.
\]
The family $F$ is the renewal instrument.  The family $R$ is the retrodictive
instrument reindexed in the same path orientation.  When these are
probabilistic branches, their outgoing sums are required to be normalization
preserving.
\end{definition}

\begin{proposition}[Multiplicative path weights]
\label{prop:multiplicative-path-weights}
For a path $\gamma=e_n\cdots e_1$, put
$F_\gamma=F_{e_n}\cdots F_{e_1}$ and similarly for $R_\gamma$.  Then
\[
 F_\gamma\eta_{s(\gamma)}=W_+(\gamma)\eta_{t(\gamma)},
 \qquad
 R_\gamma\eta_{s(\gamma)}=W_-(\gamma)\eta_{t(\gamma)},
\]
where
\[
 \boxed{W_\pm(\gamma)=\prod_{k=1}^n p_{e_k}^\pm.}
\]
In particular the path weights are multiplicative under concatenation.
\end{proposition}

\begin{proof}
Induct on path length.  If the claim holds through $e_{k-1}\cdots e_1$, then
\[
 F_{e_k}\!
 \left(\prod_{j<k}p_{e_j}^+\eta_{s(e_k)}\right)
 =\left(\prod_{j\le k}p_{e_j}^+\right)\eta_{t(e_k)}.
\]
The same argument applies to $R$.  Splitting the product at a concatenation
point gives multiplicativity.
\end{proof}

Multiplicativity is the first payoff of ray homogeneity: a resolved history is
assigned a scalar weight by multiplying local branch factors.  Before those
weights can be attached to predictive path classes, however, one must verify that
they are unchanged by the path equivalences used in the quotient.

\begin{theorem}[Descent of operational path weights]
\label{thm:path-weight-descent}
Let $\equiv$ be any path congruence.  The functions $W_\pm$ descend uniquely to
the quotient path category if and only if they are constant on $\equiv$-classes.
A sufficient condition is
\[
 \gamma\equiv\gamma'
 \quad\Longrightarrow\quad
 F_\gamma\eta_{s(\gamma)}=F_{\gamma'}\eta_{s(\gamma)},
 \quad
 R_\gamma\eta_{s(\gamma)}=R_{\gamma'}\eta_{s(\gamma)}.
\]
In particular, equality of accumulated branch transformations is sufficient.
\end{theorem}

\begin{proof}
A scalar function factors through a quotient exactly when it is constant on
quotient classes.  Under the displayed hypothesis, the two images are equal
nonzero multiples of the same terminal reference vector, so their scalar
coefficients coincide.  Equality of the full accumulated maps implies equality
on the reference ray.
\end{proof}

\begin{proposition}[Kraus invariance and outcome dependence]
\label{prop:kraus-invariance}
Let a resolved quantum branch be represented by the CP map
\[
 \mathcal F_e(\rho)=\sum_aK_{e,a}\rho K_{e,a}^*,
\]
and suppose that on the predictive reference orbit
\[
 \mathcal F_e(\rho_x)=p_e\rho_y,
 \qquad p_e>0.
\]
Then $p_e=\Tr\mathcal F_e(\rho_x)$ depends only on the CP branch map, not on its
Kraus representation.  Hence unitary or isometric Kraus mixing within a fixed
resolved outcome leaves the scalar path weights and the resulting affinity
unchanged.  By contrast, splitting, merging, or physically relabelling outcomes
changes the branch maps and possibly the resolved graph; no invariance under
those operations follows without a separate coarse-graining theorem.
\end{proposition}

\begin{proof}
Two Kraus families related by an isometric change of Kraus section define the
same completely positive map.  Since $\rho_y$ is normalized,
\[
 p_e=\Tr(p_e\rho_y)=\Tr\mathcal F_e(\rho_x),
\]
which is determined solely by $\mathcal F_e$ and the fixed reference input.
Therefore every product of such coefficients along a resolved path, and every
logarithmic ratio formed from them, is Kraus-representation invariant.  Outcome
splitting replaces one CP branch by several CP summands, while merging replaces
several branches by their sum.  These operations alter the instrument rather
than merely its Kraus section, so the preceding invariance argument does not
apply.
\end{proof}

\begin{example}[A primitive averaged channel need not renew branchwise]
\label{ex:average-not-branchwise}
Let $\rho_*=\one_2/2$ and
\[
 \sigma_0=\begin{pmatrix}3/4&0\\0&1/4\end{pmatrix},
 \qquad
 \sigma_1=\begin{pmatrix}1/4&0\\0&3/4\end{pmatrix}.
\]
Define Schr\"odinger CP subchannels
\[
 \mathcal E_i(x)=\frac12\Tr(x)\sigma_i,
 \qquad i=0,1.
\]
Their sum is the one-step primitive replacement channel
$x\mapsto\Tr(x)\rho_*$, but
$\mathcal E_i(\rho_*)=\sigma_i/2$ is not proportional to $\rho_*$.  Thus
stationarity and primitivity of the averaged channel do not imply
ray-homogeneity of its resolved branches.
\end{example}

\begin{remark}[Scope of the scalar reduction]
For a generic instrument acting on arbitrary inputs, the probability of the
next outcome depends on the entire conditional state left by preceding
branches.  Path probabilities then form a cocycle over the conditional-state
bundle rather than a scalar one-cochain on $G$.  The scalar reduction is exact
on the canonical predictive orbit because
\Cref{thm:predictive-reference-renewal} returns every resolved successor to its
uniquely represented predictive state.  Outside that orbit, or when target
labels are under-resolved, the state-bundle description remains necessary.
\end{remark}

The abstract formulation therefore identifies both the strength and the scope of
the scalar reduction.  It is invariant under changes of Kraus representation for
a fixed resolved branch, but not under physical changes of outcome resolution,
and it need not hold away from the canonical predictive orbit.  We next fix an
important convention: which forward/backward path comparison actually measures
irreversibility.

\subsection{Classical Petz path reversal and the affinity convention}
\label{subsec:classical-petz-path-reversal}

The canonical retrodictive chain is constructed to reproduce stationary
forward experiments after reversing the path.  Consequently that particular
forward/reverse ratio cannot measure irreversibility.

The classical calculation makes this point transparently.  Petz reversal is
constructed so that a stationary forward path and the oppositely traversed Petz
path have identical total weight.  A nonzero affinity must therefore compare the
forward and retrodictive laws after placing them in the same orientation.

\begin{proposition}[Stationary Petz path reversal is exact]
\label{prop:stationary-petz-path-reversal}
Let \(P=(p_{ij})\) be a finite row-stochastic matrix with faithful stationary
law \(\pi\), and let
\[
 p^\sharp_{ij}:=\frac{\pi_jp_{ji}}{\pi_i}
\]
be its stationary reverse kernel, the commutative specialization of the
KMS--Petz dual \cite{Petz1986,Petz2003}.  For every path
\(\gamma=(i_0,i_1,\ldots,i_m)\),
\[
 \boxed{
 \pi_{i_0}\prod_{k=1}^m p_{i_{k-1}i_k}
 =
 \pi_{i_m}\prod_{k=1}^m p^\sharp_{i_ki_{k-1}}.}
\]
Therefore the logarithm of stationary forward-path weight divided by the
weight of the reversed path under the Petz chain is identically zero,
whether or not \(P\) satisfies detailed balance.
\end{proposition}

\begin{proof}
For every step,
\[
 p^\sharp_{i_ki_{k-1}}
 =
 \frac{\pi_{i_{k-1}}p_{i_{k-1}i_k}}{\pi_{i_k}}.
\]
Multiplication over \(k\) telescopes all intermediate stationary weights:
\[
 \pi_{i_m}\prod_{k=1}^m p^\sharp_{i_ki_{k-1}}
 =
 \pi_{i_m}
 \prod_{k=1}^m
 \frac{\pi_{i_{k-1}}p_{i_{k-1}i_k}}{\pi_{i_k}}
 =
 \pi_{i_0}\prod_{k=1}^m p_{i_{k-1}i_k}.
\]
\end{proof}

\begin{remark}[Why affinity uses the same orientation]
\label{rem:same-orientation-affinity}
The identity in \Cref{prop:stationary-petz-path-reversal} is adjointness, not
detailed balance.  A genuine cycle obstruction must instead compare forward
and retrodictive branch laws after both have been indexed in the same path
orientation, or compare the two opposite directions under the same forward
law.  This is exactly the convention used in
\Cref{cor:petz-reciprocal-weights,def:affinity}.
\end{remark}

With the convention fixed, the next subsection extracts the genuine continuous
obstruction to reciprocity.  It depends only on ratios of opposite oriented edge
weights and is insensitive to how the record frames themselves are oriented.
\subsection{Real affinity from forward and backward branch weights}

Assume henceforth that the selected resolved branch law is ray homogeneous and
instrument compatible in the sense of
\Cref{def:ray-bi-instrument,thm:path-weight-descent}.  Let
$G=(V,E_G)$ be a finite connected bidirected graph, with reverse edge
$\bar e:y\to x$ for every $e:x\to y$, and let $w_e>0$ denote the resulting
positive scalar capacity of the oriented branch.  In a classical model this may
be a stationary edge flux; in a quantum reset instrument it is the scalar by
which the branch carries the selected reference ray to its conditional
successor ray.

Taking logarithms turns multiplicative path ratios into an additive one-cochain.
Vertex rescalings of the reference rays change this cochain by a coboundary, so
only its cycle class is physically invariant.

\begin{definition}[Edge affinity]
\label{def:affinity}
The edge affinity is the antisymmetric one-cochain
\[
 A(e):=\log\frac{w_e}{w_{\bar e}},
 \qquad A(\bar e)=-A(e).
\]
A vertex gauge $f:V\to\R$ acts by
\[
 A^f(e:x\to y)=A(e)+f(x)-f(y).
\]
The gauge-invariant class is
\[
 [A]\in H^1(G;\R).
\]
\end{definition}

\begin{proposition}[Cycle affinity and detailed balance]
\label{prop:affinity-detailed-balance}
For every oriented cycle $c=e_k\cdots e_1$,
\[
 \oint_c A:=\sum_{j=1}^{k}A(e_j)
 =\log\frac{\prod_jw_{e_j}}{\prod_jw_{\bar e_j}}.
\]
The following are equivalent:
\begin{enumerate}[label=\textup{(\roman*)},leftmargin=2.4em]
\item $[A]=0$ in $H^1(G;\R)$;
\item every cycle affinity vanishes;
\item there is a positive vertex function $\pi$ such that
      $\pi(x)w_e=\pi(y)w_{\bar e}$ for every $e:x\to y$.
\end{enumerate}
\end{proposition}

\begin{proof}
The cycle identity follows by summing logarithms.  If $[A]=0$, then for some vertex function $f$ one has
$A(e:x\to y)=f(y)-f(x)$.  Let $\pi(x)=e^{f(x)}$; then
\[
 \log\frac{\pi(x)w_e}{\pi(y)w_{\bar e}}
 =f(x)+A(e)-f(y)=0.
\]
Thus (i) implies (iii), and (iii) implies vanishing of every cycle product by
telescoping.  If every cycle integral vanishes, fix a base vertex $x_0$ and
define $f(x)$ as the integral of $A$ along any path from $x_0$ to $x$.  Cycle
vanishing makes this path independent, and then $A(e)=f(y)-f(x)$, so $[A]=0$.
\end{proof}

\begin{theorem}[Unique affinity--conductance factorization]
\label{thm:affinity-conductance}
Let $q_e>0$ be positive oriented edge capacities on a bidirected graph.  Set
\[
 c_e:=\sqrt{q_e q_{\bar e}},
 \qquad
 A_q(e):=\log\frac{q_e}{q_{\bar e}}.
\]
Then $c_{\bar e}=c_e$, $A_q(\bar e)=-A_q(e)$, and
\[
 \boxed{q_e=c_e e^{A_q(e)/2}.}
\]
This factorization is unique among a symmetric positive capacity $c$ and an
antisymmetric real one-cochain.  In particular, when $q_e=w_e$ is the resolved
branch capacity of \Cref{def:affinity}, one has $A_q=A$.
\end{theorem}

\begin{proof}
The symmetry and antisymmetry are immediate.  Moreover
\[
 \sqrt{q_e q_{\bar e}}
 \exp\!\left(\frac12\log\frac{q_e}{q_{\bar e}}\right)=q_e.
\]
If $q_e=\tilde c_e e^{\tilde A(e)/2}$ with
$\tilde c_e=\tilde c_{\bar e}>0$ and
$\tilde A(\bar e)=-\tilde A(e)$, then
$q_e q_{\bar e}=\tilde c_e^2$ and
$q_e/q_{\bar e}=e^{\tilde A(e)}$.  Positivity and injectivity of the real
exponential give $\tilde c=c$ and $\tilde A=A_q$.
\end{proof}

\begin{remark}[Affinity does not determine capacity]
The cohomology class $[A]$ fixes only the oriented ratio of the two edge
capacities.  The symmetric conductance $c_e$ is an independent unoriented
renewal-capacity datum.  The sign class introduced below does not enter the
positive factorization.
\end{remark}

The affinity class measures nonreciprocal weight.  It is continuous and real.
The Lorentzian sign class has a different origin.


Affinity has now been decomposed into an antisymmetric log ratio and an
independent symmetric conductance.  Neither datum supplies a sign-changing line
or a two-sheeted cover.  To obtain orientation, the theory needs stable resolved
records and a rule for transporting their contrasts.  The next subsection first
explains how such records can arise canonically from specified microscopic
interactions.

\subsection{Microscopic extraction of a stable resolved instrument}
\label{subsec:stable-pointer-instrument}

The record quotient below begins once a finite raw outcome family is
operationally resolved.  A channel by itself does not determine such an
outcome decomposition.  This subsection gives a conditional microscopic
bridge: a specified system--environment interaction, together with an
interaction-complete monitoring phase, canonically determines a stable
classical pointer algebra and hence a resolved instrument.  The construction
is finite dimensional and basis independent.

This is a conditional bridge rather than a reconstruction from the averaged
channel alone.  The interaction family determines which environment observables
are repeatedly monitored; their commutant contains the nondemolished information,
and the centre of that commutant is the classical part that can label stable
outcomes.

Let
\[
 V:H_{\mathrm{in}}\longrightarrow H_{\mathrm{out}}\otimes E
\]
be a minimal Stinespring isometry for a channel
\[
 \Phi_*(\rho)=\Tr_E(V\rho V^*).
\]
Let \(P\) be a finite-dimensional probe interface and let
\[
 \mathscr H=\{H_1,\ldots,H_q\},
 \qquad H_\nu=H_\nu^*\in\mathcal B(P\otimes E),
\]
be a specified family of microscopic probe--environment interactions.

\begin{definition}[Environment coefficient algebra]
\label{def:environment-coefficient-algebra}
Define
\[
 \mathcal C_E(\mathscr H)
 :=
 C^*\!\left(
 \left\{(\omega\otimes\id_E)(H_\nu):
 \omega\in\mathcal B(P)_*,\ 1\le\nu\le q\right\}
 \right)
 \subseteq\mathcal B(E).
\]
Its commutant
\[
 \mathcal F_E:=\mathcal C_E(\mathscr H)'
\]
is called the \emph{nondemolition stable algebra}, and
\[
 \mathcal Z_E:=Z(\mathcal F_E)
\]
is its stable classical pointer algebra.
\end{definition}

\begin{proposition}[Minimality and nondemolition characterization]
\label{prop:interaction-algebra-minimal}
The algebra \(\mathcal C_E(\mathscr H)\) is the unique smallest unital
\(C^*\)-subalgebra \(\mathcal C\subseteq\mathcal B(E)\) satisfying
\[
 H_\nu\in\mathcal B(P)\otimes\mathcal C
 \qquad(1\le\nu\le q).
\]
In particular it is independent of every probe basis, operator-Schmidt
decomposition, and generating list.  Moreover, for \(X\in\mathcal B(E)\),
\[
 \boxed{
 X\in\mathcal F_E
 \quad\Longleftrightarrow\quad
 [\one_P\otimes X,H_\nu]=0
 \quad\text{for every }\nu.}
\]
\end{proposition}

\begin{proof}
Choose matrix units \((e_{ab})\) of \(\mathcal B(P)\) and write
\[
 H_\nu=\sum_{a,b}e_{ab}\otimes h_{ab}^{(\nu)}.
\]
Each coefficient is a slice,
\(h_{ab}^{(\nu)}=(\omega_{ba}\otimes\id_E)(H_\nu)\), so it belongs to
\(\mathcal C_E(\mathscr H)\), and therefore
\(H_\nu\in\mathcal B(P)\otimes\mathcal C_E(\mathscr H)\).  Conversely, if
\(H_\nu\in\mathcal B(P)\otimes\mathcal C\), every slice of \(H_\nu\) belongs
to \(\mathcal C\).  Hence \(\mathcal C_E(\mathscr H)\subseteq\mathcal C\),
proving minimality and basis independence.

Finally,
\[
 [\one_P\otimes X,H_\nu]
 =\sum_{a,b}e_{ab}\otimes[X,h_{ab}^{(\nu)}].
\]
This vanishes for every \(\nu\) exactly when \(X\) commutes with every slice
coefficient, equivalently with the \(C^*\)-algebra generated by them.
\end{proof}

Equip \(\mathcal B(E)\) with the normalized Hilbert--Schmidt inner product
\[
 \langle X,Y\rangle_2:=\tau_E(X^*Y),
 \qquad \tau_E:=\frac1{\dim E}\Tr.
\]
Let \(\mathcal C_{0,\mathrm{sa}}\) be the orthogonal complement of
\(\R\one_E\) in \(\mathcal C_E(\mathscr H)_{\mathrm{sa}}\), and choose any
real orthonormal basis \(A_1,\ldots,A_m\) of that space.

\begin{theorem}[Interaction-complete stable-pointer selection]
\label{thm:stable-pointer-selection}
Define
\[
 \mathcal L_{\mathscr H}(X)
 :=-\frac12\sum_{j=1}^m[A_j,[A_j,X]].
\]
Then:

\begin{enumerate}[label=\textup{(\roman*)},leftmargin=2.4em]
\item \(\mathcal L_{\mathscr H}\) is independent of the chosen real
orthonormal basis;
\item \(D_t:=e^{t\mathcal L_{\mathscr H}}\) is a norm-continuous semigroup of
unital, trace-preserving, Hilbert--Schmidt self-adjoint completely positive
maps;
\item its fixed algebra is exactly
\[
 \operatorname{Fix}(D)=\mathcal F_E=\mathcal C_E(\mathscr H)';
\]
\item as \(t\to\infty\), \(D_t\) converges in operator norm to the
Hilbert--Schmidt orthogonal projection
\[
 \Delta_{\mathscr H}:\mathcal B(E)\longrightarrow\mathcal F_E;
\]
this limit is the unique trace-preserving conditional expectation onto
\(\mathcal F_E\).
\end{enumerate}
\end{theorem}

\begin{proof}
If \(B_k=\sum_jO_{kj}A_j\) is another real orthonormal basis, then \(O\) is
orthogonal and
\[
 \sum_k[B_k,[B_k,X]]
 =
 \sum_{i,j}\left(\sum_kO_{ki}O_{kj}\right)[A_i,[A_j,X]]
 =
 \sum_i[A_i,[A_i,X]],
\]
proving (i).

For each self-adjoint \(A_j\), the map
\(-\tfrac12[A_j,[A_j,\cdot]]\) is the bounded Lindblad generator
\[
 A_jXA_j-\frac12(A_j^2X+XA_j^2).
\]
Their sum therefore generates a norm-continuous UCP semigroup.  Each summand
annihilates \(\one_E\), and Hilbert--Schmidt cyclicity makes the generator
self-adjoint.  Trace preservation follows either from the Lindblad form or
from \(\Tr[A_j,Y]=0\), proving (ii).

For every \(X\),
\[
 -\langle X,\mathcal L_{\mathscr H}X\rangle_2
 =
 \frac12\sum_j\norm{[A_j,X]}_2^2.
\]
Hence \(\mathcal L_{\mathscr H}X=0\) if and only if
\([A_j,X]=0\) for every \(j\).  Since the \(A_j\), together with the identity,
linearly span the self-adjoint part of \(\mathcal C_E(\mathscr H)\), this is
equivalent to \(X\in\mathcal C_E(\mathscr H)'\), proving (iii).

The generator is self-adjoint and negative semidefinite on the finite
Hilbert space \(\mathcal B(E)\).  Its spectrum is therefore real and
nonpositive, and zero is isolated from the remaining finite spectrum.  The
spectral theorem gives norm convergence of \(e^{t\mathcal L_{\mathscr H}}\)
to the orthogonal projection onto its kernel, namely
\(\Delta_{\mathscr H}\).  A norm limit of UCP maps is UCP, and the limit is
unital, trace preserving, idempotent, and has range \(\mathcal F_E\).
For \(A,B\in\mathcal F_E\), every \(D_t\) is an
\(\mathcal F_E\)-bimodule map because the \(A_j\) commute with \(A,B\);
passing to the limit gives
\[
 \Delta_{\mathscr H}(AXB)=A\Delta_{\mathscr H}(X)B.
\]
Thus \(\Delta_{\mathscr H}\) is a conditional expectation.  If
\(E_1,E_2\) are trace-preserving conditional expectations onto the same
finite-dimensional subalgebra, then for \(Y\in\mathcal F_E\),
\[
 \langle Y,X-E_iX\rangle_2
 =\tau_E(Y^*X)-\tau_E(E_i(Y^*X))
 =0,
\]
using trace preservation and the bimodule property.  Both are therefore the
orthogonal projection onto \(\mathcal F_E\), proving uniqueness.
\end{proof}

\begin{theorem}[Canonical pointer instrument]
\label{thm:canonical-pointer-instrument}
Let \(z_1,\ldots,z_r\) be the minimal central projections of
\(\mathcal F_E\).  They are uniquely determined as an unordered family.
Define CP trace-nonincreasing maps
\[
 \mathcal I_\alpha(\rho)
 :=
 \Tr_E\!\left[
 (\one\otimes z_\alpha)V\rho V^*(\one\otimes z_\alpha)
 \right].
\]
Then
\[
 \sum_{\alpha=1}^r\mathcal I_\alpha=\Phi_*,
\]
so \(\{\mathcal I_\alpha\}\) is a basis-independent resolved instrument.
Its raw classical outcome algebra is canonically
\[
 \mathcal Z_E\cong C(\{1,\ldots,r\};\C).
\]
Applying conditional future equivalence to these pointer outcomes produces
the unique reduced future-distinct record algebra of
\Cref{thm:canonical-record-algebra}.  Reversible monodromy of the resulting
record sectors therefore feeds canonically into the determinant-line
construction.
\end{theorem}

\begin{proof}
The minimal central projections of a finite-dimensional \(C^*\)-algebra are
intrinsic and unique up to permutation.  Each \(\mathcal I_\alpha\) is CP,
being the composition of the isometric dilation, compression by a projection,
and partial trace.  Its trace is at most \(\Tr\rho\), so it is
trace-nonincreasing.

Since the \(z_\alpha\) are pairwise orthogonal and sum to \(\one_E\),
\[
 V\rho V^*
 =
 \sum_{\alpha,\beta}
 (\one\otimes z_\alpha)V\rho V^*(\one\otimes z_\beta).
\]
For \(\alpha\ne\beta\), the partial trace of the off-diagonal block vanishes:
in an orthonormal basis adapted to
\(E=\bigoplus_\alpha z_\alpha E\), every diagonal matrix coefficient of that
block is zero.  Therefore
\[
 \Phi_*(\rho)=\Tr_E(V\rho V^*)
 =\sum_\alpha\mathcal I_\alpha(\rho).
\]
The centre is the span of the \(z_\alpha\), hence is canonically the algebra
of functions on their unordered spectrum.  The final assertions follow by
applying \Cref{def:record-future-equivalence,thm:canonical-record-algebra}
and then the determinant transport results below.
\end{proof}

\begin{remark}[Sharp boundary of pointer extraction]
\label{rem:pointer-extraction-boundary}
The construction is canonical relative to the specified microscopic
interaction family and the interaction-complete monitoring phase.  It is not
canonical relative to the averaged system channel alone.  Different
interaction algebras may have different commutants and hence different stable
pointer centres, even when the reduced channel agrees.  Pure Hamiltonian
motion also need not converge to \(\Delta_{\mathscr H}\); actual monitoring,
phase randomization, or an equivalent irreversible coarse graining is
load-bearing.
\end{remark}


The stable-pointer theorem and canonical instrument theorem supply physically
resolved raw outcomes, but they need not yet be minimal for prediction.  Some
pointer labels may have identical conditional futures.  The next subsection
quotients precisely that redundancy and constructs the canonical algebra on which
record orientation is defined.

\subsection{Canonical predictive record algebra}

The determinant construction should not begin from arbitrarily named record
coordinates.  Once a finite family of raw outcomes is physically resolved,
its minimal future-distinct algebra is fixed by prediction.

The construction mirrors the rooted predictive quotient of
\Cref{sec:predictive}: raw record names are irrelevant, and only distinctions that
can affect future experiments are retained.  This prevents the later determinant
sign from depending on an arbitrary over-refinement of the outcome alphabet.

\begin{definition}[Conditional future equivalence of records]
\label{def:record-future-equivalence}
At a predictive vertex $x$, let $\mathcal O_x$ be a finite set of physically
resolved raw outcome records.  For every admissible future experiment $F$, let
$p_F^x(o)\in[0,1]$ be its conditional probability after observing
$o\in\mathcal O_x$.  Set
\[
 o\equiv_x o'
 \quad\Longleftrightarrow\quad
 p_F^x(o)=p_F^x(o')\ \text{for every admissible }F,
\]
and write $\Omega_x:=\mathcal O_x/{\equiv_x}$ with quotient map $q_x$.
\end{definition}

Define the real unital algebra generated by conditional future predictions,
\[
 \Rr_x^{\pred}:=
 \operatorname{alg}_{\R}\bigl(\one,\{p_F^x:F\text{ admissible}\}\bigr)
 \subseteq C(\mathcal O_x;\R).
\]

\begin{theorem}[Canonical future-distinct record algebra]
\label{thm:canonical-record-algebra}
For every finite resolved record set,
\[
 \boxed{
 \Rr_x^{\pred}=q_x^*C(\Omega_x;\R)
 =\{f\in C(\mathcal O_x;\R):
      o\equiv_x o'\Rightarrow f(o)=f(o')\}.}
\]
Consequently $C(\Omega_x;\R)$ is the unique reduced minimal finite
commutative record algebra through which all conditional future predictions
factor.
\end{theorem}

\begin{proof}
Every generator $p_F^x$ is constant on $\equiv_x$-classes, giving one
inclusion.  For the reverse inclusion, fix a class $C\in\Omega_x$.  For every
$D\ne C$, some future experiment separates $C$ and $D$.  Affinely rescale its
classwise prediction function to obtain $g_D\in\Rr_x^{\pred}$ with
$g_D|_C=1$ and $g_D|_D=0$.  Then
\[
 e_C:=\prod_{D\ne C}g_D
\]
is the indicator of $C$.  The class indicators span $C(\Omega_x;\R)$, proving
the algebra identity.

If $r:\mathcal O_x\to Y$ is any future-sufficient realization, equality
$r(o)=r(o')$ forces $o\equiv_x o'$, so $q_x$ factors uniquely through
$r(\mathcal O_x)$.  If $r$ is reduced, this factor is bijective.  This is the
claimed minimality and uniqueness.
\end{proof}

\begin{proposition}[Descent of reversible record transport]
\label{prop:record-transport-descent}
Suppose a reversible predictive process $u:x\to y$ transports raw records by
a bijection $\tau_u:\mathcal O_x\to\mathcal O_y$ and pulls future experiments
back so that
\[
 p_F^y(\tau_u(o))=p_{u^*F}^x(o).
\]
Then $\tau_u$ descends uniquely to a bijection
$\bar\tau_u:\Omega_x\to\Omega_y$.  Pullback gives a unital order isomorphism
of the canonical record algebras, and composition and inversion descend
functorially.
\end{proposition}

\begin{proof}
If $o\equiv_x o'$, then for every $F$ at $y$,
\[
 p_F^y(\tau_u(o))=p_{u^*F}^x(o)
 =p_{u^*F}^x(o')=p_F^y(\tau_u(o')),
\]
so $\tau_u(o)\equiv_y\tau_u(o')$.  The same argument for $u^{-1}$ proves that
the descended map is bijective.  Functoriality follows from uniqueness of
descent.
\end{proof}

\begin{theorem}[Uniqueness of resolved-record parity]
\label{thm:record-parity-unique}
Let $k\ge2$.  Every group homomorphism
\[
 \zeta:S_k\longrightarrow\{\pm1\}
\]
is either the trivial character or the permutation sign.  Consequently, for
reversible relabellings of $k$ future-distinct records, determinant parity is the
unique nontrivial one-dimensional real functorial invariant.
\end{theorem}

\begin{proof}
All transpositions in $S_k$ are conjugate, so a homomorphism to the abelian group
$\{\pm1\}$ takes one common value on every transposition.  Transpositions
generate $S_k$.  If their common value is $+1$, the homomorphism is trivial.  If
it is $-1$, the value on any permutation is $(-1)^m$, where $m$ is the number of
transpositions in a decomposition.  This is well defined and equals the usual
sign character.  The determinant of the permutation representation is exactly
that character; the invariant constant line has determinant $+1$, so the same
character acts on the contrast module.
\end{proof}

For a finite set $\Omega_x$, its canonical contrast module is
\[
 R_x^0:=\left\{f\in C(\Omega_x;\R):
             \sum_{\omega\in\Omega_x}f(\omega)=0\right\}.
\]
Every descended record bijection preserves $R_x^0$.  Thus the determinant line
of the contrast module is obtained from the conditional predictive quotient,
not from a choice of record names.

\begin{remark}[The remaining instrument boundary]
The theorem reconstructs the minimal record algebra after raw outcomes have
been physically resolved.  An averaged CP channel does not in general select
a preferred outcome decomposition.  The remaining input is therefore the
operational availability of resolved outcomes, not the choice of their
future-sufficient quotient.
\end{remark}

For reversible relabellings, the canonical record quotient carries one unique
nontrivial real one-dimensional invariant: permutation parity on the contrast
module.  Real records, however, may be transported noisily rather than by exact
permutations.  We therefore extend the construction to every positive
mass-preserving correspondence that remains invertible on future-distinct
contrasts.

\subsection{Determinant orientation for noisy record correspondences}
\label{subsec:noisy-record-orientation}

Reversible permutations are not the largest class carrying record orientation.
A noisy channel can blur record values while still preserving the dimension and
orientation of their contrast space.  The determinant detects precisely this
coarser topological information: its magnitude may contract, but its sign remains
defined until a contrast direction is actually lost.
Let
\[
 M_x:=\R^{\Omega_x},
 \qquad
 \mathbf u_x(m):=\sum_{\omega\in\Omega_x}m_\omega,
 \qquad
 M_x^0:=\ker\mathbf u_x.
\]
A positive mass-preserving record correspondence along a chosen orientation of
a geometric edge is a linear map
\[
 K_e:M_x\longrightarrow M_y,
 \qquad
 \mathbf u_yK_e=\mathbf u_x.
\]
It restricts to
\[
 Q_e:=K_e|_{M_x^0}:M_x^0\longrightarrow M_y^0.
\]

\begin{definition}[Determinant-regular noisy record edge]
\label{def:determinant-regular-record-edge}
A noisy record edge is \emph{determinant regular} when
$|\Omega_x|=|\Omega_y|$ and $Q_e$ is a linear isomorphism.  The reverse
transport of the associated real local system is defined algebraically as
$Q_e^{-1}$; it is not required to be positive or to coincide with a physical
reverse stochastic channel.
\end{definition}

\begin{theorem}[Signed cover from determinant-regular noisy records]
\label{thm:noisy-record-orientation}
Let $G$ be a connected graph of constant finite record rank, choose one
orientation of each geometric edge, and assign a determinant-regular positive
mass-preserving record correspondence $K_e$ to that orientation.  Then the
lines
\[
 \mathscr D_x:=\det(M_x^0)
\]
with transports $\det Q_e$ form a real line local system.  In local oriented
generators $d_x\in\mathscr D_x\setminus\{0\}$, write
\[
 (\det Q_e)d_x=a_ed_y,\qquad a_e\in\R^\times,
\]
and put $\chi_{\det}(e)=\sgn(a_e)$.  Then:
\begin{enumerate}[label=\textup{(\roman*)},leftmargin=2.4em]
\item signs multiply under path composition;
\item changing $d_x$ by $\epsilon_xd_x$, $\epsilon_x\in\{\pm1\}$, changes
      $\chi_{\det}$ by a vertex coboundary;
\item the gauge-invariant obstruction is
      \[
       \boxed{
       [\chi_{\det}]
       =w_1(\mathscr D)\in H^1(G;\Ztwo);}
      \]
\item for deterministic record permutations this class agrees with ordinary
      permutation parity on the contrast module.
\end{enumerate}
Thus operational reversibility of the noisy channel is unnecessary; only
linear nondegeneracy of its action on future-distinct record contrasts is
required.
\end{theorem}

\begin{proof}
Mass preservation gives
$K_e(M_x^0)\subseteq M_y^0$.  Determinant regularity makes $Q_e$ invertible, so
its top exterior power is a nonzero map
$\det Q_e:\mathscr D_x\to\mathscr D_y$.  Functoriality of exterior powers gives
\[
 \det(Q_fQ_e)=(\det Q_f)(\det Q_e),
\]
and therefore multiplicativity of the scalar signs along paths.  Replacing
$d_x$ by $\epsilon_xd_x$ changes the edge scalar to
$\epsilon_y^{-1}a_e\epsilon_x$; over $\{\pm1\}$ this is precisely multiplication
by a coboundary.  Hence the sign cocycle represents the first
Stiefel--Whitney class of the determinant-line local system.

For a permutation correspondence, $Q_e$ is the augmentation representation of
the permutation.  The constant line has determinant $+1$, so the determinant
on $M_x^0$ equals the determinant of the full permutation matrix, namely its
ordinary parity.
\end{proof}

\begin{example}[Noisy orientation reversal without a positive inverse]
\label{ex:noisy-binary-record}
Let $\Omega=\{0,1\}$ and
\[
 K_p=
 \begin{pmatrix}
 p&1-p\\
 1-p&p
 \end{pmatrix},
 \qquad 0\le p\le1.
\]
This is a positive bistochastic record channel.  On
$M^0=\R(1,-1)$,
\[
 Q_p(1,-1)=(2p-1)(1,-1).
\]
Hence
\[
 \chi_{\det}(K_p)=
 \begin{cases}
 +1,&p>1/2,\\
 -1,&p<1/2.
 \end{cases}
\]
For $0<p<1$ and $p\ne1/2$, the channel is not a permutation and its linear
inverse is not positive, yet it carries a well-defined orientation sign.  At
$p=1/2$ all record contrast is erased and determinant transport becomes
singular.
\end{example}

\begin{proposition}[Orientation is stable away from rank loss]
\label{prop:noisy-record-sign-stability}
In a continuous family of determinant-regular record correspondences,
$\sgn\det Q_e$ is locally constant.  It can change only at a parameter value
where
\[
 \det Q_e=0,
\]
equivalently where at least one future-distinct record contrast is lost.
\end{proposition}

\begin{proof}
The determinant is continuous and a nonzero real number cannot change sign
without passing through zero.
\end{proof}

The noisy-record theorem shows that orientation survives dissipation whenever
contrast transport remains nonsingular, and that its sign can change only at a
rank-loss wall.  The next subsection packages this local determinant data as a
line system over the graph and identifies the resulting principal double cover.

\subsection{Stable records and determinant orientation}

\begin{definition}[Future-distinct record system]
\label{def:record-system}
A future-distinct record system over $G$ consists of finite-dimensional real
record spaces $R_x$ and invertible linear transports
\[
 L_e:R_x\longrightarrow R_y,
 \qquad e:x\to y,
\]
with $L_{\bar e}=L_e^{-1}$.  The spaces are reduced: two record vectors are
identified when all future record tests agree.  The determinant line is
\[
 \det R_x:=\bigwedge^{\dim R_x}R_x.
\]
\end{definition}

The canonical construction above supplies this datum in the reversible
resolved-record phase by taking $R_x=R_x^0$ and letting $L_e$ be the descended
record transport.  The abstract definition also permits other stable real
record carriers and determinant-regular noisy transports.

At this level the construction no longer depends on a particular pointer model.
Any reduced real record carrier with invertible transport has a determinant line,
and the signs of its transition maps define the same cohomological obstruction.

\begin{definition}[Record-orientation cochain]
\label{def:orientation-cochain}
Assume the record rank is constant on the connected component.  Define
\[
 \chi(e):=\sgn\det L_e\in\{\pm1\}\cong\Ztwo.
\]
Changing the orientation of $R_x$ by $g(x)\in\{\pm1\}$ acts by
\[
 \chi^g(e:x\to y)=g(y)\chi(e)g(x).
\]
The gauge-invariant class is
\[
 [\chi]\in H^1(G;\Ztwo).
\]
\end{definition}

\begin{theorem}[Determinant line and principal signed cover]
\label{thm:determinant-cover}
The determinant transports $\det L_e$ define a real line local system on $G$.
Its first Stiefel--Whitney class is $[\chi]\in H^1(G;\Ztwo)$.  The associated
orientation cover
\[
 G^\chi=(V\times\{\pm1\},E_G^\chi),
 \qquad
 (x,\eta)\xrightarrow{e}(y,\eta\chi(e)),
\]
is a principal $\Ztwo$-cover.  It is trivial if and only if $[\chi]=0$, and it
is connected for every nonzero class.
\end{theorem}

\begin{proof}
The top exterior power is functorial, so the maps $\det L_e$ compose along
paths and define a rank-one real local system.  Choosing an orientation in
each fibre identifies its transition functions with their signs $\chi(e)$.
A change of local orientation multiplies the edge signs by the displayed
coboundary, so the isomorphism class is $[\chi]$, the first
Stiefel--Whitney class of the line system.

The displayed graph is the standard cover associated with the homomorphism
$\pi_1(G)\to\Ztwo$ obtained by multiplying $\chi$ around loops.  The deck action
$(x,\eta)\mapsto(x,-\eta)$ is free and transitive on each fibre.  A global
section exists exactly when the transition signs are a coboundary, equivalently
$[\chi]=0$.  For a connected base, a nonzero homomorphism to $\Ztwo$ is
surjective, so the corresponding two-sheeted cover is connected.
\end{proof}

\begin{corollary}[Minimal real signed enrichment]
\label{cor:minimal-signed-enrichment}
Among local enrichments that preserve the positive record cone, recover the
unsigned base by a deck quotient, and have both positive and negative local
lines, the orientation cover has minimal fibre rank two.  Its canonical
fundamental symmetry is
\[
 \J e_{(x,\eta)}=\eta e_{(x,\eta)}.
\]
Every objectwise-minimal such enrichment is classified, up to orientation
gauge, by $H^1(G;\Ztwo)$.
\end{corollary}

\begin{proof}
Indefiniteness requires at least one positive and one negative local line, so
the fibre rank is at least two.  The orientation cover realizes rank two.
Cone-preserving unitaries permute the two extremal rays, leaving only the
identity and sheet exchange.  Their edgewise choices are precisely
$\Ztwo$-valued transition functions, modulo vertex reorientations.  This gives
the stated cohomology classification and the diagonal sheet-sign operator.
\end{proof}

The determinant line therefore produces the minimal signed enrichment of the
positive predictive graph.  It remains to compare this mod-two obstruction with
the real affinity class derived earlier.  Although both are associated with
cycles, they encode different questions and must not be identified.

\subsection{Affinity and orientation are independent}

Affinity asks whether opposite path weights balance; orientation asks whether a
record frame returns with reversed determinant sign.  The following two theorems
prove independence both constructively, by realizing all four combinations, and
functorially, by ruling out any coefficient homomorphism from real affinity to
mod-two parity.

\begin{theorem}[Independence theorem]
\label{thm:independence}
The real affinity class $[A]\in H^1(G;\R)$ and the determinant orientation
class $[\chi]\in H^1(G;\Ztwo)$ are logically independent.  More precisely, on
every graph with a cycle there exist resolved models realizing each of the four
possibilities
\[
 ([A]=0\text{ or }[A]\ne0)
 \quad\times\quad
 ([\chi]=0\text{ or }[\chi]\ne0).
\]
Consequently detailed balance neither implies nor is implied by orientability
of the record local system.
\end{theorem}

\begin{proof}
It suffices to work on one oriented cycle.  Choose arbitrary positive edge
weights.  Taking $w_e=w_{\bar e}$ gives $A=0$; multiplying one forward weight
by $e^a$ with $a\ne0$ gives a nonzero cycle affinity.  Independently choose
one-dimensional real record transports $L_e=\lambda_e$ with $\lambda_e\ne0$.
If every $\lambda_e>0$, then $[\chi]=0$.  If exactly one edge transport is
negative, the product of signs around the cycle is $-1$, so $[\chi]\ne0$.
The scalar weights and determinant signs can be chosen independently, yielding
all four cases.
\end{proof}

\begin{remark}[Why the distinction matters]
A real affinity records how strongly forward and backward path weights differ.
A determinant sign records whether a future-distinct record frame returns with
its orientation reversed.  The former can vary continuously through zero; the
latter is locally constant on every determinant-regular stratum and changes
only across a rank-loss wall.  The Lorentzian cover is controlled by the
parity class, not by entropy production or nonequilibrium affinity.
\end{remark}

\begin{theorem}[No canonical mod-two reduction of real affinity]
\label{thm:no-affinity-parity-reduction}
Every additive group homomorphism
\[
 f:(\R,+)\longrightarrow\Ztwo
\]
is zero.  Hence no nontrivial map on cohomology induced by a coefficient
homomorphism can send
\[
 H^1(G;\R)\longrightarrow H^1(G;\Ztwo)
\]
and convert the real affinity class into the orientation class.  The magnitude
and parity invariants must arise from distinct components of the operational
data.
\end{theorem}

\begin{proof}
For every $x\in\R$ there is $y\in\R$ with $x=2y$.  Therefore
\[
 f(x)=f(2y)=2f(y)=0
\]
in the exponent-two group $\Ztwo$.  Thus $f$ is identically zero, and functorial
cohomology applied to a coefficient homomorphism produces only the zero map.
\end{proof}

The separation is now complete: entropy production cannot generate the signed
cover, and orientability cannot enforce detailed balance.  The real cover is
already sufficient for an indefinite fundamental symmetry.  A further choice is
needed only if its sign holonomy is to act as a phase on complex amplitudes.

\subsection{Amplitude square roots}

The real signed cover is sufficient for a Krein fundamental symmetry.  If the
same sign holonomy is retained as a phase acting on amplitudes, one must choose
a square root.

\begin{theorem}[Classification and minimality of amplitude lifts]
\label{prop:z4-lift}
Let $G$ be connected.  Identify $\Zfour$ and $\Ztwo$ with the fourth and
second roots of unity when discussing amplitude holonomy, and let
$q:\Zfour\to\Ztwo$ be reduction modulo two.  Then
\[
 q_*:H^1(G;\Zfour)\twoheadrightarrow H^1(G;\Ztwo)
\]
is surjective.  For every $[\chi]\in H^1(G;\Ztwo)$, the fibre
$q_*^{-1}([\chi])$ is a torsor for $H^1(G;\Ztwo)$.  If
$b_1(G)=r$, there are exactly $2^r$ gauge-inequivalent lifts.

If a loop has real holonomy $-1$, every amplitude lift has order-four holonomy
on that loop.  Hence $\Zfour$ is the minimal finite coefficient group that can
contain an amplitude square root of a nontrivial signed sector.  No group
homomorphism $s:\Ztwo\to\Zfour$ satisfies $q\circ s=\id$; therefore the signed
class does not canonically select one of its amplitude lifts.
\end{theorem}

\begin{proof}
Use the short exact sequence
\[
 0\longrightarrow\Ztwo\xrightarrow{\,\iota\,}\Zfour
 \xrightarrow{\,q\,}\Ztwo\longrightarrow0,
 \qquad \iota(1)=2.
\]
The associated cohomology sequence contains
\[
 H^1(G;\Ztwo)\xrightarrow{\iota_*}H^1(G;\Zfour)
 \xrightarrow{q_*}H^1(G;\Ztwo)
 \longrightarrow H^2(G;\Ztwo).
\]
A graph is a one-dimensional CW complex, so $H^2(G;\Ztwo)=0$ and $q_*$ is
surjective.  Since the map on $H^0$ induced by $q$ is surjective, exactness also
makes $\iota_*$ injective.  Thus two lifts of the same class differ freely and
transitively by $H^1(G;\Ztwo)$.  For a connected graph homotopy equivalent to a
wedge of $r$ circles, this group has $2^r$ elements.

If $u(\gamma)^2=-1$, then $u(\gamma)^4=1$ while
$u(\gamma)^2\ne1$, so $u(\gamma)$ has order four.  Conversely the order-four
phase $i$ squares to $-1$.  Finally, if a section $s$ existed, the image $s(1)$ of the order-two
element $1\in\Ztwo$ would have order dividing two.  But
$q(s(1))=1$, so $s(1)$ is an odd class in $\Zfour$ and therefore has order
four, a contradiction.
\end{proof}

\begin{remark}[The second pillar]
The retrodictive dual, the affinity class, and the signed cover are now
separated.  The KMS--Petz dual is selected by a faithful stationary state.  The
real affinity is selected by resolved forward/backward branch weights.  The
orientation class is selected by stable-record monodromy.  Only the last
produces the principal double cover and its fundamental symmetry.
\end{remark}

Section~
ef{sec:retrodictive} has therefore produced two independent cycle
invariants and assigned them distinct roles.  The KMS--Petz construction and
resolved branch weights give the real affinity class, whereas stable record
transport gives the determinant class and its signed cover; amplitude square
roots exist but are not canonically selected.  The next section asks what further
phase properties turn these finite ingredients into modular and Clifford data of
Lorentzian type.

% =====================================================================
\section{The signed modular Lorentzian phase}
\label{sec:phase}

The previous sections reconstruct the finite positive memory and its possible
signed cover.  Lorentzian geometry requires a further recurrent phase: some
predictive transformations must be reversible, renewal depth must acquire a
modular scale, spatial records must have a nondegenerate second moment, and the
translational diamond defect must lie in its exchange-symmetric phase.  These
are phase properties, not universal consequences of complete positivity.

The construction now changes character.  Sections~
ef{sec:predictive} and

ef{sec:retrodictive} identified canonical structures once their operational
domains were specified.  Here we select recurrent components and impose
checkable phase criteria.  Reversible units yield a projective Clifford factor,
the signed cover supplies its temporal row, pressure selects a modular scale, and
record moments supply metric and torsion data.  The closing hierarchy states
exactly which conclusions are structural, phase dependent, or delegated to the
companion continuum theorem.

\subsection{Canonical reductions on recurrent and Perron components}

\begin{proposition}[Component reduction for stochastic and weighted transfers]
\label{prop:terminal-component}
Let $G$ be a finite directed graph and let $C_1,\ldots,C_m$ be its strongly
connected components.
\begin{enumerate}[label=\textup{(\roman*)},leftmargin=2.4em]
\item The condensation graph of the $C_j$ is finite and acyclic.
\item Every matrix $B$ supported on the edges of $G$ is block upper triangular
      after a topological ordering of the components, and
      \[
       \rho(B)=\max_{1\le j\le m}\rho(B|_{C_j}).
      \]
\item If $P$ is stochastic, every stationary probability measure is supported
      on the union of the closed communicating classes, equivalently the
      terminal components of the positive-probability support graph.
\item Under a stationary Markov law, almost every trajectory lies in one closed
      recurrent class after a finite random time.  This statement is not true
      for arbitrary deterministic infinite paths in the underlying graph.
\end{enumerate}
Consequently pressure analysis is performed on an irreducible component block
attaining the maximal Perron radius, whereas a stationary stochastic phase is
reduced to a selected closed recurrent class.
\end{proposition}

\begin{proof}
Collapsing each strongly connected component gives a directed acyclic graph;
otherwise a directed cycle among components would make their union strongly
connected.  A topological order therefore makes every supported matrix block
upper triangular.  The spectrum of a block upper triangular matrix is the union
of the spectra of its diagonal blocks, proving (ii).

Let $P$ be stochastic and let $\pi P=\pi$.  In the finite condensation graph,
a nonterminal component has a directed path to a terminal component.  Summing
the stationarity equation over a downward-closed set and iterating along the
acyclic order shows that a stationary measure assigns zero mass to every
transient component; equivalently, finite-state Markov-chain recurrence theory
places all stationary mass on closed communicating classes
\cite{Seneta2006}.  Starting from a stationary law, the chain is in the recurrent
support almost surely and visits transient states only finitely often, which
proves (iv).  A deterministic path may ignore an available exit forever, so no
pathwise statement holds without probabilistic qualification.
\end{proof}

Global strong connectivity is therefore unnecessary.  The relevant reduction
must, however, match the theorem being used: closed recurrent classes for
stationary Markov dynamics and Perron-maximizing irreducible blocks for general
weighted pressure transfers.

The component reduction prevents two distinct notions of recurrence from being
conflated.  Stationary stochastic arguments belong on closed communicating
classes, whereas pressure belongs on irreducible blocks attaining the maximal
Perron radius.  After making the appropriate reduction, we can examine the
reversible part of the renewal dynamics and the algebra generated by its
projective implementers.

\subsection{Reversible units and the primitive Clifford factor}

Let $K_{\rev}$ be a finite elementary abelian two-group of realized reversible
labels acting on a finite recurrent predictive algebra.  By
\Cref{cor:units-automorphisms}, every label acts by a $*$-automorphism.  After
selecting one invariant central orbit, quotienting the action kernel, and
choosing a simple block, we obtain a faithful action on a matrix factor
$M_N(\C)$.

Automorphisms of a matrix factor are inner, but their unitary implementers are
defined only up to phase.  This unavoidable phase ambiguity is not a defect: its
commutator records the alternating binary form from which the Clifford algebra
will emerge.

\begin{theorem}[Automatic projective lift]
\label{thm:projective-lift}
Choose unitary implementers $R_\alpha$ of the automorphisms,
\[
 \Phi_\alpha=\Ad(R_\alpha),
 \qquad \alpha\in K_{\rev},
 \qquad R_0=\one.
\]
Then there is a normalized scalar two-cocycle
$\sigma\in Z^2(K_{\rev},U(1))$ such that
\[
 R_\alpha R_\beta
 =\sigma(\alpha,\beta)R_{\alpha+\beta}.
\]
Its antisymmetric ratio is a gauge-invariant alternating form
\[
 \sigma(\alpha,\beta)\sigma(\beta,\alpha)^{-1}
 =(-1)^{\omega_{\spat}(\alpha,\beta)},
 \qquad
 \omega_{\spat}\in\Alt^2(K_{\rev},\F_2).
\]
\end{theorem}

\begin{proof}
Every automorphism of a complex matrix factor is inner.  The unitaries
$R_\alpha R_\beta$ and $R_{\alpha+\beta}$ implement the same automorphism,
so their ratio lies in the centre and is a scalar phase.  Associativity of the
unitary product is exactly the two-cocycle identity for $\sigma$.  Replacing
$R_\alpha$ by $\mu(\alpha)R_\alpha$ changes $\sigma$ by the coboundary
$\delta\mu$.

Because $K_{\rev}$ is abelian, the antisymmetric ratio is a bicharacter.
Because $2\alpha=0$, its square is one, so it takes values in $\{\pm1\}$ and
defines an $\F_2$-valued alternating form.  Coboundaries are symmetric on an
abelian group and cancel from the ratio, proving gauge invariance.
\end{proof}

The automatic projective lift extracts the spatial commutator data from the
reversible predictive action alone.  It has not yet identified which generator
should be timelike.  That information comes from the signed cover: sheet holonomy
determines how the fundamental symmetry commutes with every reversible cycle.

The signed cover supplies a temporal generator.  Let $J$ be the sheet-sign
fundamental symmetry and suppose a reversible cycle class $\alpha$ shifts the
signed sheet by the holonomy pairing with $[\chi]$.

\begin{theorem}[Canonical temporal row]
\label{thm:temporal-row}
For every realized reversible cycle class $\alpha\in K_{\rev}$,
\[
 JR_\alpha
 =(-1)^{\langle[\chi],\alpha\rangle}R_\alpha J.
\]
Adjoin a temporal label $t$ with $R_t=J$ and set
$K_{\mathrm{full}}=K_{\rev}\oplus\langle t\rangle$.  The full commutator form
is
\[
 \omega((\alpha,\epsilon),(\beta,\delta))
 =\omega_{\spat}(\alpha,\beta)
  +\epsilon\langle[\chi],\beta\rangle
  +\delta\langle[\chi],\alpha\rangle.
\]
Hence the temporal row is fixed by the signed cover and is not an independent
projective input.
\end{theorem}

\begin{proof}
A lifted cycle beginning at $(x,\eta)$ ends at
$(x,\eta(-1)^{\langle[\chi],\alpha\rangle})$.  Since $J$ multiplies a basis
state by its sheet sign, moving $J$ through the lifted revision produces the
displayed scalar.  Define
$R_{\alpha+\epsilon t}:=R_\alpha J^\epsilon$ and move $J^\epsilon$ through
$R_\beta$.  Combining the resulting sign with the spatial projective cocycle
of \Cref{thm:projective-lift} gives the formula for the full antisymmetric
commutator.
\end{proof}

The temporal-row theorem joins the spatial projective form and the determinant
holonomy into one alternating form on the enlarged label group.  Some labels may
still commute with everything and therefore contribute only central redundancy.
Quotienting this radical leaves the primitive nondegenerate sector.

\begin{theorem}[Radical reduction and primitive Clifford factor]
\label{thm:clifford-factor}
Let
\[
 \rad\omega:=\{a\in K_{\mathrm{full}}:\omega(a,b)=0
                 \text{ for all }b\in K_{\mathrm{full}}\}
\]
and let $K_{\mathrm{prim}}=K_{\mathrm{full}}/\rad\omega$.  If
$\dim_{\F_2}K_{\mathrm{prim}}=2m$, then the noncentral factor of the twisted
group algebra is
\[
 \C_\sigma[K_{\mathrm{prim}}]
 \cong M_{2^m}(\C)
 \cong\Cl_{2m}(\C).
\]
The faithful irreducible projective representation is unique up to unitary
equivalence.
\end{theorem}

\begin{proof}
The induced alternating form on $K_{\mathrm{prim}}$ is nondegenerate by
construction and therefore has even dimension.  Choose a symplectic basis
$(p_1,q_1,\ldots,p_m,q_m)$.  After scalar rephasing, the corresponding
implementers satisfy
\[
 R_{p_j}R_{q_j}=-R_{q_j}R_{p_j},
\]
while generators belonging to different pairs commute.  Each pair generates
$M_2(\C)$, and the full twisted algebra is their tensor product
$M_2(\C)^{\otimes m}=M_{2^m}(\C)$.  This is the complex Clifford algebra on
$2m$ generators.  The finite Stone--von Neumann theorem gives uniqueness of
the irreducible representation of the nondegenerate finite Heisenberg
extension.
\end{proof}

The radical quotient identifies the finite matrix factor generated by the
noncommuting revisions.  When the determinant class pairs nontrivially with a
reversible cycle, the abstract temporal row also produces a concrete operator of
negative square, as the following corollary shows.

\begin{corollary}[Intrinsic timelike generator]
\label{cor:timelike-generator}
If some $\alpha_0\in K_{\rev}$ satisfies
$\langle[\chi],\alpha_0\rangle=1$, normalize $J$ and $R_{\alpha_0}$ as
Hermitian involutions.  Then
\[
 \gamma^0:=JR_{\alpha_0}
\]
is anti-Hermitian and
\[
 (\gamma^0)^2=-\one.
\]
\end{corollary}

\begin{proof}
The temporal-row relation gives
$JR_{\alpha_0}=-R_{\alpha_0}J$.  Therefore
\[
 (JR_{\alpha_0})^*=R_{\alpha_0}J=-JR_{\alpha_0},
 \qquad
 (JR_{\alpha_0})^2=-J^2R_{\alpha_0}^2=-\one.
\]
\end{proof}

A nontrivial signed cover and a timelike Clifford generator establish the
\emph{possibility} of temporal orientation.  They do not choose one of the two
sheets, nor do they guarantee that future experiments can distinguish the choice.
The next subsection separates these three logically different levels.

\subsection{Selection and predictive visibility of temporal orientation}

Let $\vartheta$ denote the deck automorphism on the signed observable algebra;
thus $\vartheta(J)=-J$.  An observable $a$ is deck odd when
$\vartheta(a)=-a$.

The first result is a no-go theorem: finite primitive dynamics that respects deck
symmetry selects the symmetric state, not an orientation.  The second result asks
a different question---whether, once an oriented state has been selected by some
additional mechanism, the predictive test algebra is rich enough to observe it.

\begin{theorem}[Finite primitive symmetric phases cannot select an orientation]
\label{thm:no-finite-orientation-selection}
Let $T$ be a primitive finite-dimensional renewal transfer on the signed state
space and assume deck covariance,
\[
 T\circ\vartheta_*=\vartheta_*\circ T.
\]
Then its unique stationary state $\rho_*$ is deck invariant and
\[
 \Tr(\rho_*a)=0
\]
for every deck-odd observable $a$, in particular for $a=J$.
\end{theorem}

\begin{proof}
By \Cref{thm:primitive-stationary-state}, $T$ has a unique stationary density.
Deck covariance makes $\vartheta_*(\rho_*)$ stationary as well, so uniqueness
gives $\vartheta_*(\rho_*)=\rho_*$.  For deck-odd $a$,
\[
 \Tr(\rho_*a)=\Tr(\vartheta_*(\rho_*)a)
 =\Tr(\rho_*\vartheta(a))=-\Tr(\rho_*a),
\]
and the expectation vanishes.
\end{proof}

\begin{corollary}[Possible orientation-selection mechanisms]
\label{cor:orientation-selection-mechanisms}
A nonzero stationary deck-odd order parameter requires at least one departure
from a finite primitive deck-symmetric phase: explicit deck-symmetry breaking,
stationary-state degeneracy, boundary or superselection data, or a singular
infinite-volume/infinite-depth limit in which uniqueness fails.
\end{corollary}

\begin{theorem}[Deck-odd observability criterion]
\label{thm:deck-observability}
Let $\omega_-:=\omega_+\circ\vartheta$.  On the deck-even algebra
$\A^\vartheta=\{a:\vartheta(a)=a\}$ one has
\[
 \omega_-|_{\A^\vartheta}=\omega_+|_{\A^\vartheta}.
\]
Hence a predictive quotient using only deck-even future tests identifies the
two orientations.  Conversely, if an admissible deck-odd observable $a$
satisfies $\omega_+(a)\ne0$, then
\[
 \omega_-(a)=-\omega_+(a),
\]
so the two phases are predictively distinguishable.
\end{theorem}

\begin{proof}
For deck-even $a$,
$\omega_-(a)=\omega_+(\vartheta(a))=\omega_+(a)$.  For deck-odd $a$, the same
identity gives $\omega_-(a)=-\omega_+(a)$.  A single admissible test with
different expectations separates the predictive classes.
\end{proof}

\begin{remark}[Three distinct levels]
A nontrivial class $[\chi]$ gives the possibility of temporal orientation; a
state-level phase chooses one sheet; a deck-odd admissible test makes that
choice operationally visible.  None of these statements implies either of the
other two.
\end{remark}

We have now distinguished topology, state selection, and observability.  The
signed cover supplies a binary temporal direction, but not a rate at which renewal
depth should be read as time.  The next subsection derives a dimensionless
modular scale from Perron--Frobenius pressure and realizes it on an unbounded
graded path space.

\subsection{Perron pressure and the affine modular clock}

Let the selected recurrent component carry positive branch capacities
$q_e>0$ and positive depths $\ell_e\ge\ell_0>0$.  Define the weighted transfer
matrix
\[
 B(s)_{xy}:=\sum_{e:x\to y}q_e e^{-s\ell_e}.
\]
When these capacities come from the resolved bidirected branch law, write
$q_e=c_e e^{A(e)/2}$ as in \Cref{thm:affinity-conductance}.

There are four steps.  Gauge covariance first shows that the pressure is
independent of local normalizations.  Supercriticality then selects a unique root
$\beta$.  The associated Doob transform gives a stationary edge process whose
mean depth and entropy production admit an exact identity.  Finally, a graded path
cover turns the scalar root into an affine modular flow without contradicting the
finite-cycle obstruction.

\begin{proposition}[Gauge covariance of pressure and the Doob law]
\label{prop:pressure-gauge-invariance}
Let a vertex gauge act by
$A^f(e:x\to y)=A(e)+f(x)-f(y)$ and put
$q_e^f=c_e e^{A^f(e)/2}$.  With
$D_f=\operatorname{diag}(e^{f(x)/2})$,
\[
 B^f(s)=D_f B(s)D_f^{-1}.
\]
Consequently the pressure $\log\rho(B(s))$, its zero $\beta$, and its
derivatives are gauge invariant.  If $h,\nu$ are Perron vectors at the root,
then choosing $h^f=D_f h$ and $(\nu^f)^T=\nu^T D_f^{-1}$ leaves the Doob edge
law
\[
 p_e=q_e e^{-\beta\ell_e}\frac{h_y}{h_x}
\]
and the stationary vertex law $\pi_x=\nu_x h_x$ unchanged.
\end{proposition}

\begin{proof}
For an edge $e:x\to y$,
$q_e^f=e^{f(x)/2}q_e e^{-f(y)/2}$, which gives the diagonal similarity after
summing parallel edges.  Similarity preserves the spectrum and hence the
pressure data.  Substitution of the transformed Perron vectors gives
$p_e^f=p_e$ and $\pi_x^f=\pi_x$.
\end{proof}

The similarity relation shows that pressure depends only on gauge-invariant cycle
data, even though individual oriented capacities depend on local choices.  We can
therefore use its spectral radius to select an intrinsic exponent.

\begin{theorem}[Pressure-root selection]
\label{thm:pressure-root}
Assume the recurrent component is strongly connected and let
$r(s)=\rho(B(s))$.  Then $r$ is continuous and strictly decreasing on
$[0,\infty)$, with $r(s)\to0$ as $s\to\infty$.  Consequently:
\begin{enumerate}[label=\textup{(\roman*)},leftmargin=2.4em]
\item if $r(0)>1$, there is a unique $\beta>0$ with $r(\beta)=1$;
\item if $r(0)=1$, the unique nonnegative root is $\beta=0$;
\item if $r(0)<1$, no nonnegative root exists.
\end{enumerate}
At the root the Perron eigenvalue is simple and has strictly positive left and
right eigenvectors.
\end{theorem}

\begin{proof}
For $s_2>s_1$, every nonzero matrix entry satisfies
$B(s_2)_{xy}<B(s_1)_{xy}$.  Irreducibility and the strict monotonicity part of
Perron--Frobenius imply $r(s_2)<r(s_1)$.  Continuity follows from continuity of
the finite matrix entries and of the spectral radius.  The positive lower
bound on $\ell_e$ makes every entry tend to zero, so $r(s)\to0$.  The three
alternatives follow from the intermediate value theorem and strict
monotonicity.  Simplicity and positivity of the Perron eigenvectors are the
irreducible Perron--Frobenius theorem.
\end{proof}

The pressure-root theorem identifies exactly when a positive modular exponent
exists: it is a supercritical phase property, not a universal consequence of
recurrence.  At the root, the Perron vectors canonically normalize the weighted
transfer into a stationary stochastic law.

\begin{theorem}[Stationary pressure law and mean depth]
\label{constr:pressure-law}
Let $h>0$ and $\nu>0$ be right and left Perron eigenvectors at $\beta$,
normalized by $\nu^Th=1$.  For an edge $e:x\to y$, define
\[
 p_e:=q_e e^{-\beta\ell_e}\frac{h_y}{h_x},
 \qquad \pi_x:=\nu_xh_x.
\]
Then $p$ is a stochastic edge law and $\pi$ is stationary.  If
$P(s):=\log r(s)$, its stationary mean depth satisfies
\[
 \boxed{
 \mu_\ell:=\sum_e\pi_{s(e)}p_e\ell_e=-P'(\beta)>0.}
\]
\end{theorem}

\begin{proof}
The row normalization is
\[
 \sum_{e:s(e)=x}p_e=h_x^{-1}(B(\beta)h)_x=1.
\]
For a vertex $y$,
\[
 \sum_{e:t(e)=y}\pi_{s(e)}p_e
 =h_y\sum_x\nu_xB(\beta)_{xy}=\nu_yh_y=\pi_y.
\]
The Perron root is simple, so the standard simple-eigenvalue perturbation
formula gives $r'(\beta)=\nu^TB'(\beta)h$; see, for example,
\cite{Seneta2006}.  Since $r(\beta)=1$,
\[
 -P'(\beta)=-r'(\beta)
 =\sum_e\nu_{s(e)}q_e e^{-\beta\ell_e}h_{t(e)}\ell_e
 =\sum_e\pi_{s(e)}p_e\ell_e.
\]
Positivity follows from $\ell_e>0$.
\end{proof}

The Doob law gives the probabilistic process selected by the pressure root, and
the derivative of pressure gives its average renewal depth.  To interpret
$\beta$, we now decompose each capacity into symmetric conductance and affinity
and compare the resulting stationary edge fluxes.

The pressure root also relates the modular scale to the information carried by
the stationary edge process.  Assume in the remainder of this subsection that
the selected edge set is bidirected, $\ell_{\bar e}=\ell_e$, and
\[
 q_e=c_e e^{A(e)/2},
 \qquad c_{\bar e}=c_e>0,
 \qquad A(\bar e)=-A(e),
\]
as in \Cref{thm:affinity-conductance}.

\begin{definition}[Stationary entropy production and entropy rate]
\label{def:pressure-entropy-data}
Let
\[
 F_e:=\pi_{s(e)}p_e.
\]
Then $(F_e)_e$ and $(F_{\bar e})_e$ are probability distributions on the
oriented edge set.  Define
\[
 \mathscr A(e):=\log\frac{F_e}{F_{\bar e}},
 \qquad
 \sigma:=\sum_eF_e\mathscr A(e),
\]
and
\[
 h_{\mathrm{rate}}:=-\sum_eF_e\log p_e,
 \qquad
 \kappa_c:=\sum_eF_e\log c_e.
\]
The quantity $\sigma$ is the relative entropy
$D_{\mathrm{KL}}(F\|F\circ\bar{\phantom e})$ and is therefore nonnegative.
\end{definition}

\begin{lemma}[Input affinity and stationary affinity are cohomologous]
\label{lem:stationary-affinity-cohomology}
Let $h,\nu$ be the Perron vectors of
\Cref{constr:pressure-law}, and set
\[
 v(x):=\log\frac{h_x}{\nu_x}.
\]
Then
\[
 \mathscr A(e)=A(e)+v(t(e))-v(s(e)).
\]
Consequently
\[
 \boxed{
 \sum_eF_eA(e)=\sum_eF_e\mathscr A(e)=\sigma\ge0.}
\]
\end{lemma}

\begin{proof}
For $e:x\to y$,
\[
 F_e
 =\pi_xq_e e^{-\beta\ell_e}\frac{h_y}{h_x}
 =\nu_xq_e e^{-\beta\ell_e}h_y.
\]
Using $\ell_{\bar e}=\ell_e$ gives
\[
 \log\frac{F_e}{F_{\bar e}}
 =\log\frac{q_e}{q_{\bar e}}
  +\log\frac{\nu_xh_y}{\nu_yh_x}
 =A(e)+v(y)-v(x).
\]
Stationarity implies that the expectation of every vertex coboundary vanishes:
\[
 \sum_eF_e(v(t(e))-v(s(e)))
 =\sum_y\pi_yv(y)-\sum_x\pi_xv(x)=0.
\]
Hence the two stationary means agree.  Finally,
\[
 \sigma
 =\sum_eF_e\log\frac{F_e}{F_{\bar e}}
 =D_{\mathrm{KL}}(F\|F\circ\bar{\phantom e})\ge0
\]
by the log-sum inequality.
\end{proof}

The preceding lemma shows that the input affinity and the affinity measured by
the stationary edge process differ only by a vertex coboundary.  Their stationary
means therefore agree and equal the nonnegative entropy production.  Substituting
this fact into the Doob normalization yields the central balance law.

\begin{theorem}[Entropy--affinity--depth identity]
\label{thm:entropy-affinity-depth}
Under the preceding bidirected pressure hypotheses,
\[
 \boxed{
 \beta\mu_\ell
 =h_{\mathrm{rate}}+\kappa_c+\frac12\sigma.}
\]
Equivalently,
\[
 \boxed{
 \beta
 =
 \frac{
 h_{\mathrm{rate}}
 +\mathbb E_F[\log c_e]
 +\tfrac12\mathbb E_F[A(e)]}
 {\mu_\ell}.}
\]
In the neutral-conductance normalization $c_e\equiv1$,
\[
 \boxed{
 \beta=\frac{h_{\mathrm{rate}}+\sigma/2}{\mu_\ell}.}
\]
Thus nonequilibrium affinity contributes a positive one-half correction to the
entropy-per-depth law, while the symmetric conductance remains an independent
unoriented capacity datum.
\end{theorem}

\begin{proof}
Taking logarithms of the Doob law gives
\[
 \log p_e
 =\log q_e-\beta\ell_e
  +\log h_{t(e)}-\log h_{s(e)}.
\]
Average with respect to $F$.  By stationarity, the final coboundary has zero
mean, so
\[
 \beta\mu_\ell
 =-\sum_eF_e\log p_e+\sum_eF_e\log q_e
 =h_{\mathrm{rate}}+\sum_eF_e\log q_e.
\]
The unique affinity--conductance factorization yields
\[
 \log q_e=\log c_e+\frac12A(e).
\]
Substitution and \Cref{lem:stationary-affinity-cohomology} prove the first
identity.  Division by $\mu_\ell>0$ gives the second, and $c_e\equiv1$ gives
the final specialization.
\end{proof}

\begin{remark}[Why the factor one half is forced]
The affinity is the full oriented log ratio
$A(e)=\log(q_e/q_{\bar e})$.  A single oriented pressure capacity equals the
geometric mean $c_e$ times the square root of that ratio.  The coefficient
$1/2$ is therefore algebraic, not conventional.
\end{remark}

The identity explains what fixes the modular exponent: predictive entropy,
symmetric renewal capacity, and nonequilibrium circulation per unit mean depth.
It still remains to realize that exponent as a consistent flow.  A finite return
cycle cannot carry the same nonzero eigenvalue on every edge, so the clock must
remember accumulated path depth.

\begin{construction}[Graded path-clock cover]
\label{constr:path-clock}
Let $\mathcal P_{\mathrm{fin}}(G)$ be the set of finite directed paths in the
selected recurrent graph and put
\[
 \Hh_{\mathrm{clk}}:=\ell^2(\mathcal P_{\mathrm{fin}}(G)).
\]
For a path $w=e_1\cdots e_k$, set
\[
 L(w):=\sum_{j=1}^k\ell_{e_j},
 \qquad N\delta_w=L(w)\delta_w.
\]
If appending $e$ to $w$ is composable, define $S_e\delta_w=\delta_{we}$, and
set $S_e\delta_w=0$ otherwise.  Then $[N,S_e]=\ell_eS_e$ on the finite-path
core.
\end{construction}

\begin{theorem}[Pressure-selected affine path clock]
\label{thm:pressure-path-clock}
Let $\beta$ be the pressure root of \Cref{thm:pressure-root} and define the
positive self-adjoint operator
\[
 \Delta_\beta:=e^{\beta N}
\]
on the graded path-clock cover.  Then every path-creation operator is a modular
eigenoperator:
\[
 \boxed{
 \Delta_\beta^{it}S_e\Delta_\beta^{-it}
 =e^{it\beta\ell_e}S_e.}
\]
Thus the pressure root determines an affine modular clock without assigning a
nonzero common exponent directly to every edge of the finite recurrent base.
\end{theorem}

\begin{proof}
For every basis vector for which $we$ is defined,
\[
 \Delta_\beta^{it}S_e\Delta_\beta^{-it}\delta_w
 =e^{-it\beta L(w)}e^{it\beta L(we)}\delta_{we}
 =e^{it\beta\ell_e}S_e\delta_w.
\]
Both sides vanish when $we$ is not composable.  The finite-path span is a core,
so the identity extends to the bounded creation operators.
\end{proof}

The graded cover resolves the finite-cycle obstruction by separating the finite
recurrent base from the unbounded depth coordinate.  Every edge creation operator
then has the pressure-selected modular frequency.  The next theorem shows that,
in an irreducible realization, no different compatible Hamiltonian is possible
except for an additive scalar.

\begin{theorem}[Rigidity of a compatible modular Hamiltonian]
\label{thm:affine-clock}
Let $N$ and $S_e$ satisfy $[N,S_e]=\ell_eS_e$.  If a faithful modular flow on
a selected factor obeys
\[
 \sigma_t(S_e)=e^{it\beta\ell_e}S_e
\]
and the joint commutant of the elementary resets is scalar, then
\[
 H_{\mathrm{mod}}=\beta N+c\one.
\]
Thus every compatible irreducible realization of the pressure-selected clock
is affine in renewal depth, up to an additive scalar.
\end{theorem}

\begin{proof}
Writing $\sigma_t=\Ad(e^{itH_{\mathrm{mod}}})$ and differentiating at zero gives
$[H_{\mathrm{mod}}-\beta N,S_e]=0$ for every $e$.  The scalar-commutant
hypothesis gives $H_{\mathrm{mod}}-\beta N=c\one$.
\end{proof}

\begin{proposition}[Finite-cycle obstruction]
\label{prop:finite-cycle-clock-obstruction}
Suppose $e_1,\ldots,e_m$ form a closed reset cycle, each $S_{e_j}$ is a modular
eigenoperator with exponent $\beta_j$, and
\[
 S_{e_m}\cdots S_{e_1}=cP,
 \qquad c\ne0,
\]
where $P$ is a modular-invariant one-dimensional corner.  Then
$\sum_j\beta_j=0$.  In particular, a common exponent on such a cycle must be
zero.
\end{proposition}

\begin{proof}
The product is a modular eigenoperator with exponent $\sum_j\beta_j$, whereas
$cP$ is modular invariant.  Since the product is nonzero,
$e^{it\sum_j\beta_j}=1$ for every $t$, so the sum vanishes.
\end{proof}

\begin{remark}[Critical state versus trace density]
At $r(\beta)=1$, the finite-path partition series is critical and need not
define a trace-class density on $\Hh_{\mathrm{clk}}$.  The natural normalized
object is the Perron boundary Markov state of
\Cref{constr:pressure-law}, or a faithful semifinite modular weight.  For
$s>\beta$, the weighted path series is subcritical and convergent.
\end{remark}

The pressure theorem selects a dimensionless modular exponent and the graded
cover realizes its flow.  Converting the resulting parameter into seconds
requires a physical unit calibration and is not part of the dimensionless
reconstruction.

Pressure has supplied an intrinsic dimensionless clock, while the signed and
projective structures supplied its temporal algebraic direction.  Lorentzian data
also require a nondegenerate spatial quadratic form and compatible comparison
between neighbouring frames.  These arise next from second moments of stable
direction records.

\subsection{Spatial second moments and metric comparison}

Let $V_{\spat}\cong\R^d$ be the real span of direction records on the selected
recurrent component.  A local stationary direction law $\nu_x$ has second
moment
\[
 M_x:=\int_{V_{\spat}}vv^{\mathsf T}\,\dd\nu_x(v).
\]

The second moment is the lowest-order statistic that can define a positive
quadratic geometry without assuming coordinates.  Symmetry may force it to be
isotropic, but only positive definiteness is needed; covariance under transport
then determines the orthogonal comparison maps.

\begin{proposition}[Symmetry forces ellipticity]
\label{prop:symmetry-ellipticity}
Let a compact group $K$ act orthogonally and irreducibly on $V_{\spat}$.  If
$\nu$ is a finite positive $K$-invariant measure, not supported at the origin,
and has finite second moment, then
\[
 M=\int vv^{\mathsf T}\,\dd\nu(v)=cI
\]
for a scalar $c>0$.
\end{proposition}

\begin{proof}
Invariance gives $kMk^{-1}=M$ for every $k\in K$.  The matrix $M$ is
self-adjoint and positive.  By real Schur theory, the self-adjoint part of the
commutant of an irreducible orthogonal representation is scalar: any additional
complex or quaternionic generators are skew-adjoint.  Thus $M=cI$.  Since
$\Tr M=\int\norm{v}^2\dd\nu(v)>0$, one has $c>0$.
\end{proof}

Literal isotropy is not required.  Positive definiteness of $M_x$ is enough to
define a coframe; constant anisotropy can be absorbed into a spatial frame.
The next proposition derives the isometric comparison used downstream.

\begin{proposition}[Second-moment covariance derives orthogonal comparison]
\label{prop:metric-comparison}
Let $M_x,M_y\succ0$ be neighbouring second moments and let
$T_e:V_{\spat,x}\to V_{\spat,y}$ be an invertible record transport satisfying
\[
 T_eM_xT_e^{\mathsf T}=M_y.
\]
Then
\[
 O_e:=M_y^{-1/2}T_eM_x^{1/2}
\]
is orthogonal.  If the determinant orientation of the record system is
preserved, $O_e\in\SO(d)$.  Thus metric compatibility of neighbouring renewal
frames follows from covariance of their quadratic statistics.
\end{proposition}

\begin{proof}
Directly,
\[
 O_eO_e^{\mathsf T}
 =M_y^{-1/2}T_eM_xT_e^{\mathsf T}M_y^{-1/2}=I.
\]
Hence $O_e\in O(d)$.  Its determinant sign is the transported determinant-line
orientation, so orientation preservation gives determinant $+1$.
\end{proof}

The second-moment results supply the spatial coframe and metric-compatible
comparison transport.  They do not determine whether elementary translations
close around a plaquette.  That first-moment question is independent of metric
compatibility and is addressed by the exchange symmetry of two-step diamond
records.

\subsection{Exchange symmetry and translational plaquette closure}

The last upstream input required for the curved Cartan theorem is a
first-moment closure law.  It does not follow from associativity or ordinary
one-step detailed balance.  It follows from a stronger symmetry of one common
microscopic two-step law.

The guiding distinction is parity under exchange of the two path orders.  A
translational closure defect is exchange odd and must vanish in an
exchange-invariant stationary state; a rotational defect is exchange even and may
remain nonzero.  This is the finite precursor of the separation between torsion
and curvature.

Fix two elementary directions $i\ne j$.  Let $\Omega_{ij}$ be the finite space
of resolved microscopic diamond records, including the two ordered histories
$ij$ and $ji$, their terminal-location records, and any internal environment
labels retained by future tests.  Let
$\vartheta:\Omega_{ij}\to\Omega_{ij}$ exchange the two path orders.

\begin{definition}[Exchange-covariant diamond process]
\label{def:diamond-process}
A Markov kernel $K_{ij}$ on $\Omega_{ij}$ is exchange covariant when
\[
 K_{ij}(\vartheta\omega,\vartheta A)=K_{ij}(\omega,A)
\]
for every state $\omega$ and event $A$.  A translational defect
$d_{ij}^{(h)}:\Omega_{ij}\to V_{\spat}$ is exchange odd when
\[
 d_{ij}^{(h)}(\vartheta\omega)=-d_{ij}^{(h)}(\omega).
\]
\end{definition}

\begin{theorem}[Stationary exchange closure]
\label{thm:exchange-closure}
Assume $K_{ij}$ is primitive and exchange covariant.  Let $\pi_{ij}$ be its
unique stationary law.  Then $\pi_{ij}$ is exchange invariant and every
exchange-odd translational defect has zero stationary mean:
\[
 \int d_{ij}^{(h)}\,\dd\pi_{ij}=0.
\]
If the developed first-moment defect has the expansion
\[
 \int d_{ij}^{(h)}\,\dd\pi_{ij}
 =h^2T_{ij}+O(h^3),
\]
then $T_{ij}=O(h)$ and, in the continuum limit, $T_{ij}\to0$.  If the expansion
is exact through second order with an $h$-independent coefficient, then
$T_{ij}=0$ identically.  No exchange-even rotational defect is constrained.
\end{theorem}

\begin{proof}
Push the stationary law forward by $\vartheta$.  Exchange covariance gives
\[
 (\vartheta_*\pi_{ij})K_{ij}
 =\vartheta_*(\pi_{ij}K_{ij})
 =\vartheta_*\pi_{ij},
\]
so $\vartheta_*\pi_{ij}$ is stationary.  Primitivity gives uniqueness, hence
$\vartheta_*\pi_{ij}=\pi_{ij}$.  Therefore
\[
 \int d_{ij}^{(h)}\,\dd\pi_{ij}
 =\int d_{ij}^{(h)}\circ\vartheta\,\dd\pi_{ij}
 =-\int d_{ij}^{(h)}\,\dd\pi_{ij},
\]
and the integral vanishes.

Substituting the asymptotic expansion gives
$0=h^2T_{ij}+O(h^3)$, hence $T_{ij}=O(h)$.  If $T_{ij}$ is the fixed
second-order coefficient, division by $h^2$ and passage to $h\downarrow0$
gives $T_{ij}=0$.  The argument uses oddness and says nothing about an
even rotational observable.
\end{proof}

The stationary exchange theorem gives a direct symmetry argument for closure.
The same condition can also be characterized variationally.  Introducing a source
conjugate to the exchange-odd defect turns its expectation into the gradient of a
convex pressure, so zero source becomes the Euler--Lagrange point.

\begin{definition}[Exchange-source pressure]
\label{def:exchange-source-pressure}
Let $\Omega_{ij}$ be a finite resolved diamond space with exchange involution
$\vartheta$.  Let $S_h:\Omega_{ij}\to\R$ be exchange even and let
$D_h:\Omega_{ij}\to V_{\spat}$ be exchange odd:
\[
 S_h\circ\vartheta=S_h,
 \qquad D_h\circ\vartheta=-D_h.
\]
For $\lambda\in V_{\spat}^*$ define
\[
 \Psi_h(\lambda)
 :=\log\sum_{\omega\in\Omega_{ij}}
 \exp\!\left(-S_h(\omega)+\langle\lambda,D_h(\omega)\rangle\right),
\]
and let $\pi_{h,\lambda}$ be the associated Gibbs law.
\end{definition}

\begin{theorem}[Zero exchange source is the torsion-free stationary point]
\label{thm:exchange-pressure-stationarity}
For the pressure of \Cref{def:exchange-source-pressure},
\[
 \nabla\Psi_h(\lambda)
 =\mathbb E_{\pi_{h,\lambda}}D_h,
 \qquad
 \nabla^2\Psi_h(\lambda)
 =\operatorname{Cov}_{\pi_{h,\lambda}}(D_h)\succeq0.
\]
Moreover $\Psi_h(-\lambda)=\Psi_h(\lambda)$.  Hence $\lambda=0$ is a global
minimum and
\[
 \nabla\Psi_h(0)=0
 \quad\Longleftrightarrow\quad
 \mathbb E_{\pi_{h,0}}D_h=0.
\]
The minimum is unique modulo the annihilator of
\[
 \Span\{D_h(\omega)-D_h(\omega'):\omega,\omega'\in\Omega_{ij}\}.
\]
When $D_h$ is the developed translational diamond defect, zero source is
therefore exactly the Euler--Lagrange condition for first-moment plaquette
closure.
\end{theorem}

\begin{proof}
Because the sum is finite, differentiation under the sum is legitimate.  The
first derivative of the log partition function is the Gibbs expectation of
$D_h$, and the second derivative is its covariance bilinear form.  The latter is
positive semidefinite, so $\Psi_h$ is convex.

The change of variables $\omega\mapsto\vartheta\omega$, together with the
parities of $S_h$ and $D_h$, gives
\[
 \Psi_h(-\lambda)
 =\log\sum_\omega
 e^{-S_h(\vartheta\omega)+\langle\lambda,D_h(\vartheta\omega)\rangle}
 =\Psi_h(\lambda).
\]
Thus the convex function is even and has vanishing gradient at the origin,
which is consequently a global minimum.  Its Hessian vanishes on a covector
$\lambda$ exactly when $\langle\lambda,D_h\rangle$ is constant on the support;
these covectors form the stated annihilator.  Outside that subspace the pressure
is strictly convex.  Finally, taking $D_h=d_{ij}^{(h)}$ identifies the gradient
with the stationary translational first moment, so the zero-source stationarity
equation is precisely plaquette closure.
\end{proof}

\begin{remark}[Why ordinary detailed balance is insufficient]
A reversible one-step edge law may coexist with a stationary distribution that
selects one of two exchange-related diamond orders.  Such a phase has nonzero
exchange-odd first moment and therefore torsion.  The exact upstream input is
not generic stationarity or detailed balance but selection of the
exchange-invariant, zero-source translational phase.  This condition is
falsifiable and admits torsional counterphases.
\end{remark}

Exchange covariance therefore supplies two equivalent descriptions of the
torsion-free phase: invariance of the stationary diamond law and stationarity at
zero exchange source.  What remains to be shown is that weaker conditions, such
as ordinary stationarity or one-step detailed balance, do not already force this
choice.  The following countermodel establishes that sharp boundary.

\subsection{A sharp torsion phase countermodel}
\label{subsec:torsion-countermodel}

The exchange theorem gives exact closure in an exchange-invariant stationary
state.  The next model proves that ordinary stationarity and one-step detailed
balance do not force that state to be selected.

The model is intentionally elementary.  Its Ising order parameter records the
choice between the two diamond orders, while an independent isotropic direction
record keeps the spatial second moment nondegenerate.  This isolates torsion
selection from both reversibility and metric regularity.

\begin{construction}[Local torsion-record heat-bath field]
\label{constr:torsion-ising}
On $\Z^2$ let $\tau_x\in\{\pm1\}$ and consider the nearest-neighbour zero-field
Ising specification with coupling $\theta\ge0$.  Its single-site conditional law
is
\[
 q_x(s\mid\tau)
 =\frac{\exp\!\left(\theta s\sum_{y\sim x}\tau_y\right)}
        {2\cosh\!\left(\theta\sum_{y\sim x}\tau_y\right)}.
\]
Let global sign reversal $\vartheta\tau=-\tau$ exchange the two orders of one
microscopic diamond.  Fix $v\ne0$ and $\kappa>0$, and define the exchange-odd
translational defect
\[
 d_{ij}^{(h)}(\tau)=\kappa h^2\tau_0v.
\]
For every stationary Gibbs state $\mu$,
\[
 \Pi_{\mathrm{tr}}\mathcal P_{ij}^{(h)}(\mu)
 =\kappa h^2\mu(\tau_0)v.
\]
Independently attach to each update a deck-even direction record taking the six
values $\{\pm e_1,\pm e_2,\pm e_3\}$ with equal probability.  Its second moment
is $I_3/3$, so the model has a fixed nondegenerate spatial second moment in every
stationary phase without changing the torsion defect.
\end{construction}

\begin{theorem}[Torsion-free and torsional stationary phases coexist]
\label{thm:torsion-phase-coexistence}
The heat-bath law of \Cref{constr:torsion-ising} is local, finite range, and
satisfies one-step detailed balance in every Gibbs state.  Moreover:
\begin{enumerate}[label=\textup{(\roman*)},leftmargin=2.4em]
\item if $4\tanh\theta<1$, the infinite-volume Gibbs state is unique,
      exchange invariant, and
      $\Pi_{\mathrm{tr}}\mathcal P_{ij}^{(h)}=0$;
\item if $\theta\ge\tfrac12\log12$, there are two extremal deck-related Gibbs
      states $\mu_\theta^\pm$ satisfying
      \[
       \pm\mu_\theta^\pm(\tau_0)\ge\frac{203}{216},
      \]
      and hence
      \[
       \pm\,v\cdot
       \Pi_{\mathrm{tr}}\mathcal P_{ij}^{(h)}(\mu_\theta^\pm)
       \ge \kappa h^2\frac{203}{216}\norm{v}^2>0;
      \]
\item the symmetric mixture
      $\overline\mu_\theta=(\mu_\theta^++\mu_\theta^-)/2$ is stationary and
      detailed balanced at the same coupling but has zero translational defect.
\end{enumerate}
Thus one microscopic detailed-balanced law admits both torsion-free and
torsional stationary phases.
\end{theorem}

\begin{proof}
Let $c_{xy}$ be the maximal total-variation change of the conditional law at
$x$ when only the neighbouring spin at $y$ is changed.  The plus probability in
local field $h$ is $(1+\tanh(\theta h))/2$; changing one neighbour changes $h$
by two, and the maximal difference occurs between fields $-1$ and $1$.  Hence
$c_{xy}\le\tanh\theta$ for $x\sim y$ and $c_{xy}=0$ otherwise.  Thus
$\alpha:=\sup_x\sum_yc_{xy}\le4\tanh\theta$.  In a finite box, update the same
random site in two heat-bath chains and use a maximal coupling of their
conditional Bernoulli laws.  If $D$ is their Hamming distance, then
\[
 \mathbb E[D'\mid\tau,\tau']
 \le \left(1-\frac{1-\alpha}{|\Lambda|}\right)D.
\]
When $\alpha<1$, iteration and exhaustion of the lattice show that the influence
of arbitrary boundary conditions on every fixed cylinder event tends to zero;
there is therefore a unique infinite-volume Gibbs state (the standard
Dobrushin argument; see \cite{Georgii2011}).  The unique state is invariant under
$\tau\mapsto-\tau$, so its one-point function and the displayed defect vanish.

For the low-temperature statement, impose plus boundary conditions in a finite
box.  If $\tau_0=-1$, the minus cluster containing the origin has an outer dual
contour $\gamma$.  Flipping all spins inside $\gamma$ changes every crossing
bond from disagreement to agreement and increases the Gibbs weight by
$e^{2\theta|\gamma|}$.  The number $N_n$ of simple dual contours of length $n$
that surround the origin is bounded by $4n3^{n-1}$.  Hence, with
$r=3e^{-2\theta}$,
\[
 \mu_{\theta}^{+}(\tau_0=-1)
 \le \sum_{n\ge4}4n3^{n-1}e^{-2\theta n}
 =\frac43\frac{r^4(4-3r)}{(1-r)^2}.
\]
If $\theta\ge\tfrac12\log12$, then $r\le1/4$ and the right-hand side is at most
$13/432$.  Therefore
\[
 \mu_\theta^+(\tau_0)
 =1-2\mu_\theta^+(\tau_0=-1)
 \ge\frac{203}{216}.
\]
Monotone boundary limits give an extremal plus state; its global flip is the
extremal minus state with opposite magnetization.  Every Gibbs state is
stationary and detailed balanced for the heat-bath dynamics because the update
uses its exact single-site conditional distribution.  Substitution in the
defect formula proves (ii).  Convex combinations preserve stationarity and
detailed balance, while the opposite magnetizations cancel in the symmetric
mixture, proving (iii).
\end{proof}

\begin{corollary}[Torsion freedom is a genuine selection principle]
\label{cor:torsion-selection}
No theorem using only local finite-range UCP or Markov dynamics, stationarity,
one-step detailed balance, microscopic exchange covariance, bounded terminal
records, and a nondegenerate spatial second moment can imply translational
plaquette closure in every stationary phase.  Closure is selected, for example,
by uniqueness plus exchange covariance, restriction to exchange-invariant
states, or zero exchange source.
\end{corollary}

\begin{proof}
The two extremal low-temperature states in
\Cref{thm:torsion-phase-coexistence} satisfy all the listed properties and have
nonzero $O(h^2)$ translational defect, while their symmetric mixture satisfies
the same microscopic laws and has zero defect.
\end{proof}

The coexistence theorem proves that torsion freedom is a genuine phase selection:
the same detailed-balanced microscopic law can support symmetric torsion-free and
exchange-broken torsional stationary states.  We can now summarize the entire
dependency chain without hiding model choices or continuum assumptions inside the
foundational reconstruction.

\subsection{Logical hierarchy and the upstream-to-downstream bridge}
\label{subsec:logical-hierarchy}

Every premise of a theorem is a hypothesis, but not every hypothesis is a
foundational axiom.  We use the following distinctions.

This vocabulary is needed to read the final theorems correctly.  The first theorem
collects structures reconstructed uniformly from the operational package.  The
phase assumption then lists additional checkable properties of a chosen model.
The second theorem derives the finite signed modular Clifford data in that phase,
and the third invokes the separate analytic continuum theorem only after those
finite inputs have been established.

\begin{definition}[Axioms, models, phases, and selections]
\label{def:logical-roles}
\begin{enumerate}[label=\textup{(\roman*)},leftmargin=2.4em]
\item A \emph{foundational axiom} is imposed uniformly on every admissible
      operational theory; here these are the clauses of
      \Cref{ass:reconstruction-package} together with the specified finite
      domain and ancillary access.
\item A \emph{microscopic model specification} chooses a particular resolved
      graph, branch instrument, capacities, depths, reversible actions,
      direction records, diamond law, and parameter values inside that theory.
\item A \emph{phase criterion} is a checkable predicate on a specified model or
      parameter region, such as nontrivial determinant holonomy,
      supercriticality, projective rank, or ellipticity.
\item A \emph{physical selection} chooses one recurrent component, extremal
      state, boundary condition, zero-source state, or continuum scaling net
      when the model admits several possibilities.
\end{enumerate}
\end{definition}

\begin{theorem}[I. Structural reconstruction]
\label{thm:structural-reconstruction}
Assume the finite operational reconstruction package and complete ancillary
stability.  Then every admissible model canonically determines:
\begin{enumerate}[label=\textup{(\roman*)},leftmargin=2.4em]
\item the rooted future quotient $\Qq_+$ and its unique reduced predictive
      renewal memory;
\item the complete monoidal process quotient $\Qq_{\parallel}$, equal to the
      ordinary quotient precisely in the parallel-complete theory;
\item a finite oriented complex predictive $C^*$-algebra and a finite alphabet
      of UCP renewal maps;
\item on every selected recurrent support for which a primitive weighted
      transfer and channelwise stationary faithful state exist, the canonical
      KMS--Petz retrodictive family.
\end{enumerate}
No signed, supercritical, orientation-broken, elliptic, or torsion-free phase is
asserted by this theorem.
\end{theorem}

\begin{proof}
Items (i)--(ii) are
\Cref{prop:largest-congruence,thm:parallel-completion,thm:universal-quotient,thm:minimal-predictive-memory}.
The complex algebra and UCP maps are
\Cref{thm:complex-algebra,thm:operational-ucp}.  A primitive weighted transfer
selects the unique faithful state by
\Cref{thm:primitive-stationary-state}; channelwise stationarity and
\Cref{thm:petz-properties} give the common UCP retrodictive family.
\end{proof}

The structural theorem deliberately stops before any Lorentzian claim.  It
provides the common arena in which signed and unsigned, critical and subcritical,
reversible and dissipative models can all live.  The following criteria select the
particular recurrent phase whose finite data match the input requirements of the
companion construction.

\begin{assumption}[Checkable signed Lorentzian phase criteria]
\label{ass:lorentzian-phase}
A specified microscopic model $M$ has a selected irreducible recurrent block
satisfying:
\begin{enumerate}[label=\textup{(C\arabic*)},leftmargin=3em]
\item its resolved outcomes are either supplied operationally or extracted from
      an interaction-complete stable-pointer phase as in
      \Cref{thm:canonical-pointer-instrument}; their future-distinct record
      transport is determinant regular and has nonzero
      $[\chi]\in H^1(G;\Ztwo)$, with an admissible deck-odd test;
\item its realized reversible subgroup has, after kernel and radical reduction,
      a nonzero temporal row and the required spatial projective rank;
\item after the canonical successor refinement of
      \Cref{constr:successor-refinement}, its resolved branches are represented by
      their conditional predictive successors in a reduced future-tomographic
      realization, and the Petz-reindexed opposite branches satisfy the
      corresponding reverse-renewal condition; the resulting reciprocal scalar
      weights have positive depths and define a supercritical Perron block, so
      the unique pressure root $\beta>0$ exists;
\item it admits two deck-related extremal orientation states, and one of them is
      selected rather than their symmetric mixture;
\item its spatial second moment is positive definite and covariant between
      neighbouring cells;
\item its diamond state is exchange invariant, or equivalently selected at zero
      exchange source, and has the regularity required by the intended
      continuum theorem.
\end{enumerate}
\end{assumption}

\begin{theorem}[II. Signed modular Clifford phase]
\label{thm:lorentzian-phase}
Let $M$ be a microscopic model supplied by
\Cref{thm:structural-reconstruction} and satisfying
\Cref{ass:lorentzian-phase}.  Then $M$ canonically supplies:
\begin{enumerate}[label=\textup{(\roman*)},leftmargin=2.5em]
\item a real affinity class $[A]\in H^1(G;\R)$ and the unique factorization
      $q_e=c_e e^{A(e)/2}$ of each reciprocal capacity;
\item the independent determinant class $[\chi]\in H^1(G;\Ztwo)$, its principal
      signed cover, and the classified noncanonical $\Zfour$ amplitude lifts;
\item a projective reversible action whose radical quotient is a primitive
      complex Clifford factor, together with an intrinsic anti-Hermitian
      generator $\gamma^0$ satisfying $(\gamma^0)^2=-\one$;
\item the unique pressure-selected graded modular clock
      $\Delta_\beta=e^{\beta N}$;
\item a positive spatial coframe with orthogonal comparison transport;
\item exact stationary first-moment plaquette closure in the selected
      exchange-invariant phase, with no constraint on the exchange-even
      rotational defect.
\end{enumerate}
Thus $M$ lies in a signed modular Clifford--Lorentzian universality class.  The
finite criterion set is nonempty by
\Cref{cor:full-phase-nonempty}.
\end{theorem}

\begin{proof}
Conditional predictive unifilarity, canonical successor refinement, and future
tomography give branchwise reference renewal by
\Cref{prop:successor-refinement,thm:predictive-reference-renewal}; hence scalar
path weights are multiplicative by
\Cref{cor:predictive-reference-path-law,prop:multiplicative-path-weights}.
Branchwise Petz reversal gives reciprocal scalar fluxes by
\Cref{prop:branchwise-petz-reference-renewal,cor:petz-reciprocal-weights}; the
reverse-renewal clause of \Cref{ass:lorentzian-phase} makes the reindexed Petz
branches ray homogeneous.  Then
\Cref{prop:affinity-detailed-balance,thm:affinity-conductance} gives (i).
The canonical record quotient, uniqueness of parity, determinant line, signed
cover, amplitude lifts, and independence from affinity are
\Cref{thm:canonical-record-algebra,prop:record-transport-descent,thm:record-parity-unique,thm:noisy-record-orientation,thm:determinant-cover,prop:z4-lift,thm:independence,thm:no-affinity-parity-reduction}.
Projective lifting, temporal exchange, radical reduction, and the timelike
generator are
\Cref{thm:projective-lift,thm:temporal-row,thm:clifford-factor,cor:timelike-generator}.
Supercriticality gives the pressure root by \Cref{thm:pressure-root}; the
stationary law and its entropy--affinity--depth identity are
\Cref{constr:pressure-law,thm:entropy-affinity-depth}, and the graded clock and
its rigidity are \Cref{thm:pressure-path-clock,thm:affine-clock}.  Positive second moments and
covariance give (v) by \Cref{prop:metric-comparison}.  Finally (vi) is
\Cref{thm:exchange-closure,thm:exchange-pressure-stationarity}.  The deck-odd test retains the selected temporal
sign by \Cref{thm:deck-observability}.
\end{proof}

The second theorem is the finite culmination of this paper.  It assembles the
independent ingredients derived throughout Sections~
ef{sec:predictive}--

ef{sec:phase} into a signed modular Clifford--Lorentzian phase, while still
remaining entirely finite.  A smooth curved spacetime sector requires one final,
explicitly separate step: selection of a scaling net satisfying the analytic
hypotheses of the companion theorem.

\begin{theorem}[III. Realized curved sector, conditional on the companion theorem]
\label{thm:realized-curved-sector}
Let $M$ satisfy \Cref{thm:lorentzian-phase}.  Choose one recurrent component,
one deck-oriented extremal state, an exchange-invariant or zero-source diamond
state, and a continuum scaling net satisfying the analytic hypotheses of the
companion Lorentzian spectral-geometry theorem
\cite{PelissierLorentzian2026}.  Then that theorem applies to the finite datum
reconstructed here and yields the corresponding oriented torsion-free
Lorentzian Dirac continuum sector.  When the specified model admits additional
components or stationary phases, alternative selections may instead realize the
opposite temporal orientation, a symmetric unoriented mixture, or a torsional
sector; the existence of such alternatives is model dependent.
\end{theorem}

\begin{proof}
\Cref{thm:lorentzian-phase} supplies exactly the finite inputs used by the
companion theorem: UCP predictive memory, signed/Krein cover, primitive
projective Clifford sector, graded affine modular clock, positive
second-moment coframe, orthogonal comparison, and first-moment closure.  The
selected scaling net supplies the remaining analytic domain assumptions.
Opposite temporal sectors, when both exist, are related by deck reversal.
The countermodel of \Cref{thm:torsion-phase-coexistence} separately proves that
the structural framework contains microscopic laws for which different
stationary selections have zero or nonzero translational defect; it is not used
to assert this multiplicity for every model $M$.
\end{proof}

\begin{corollary}[Exact logical strength]
\label{cor:genuine-phases}
The paper does not prove that the operational axioms force every model to become
Lorentzian spacetime.  It proves:
\begin{enumerate}[label=\textup{(\alph*)},leftmargin=2.4em]
\item the axioms reconstruct a broad finite complex UCP predictive framework;
\item within that framework, explicit checkable criteria define a nonempty
      signed modular Clifford phase;
\item physical component, state, and scaling selections inside that phase allow
      the companion curved-sector theorem to be applied.
\end{enumerate}
The unsigned, subcritical, nonreversible, unoriented, spatially degenerate,
torsional, and nonsmooth counterphases remain admissible models of the same
structural theory.
\end{corollary}

The corollary is the practical reading guide for the paper's logical strength.
The operational axioms reconstruct the finite framework; phase criteria identify
a nonempty Lorentzian-type universality class; and additional physical and
analytic selections realize a particular continuum sector.  Counterphases are not
failures of the construction but part of its stated scope.

\begin{remark}[No circularity]
No Lorentzian metric, causal cone, signed cover, Clifford algebra, or physical
time coordinate is used to construct the rooted quotient, the complex algebra,
the UCP maps, the KMS dual, the record determinant, or the pressure root.  The
continuum theorem is invoked only after all finite signed modular data have been
reconstructed.  Conversely, the phase and state selections are not relabelled
as universal consequences.
\end{remark}

% =====================================================================
\section{Interpretation and outlook}
\label{sec:interpretation}
% =====================================================================

\subsection{From Predictive Relations to Temporal and Geometric Structure}
\label{subsec:intepretation}

\paragraph{Prediction as an operational primitive} The reconstruction begins below the level at which spacetime, quantum states,
or even a preferred notion of history have been specified.  Its primitive data
are composable operations, closed-experiment weights, and a family of
admissible future continuations and tests.  Two pasts are identified when no
such future experiment distinguishes them.  The resulting quotient therefore
retains neither a microscopic trajectory nor a complete record of what has
happened, but exactly the information that can still affect what happens next.
In this precise sense, the basic object of the theory is an effective
difference under continuation.

This viewpoint is related to the predictive equivalence underlying causal-state
constructions in computational mechanics
\cite{CrutchfieldYoung1989,ShaliziCrutchfield2001}, but the present setting is
both operational and compositional.  The objects being compared are generated
processes rather than sample paths of a fixed stochastic model, and the quotient
must remain compatible with admissible sequential and parallel contexts.
Moreover, outcome unifilarity is not imposed on a hidden-state presentation.
It appears only after successful histories have been quotiented by their
conditional future profiles.  In a future-tomographic operator realization,
that quotient-level determinism forces each resolved branch to return the
reference state of its unique predictive successor.  The scalar renewal law is
therefore a consequence of predictive reduction, rather than a modelling
assumption about microscopic dynamics.

The universality of the quotient gives a limited but mathematically definite
sense in which the retained predictive content is representation independent.
Every future-sufficient realization factors through it, whereas the choice of a
particular microscopic representative does not follow from the quotient alone.
Thus the construction does not assign ontological priority to one hidden
history.  It assigns priority to those distinctions between histories that
remain operationally effective.

The quantum operator structure enters only at a later stage.  Under the finite
sharp-purification and interference hypotheses used in the reconstruction
programme \cite{ChiribellaDAPS2009,ChiribellaScandolo2015,
BarnumLeeScandoloSelby2017,BarnumUdudecVdW2023}, the reduced predictive state
spaces acquire Euclidean Jordan structure.  The two-path phase law then selects
complex matrix type and an associative orientation, while compatibility with
arbitrary ancillary extensions upgrades positive deterministic dynamics to
UCP dynamics.  In this hierarchy, neither Hilbert space nor complete positivity
is primitive.  They are the finite operator form taken by a predictive theory
satisfying specific purification, interference, and compositional stability
conditions.  This is also why the complex selector and parallel completeness
must be kept logically distinct; tensor-product results such as
\cite{HancheOlsen1985} concern a different interface from the single-system
phase law.

\paragraph{Retrodiction, orientation, and temporal structure} Retrodiction adds a second structure without undoing the first.  Relative to a
faithful stationary reference state, the KMS--Petz map is the canonical adjoint
of forward prediction \cite{AccardiCecchini1982,Petz1986,Petz2003}.  It reverses
compositional order, but it is not generally an inverse channel and does not by
itself define a microscopic reversal of time.  Forward and retrodictive branch
weights may nevertheless fail to balance in the same path orientation.  Their
logarithmic ratio defines the real affinity class familiar from network
thermodynamics \cite{Schnakenberg1976,LebowitzSpohn1999,AndrieuxGaspard2007}.

The signed Lorentzian structure has a different origin.  It records whether a
frame of future-distinct records returns with reversed orientation after
transport around a recurrent cycle.  The corresponding determinant line is a
real local system, with its first Stiefel--Whitney class in
$H^1(G;\Ztwo)$ \cite{ZingerDeterminant2016}.  This parity is invisible to
positive scalar weights alone.  Conversely, a nonzero affinity need not alter
the determinant class.  The strict independence of these two invariants is
therefore more than a difference of notation: nonequilibrium circulation and
record orientation belong to different coefficient systems and classify
different operational effects.

The two sheets of the resulting cover should not be read as two universes, nor
as matter and antimatter sectors.  They are the minimal real carrier of an
orientation distinction forgotten by the unsigned predictive quotient.  The
positive theory records what the two orientations have in common; the signed
extension records how they differ.  When a reversible predictive cycle
exchanges the sheets, its implementer anticommutes with the fundamental
symmetry.  Their product then has negative square and supplies the intrinsic
timelike Clifford generator.  Temporal sign is therefore not introduced as an
additional positive coordinate.  It arises as an algebraic opposition between
two ways of carrying the same unsigned predictive content.

A physical arrow of time requires a further state selection.  A finite
primitive transfer that is symmetric under sheet exchange has a unique
sheet-symmetric stationary state, so it cannot select a nonzero deck-odd order
parameter.  An oriented phase must instead involve explicit symmetry breaking,
stationary-state degeneracy, boundary or superselection data, or a singular
large-system or large-depth limit.  The selected sign is operationally retained
only if a deck-odd future test remains admissible.  The reconstruction therefore
locates the logical ingredients from which an arrow could arise, but it does not
identify the thermodynamic or cosmological arrow with the existence of the
signed cover itself.

The pressure root supplies another temporal ingredient.  Its role is to turn
renewal depth into an intrinsic dimensionless modular scale.  The exact identity
\[
 \beta\mu_\ell
 =h_{\mathrm{rate}}+\mathbb E[\log c_e]+\frac12\sigma
\]
separates three contributions: predictive entropy, symmetric renewal capacity,
and nonequilibrium circulation.  The corresponding affine action is realized
on the graded path cover, where unbounded depth is separated from the finite
recurrent base.  This realizes, in a finite renewal setting, the relation
between modular flow and intrinsic time advocated in the thermal-time
programme \cite{ConnesRovelli1994}.  It does not, however, fix an absolute
second or energy unit.  Dimensionful calibration remains external to the
pressure theorem.

\paragraph{Spatial geometry and the status of torsion freedom} The spatial metric is likewise reconstructed from operational statistics rather
than imposed at the primitive level.  A positive second moment of stable
direction records defines a quadratic frame, and covariance of this moment
under transport forces neighbouring normalized frames to be compared
orthogonally.  Metric compatibility is therefore encoded by preservation of
quadratic predictive data.

The microscopic diamond law then separates translation from rotation.  The
translational defect is odd under exchange of the two path orders, whereas the
rotational defect is even.  An exchange-invariant stationary phase cancels the
former but leaves the latter unconstrained.  In continuum language, this is the
required distinction between torsion and curvature
\cite{Hehl1976}: the Levi--Civita sector is selected by vanishing translational
closure defect, not by suppressing rotational holonomy.

The Ising countermodel shows why this selection cannot be hidden inside ordinary
stationarity or one-step detailed balance.  One and the same local reversible
law can admit a symmetric torsion-free stationary state and extremal
exchange-broken states with nonzero translational defect.  Torsion freedom is
therefore a genuine phase or state-selection condition.  This strengthens the
logical interpretation of the curved reconstruction: the theory does not define
torsional sectors away, but identifies the additional microscopic symmetry
that selects the torsion-free branch.

\paragraph{Scope of the reconstruction} The result is a conditional sufficiency theorem, not a uniqueness theorem for
fundamental physics.  It shows that a stated operational package together with
a stated recurrent phase supplies the finite signed, modular, Clifford, and
metric data required by the companion Lorentzian construction
\cite{PelissierLorentzian2026}.  It does not show that every operational theory
enters that phase, that the phase is dynamically inevitable, or that its
continuum limit is unique.

Several limitations are structural rather than merely technical.  An averaged
channel does not select its resolved instrument without additional interaction
or monitoring data.  A nontrivial signed cover does not select one of its
orientations.  A positive pressure root fixes a dimensionless exponent but not
a laboratory unit.  A nondegenerate record moment does not by itself derive
three spatial dimensions, and finite signed modular data do not guarantee a
smooth Dirac continuum.  The unsigned, subcritical, nonreversible, spatially
degenerate, unoriented, and torsional regimes remain consistent alternatives.
Their presence is important: it makes each phase condition falsifiable rather
than definitional.

The conceptual claim should therefore be read narrowly.  What lies beneath the
geometric starting object of the companion paper need not be another geometry.
It can be a compositional system of distinctions whose future consequences can
be compressed, compared with their retrodictive adjoints, transported with or
without orientation reversal, and organized by recurrent statistical laws.  In
the selected phase studied here, these operations produce the finite algebraic
precursors of Lorentzian propagation.

\subsection{Outlook and testable interfaces}
\label{subsec:outlook}

The intermediate constructions have uses that do not depend on taking a
continuum spacetime limit.  They suggest several concrete directions in which
the framework can be tested, refined, or applied.

\paragraph{Predictive reconstruction from experimental data.}
The conditional quotient and the tomographic frame estimate turn predictive
model reduction into a quantitative inference problem.  Experimental branch
data can be grouped by approximate equality of conditional future statistics,
and the resulting successor states can be tested against the accumulated
trace-norm error bound.  This extends causal-state and quantum-memory
reconstruction \cite{ShaliziCrutchfield2001,GuWiesnerRieperVedral2012,
BinderThompsonGu2018} from model fitting toward certification of a minimal
predictive instrument.  A natural next step is a finite-sample theory that
combines statistical confidence regions with the deterministic perturbation
bounds proved here.

\paragraph{Independent control of affinity and record orientation.}
The cohomological independence of $[A]\in H^1(G;\R)$ and
$[\chi]\in H^1(G;\Ztwo)$ yields a direct experimental prediction: continuous
changes in nonequilibrium driving can alter entropy production without changing
the record-orientation sector, while a determinant crossing can flip the
orientation class without requiring a singularity in the scalar affinity.
Engineered Markov, photonic, or measurement-feedback networks provide natural
settings in which edge fluxes and record-transport determinants could be
measured simultaneously.  Such experiments would test the separation between
thermodynamic and signed information at the finite-network level.

\paragraph{Pointer and record engineering.}
The stable-pointer construction is canonical relative to a specified
interaction family.  Read in reverse, it provides a design principle: choose
system--environment couplings whose interaction algebra has a commutant with
the desired stable centre.  This connects the reconstruction to einselection,
decoherence-free subsystems, and redundant record formation
\cite{Zurek1981,Zurek2003,Schlosshauer2004,BlanchardOlkiewicz2003,
Evans1977,FrigerioVerri1982,FagnolaSassoUmanita2019,TicozziViola2008,
BlumeKohoutZurek2006}.  Determinant regularity then quantifies the noise margin
within which the orientation carried by those records survives.

\paragraph{Thermal quantum dynamics.}
The KMS--Petz branch reversal and cancellation current are directly relevant to
KMS-detailed-balanced quantum Markov dynamics.  In particular, the
entropy--affinity--depth identity gives an exact stationarity diagnostic for
resolved implementations of thermal samplers
\cite{FagnolaUmanita2010,ChenKastoryanoGilyen2023}.  Deviations from the identity
would distinguish imperfect sampling, incorrect branch resolution, and genuine
nonstationarity.  It would be useful to formulate corresponding finite-time and
finite-sample inequalities.

\paragraph{The geometric continuation.}
For the spacetime programme, the next mathematical tasks are more sharply
localized by the present reconstruction.  One must classify microscopic models
that enter the required nontrivial signed and projective phase, derive rather
than postulate the relevant spatial rank when possible, control the scaling of
the exchange-closure error, and prove the analytic compactness and
strong-resolvent hypotheses needed by the companion continuum theorem.  A
separate metrological bridge is also required to convert the dimensionless
modular parameter into physical units.  These are downstream questions; they
no longer concern the existence or canonical meaning of the finite renewal
object itself.

Taken together, the results suggest a research programme rather than a closed
ontology.  Prediction supplies the canonical reduction, retrodiction supplies
a stationary comparison, stable records supply orientation, recurrence supplies
projective and modular structure, and coarse-grained record statistics supply
metric and connection data.  The principal open question is how broadly these
structures arise in microscopic models without being imposed.  The present
paper makes that question mathematically accessible by isolating the exact
interfaces at which each additional structure enters and the counterphases in
which it fails.


% =====================================================================
\appendix

\section{Finite-dimensional operator-system facts}
\label{app:operator-systems}

This appendix records the elementary finite-dimensional facts used in
\Cref{sec:predictive}.

\begin{lemma}[Choi criterion in the present setting]
\label{lem:choi}
Let $\Phi:M_n(\C)\to M_m(\C)$ be linear.  Define its Choi matrix
\[
 C_\Phi:=\sum_{i,j=1}^{n}E_{ij}\otimes\Phi(E_{ij})
 \in M_n(\C)\otimes M_m(\C).
\]
Then $\Phi$ is completely positive if and only if $C_\Phi\ge0$.
\end{lemma}

\begin{proof}
If $\Phi$ is completely positive, then
\[
 C_\Phi=(\id_n\otimes\Phi)(|\Omega\rangle\langle\Omega|)\ge0,
 \qquad
 |\Omega\rangle=\sum_i e_i\otimes e_i.
\]
Conversely suppose $C_\Phi\ge0$ and write
$C_\Phi=\sum_a|v_a\rangle\langle v_a|$.  Reshape each
$v_a\in\C^n\otimes\C^m$ into a matrix $V_a:\C^n\to\C^m$.  Comparing matrix
entries gives
\[
 \Phi(x)=\sum_aV_axV_a^*.
\]
This Kraus form is completely positive.
\end{proof}

\begin{lemma}[Unitality and stationarity]
\label{lem:unital-stationary}
For a linear map $\Phi:\A\to\A$ on a finite-dimensional $C^*$-algebra, the
trace adjoint satisfies
\[
 \Phi(\one)=\one
 \quad\Longleftrightarrow\quad
 \Tr(\Phi_*(x))=\Tr(x)
\]
for all $x$.  More generally,
$\Phi_*(\rho)=\rho$ if and only if the state
$\varphi_\rho(a)=\Tr(\rho a)$ is invariant under $\Phi$.
\end{lemma}

\begin{proof}
By definition,
\[
 \Tr(\Phi_*(x))=\Tr(\Phi_*(x)\one)=\Tr(x\Phi(\one)).
\]
This equals $\Tr(x)$ for all $x$ exactly when $\Phi(\one)=\one$.  The second
statement is
\[
 \varphi_\rho(\Phi(a))
 =\Tr(\rho\Phi(a))
 =\Tr(\Phi_*(\rho)a),
\]
which equals $\varphi_\rho(a)$ for all $a$ exactly when
$\Phi_*(\rho)=\rho$.
\end{proof}

\section{Determinant walls and parity}
\label{app:determinant-walls}

The orientation class is locally constant only while the record transport is
invertible.  At rank-loss walls its continuation requires path data.

\begin{proposition}[Transverse determinant crossing]
\label{prop:det-crossing}
Let $L(t)$ be a $C^1$ path of real $r\times r$ matrices with
$\det L(t_0)=0$ and
\[
 \frac{\dd}{\dd t}\det L(t)\bigg|_{t=t_0}\ne0.
\]
Then $\sgn\det L(t)$ changes sign across $t_0$.  The parity of a path with
finitely many transverse crossings is
\[
 (-1)^{N_{\mathrm{cross}}},
\]
independent of sufficiently small transverse perturbations that fix the
endpoints.
\end{proposition}

\begin{proof}
A nonzero derivative makes $\det L(t)$ a scalar function with a simple zero,
so it has opposite signs on the two sides of $t_0$.  Under a sufficiently
small endpoint-fixed transverse perturbation, crossings can be created or
annihilated only in pairs; therefore their number modulo two is invariant.
This is the finite-dimensional parity principle of
\cite{FitzpatrickPejsachowicz1991,DollSchulzBaldesWaterstraat2019}.
\end{proof}

\begin{proposition}[First nonzero determinant jet]
\label{prop:det-jet}
Let $L(t)$ be $C^k$ near $t_0$ and suppose
\[
 \det L(t_0)=\cdots=\frac{\dd^{m-1}}{\dd t^{m-1}}\det L(t_0)=0,
 \qquad
 \frac{\dd^m}{\dd t^m}\det L(t_0)\ne0.
\]
Then the determinant sign changes across $t_0$ if and only if $m$ is odd.
\end{proposition}

\begin{proof}
Taylor's theorem gives
\[
 \det L(t)=c(t-t_0)^m+o(|t-t_0|^m),
 \qquad c\ne0.
\]
For sufficiently small $|t-t_0|$, the leading term controls the sign.  It has
opposite signs on the two sides exactly when $m$ is odd.
\end{proof}

\begin{definition}[Endpoint path parity]
\label{def:determinant-path-parity}
Let \(V,W\) be real vector spaces of the same finite dimension and let
\[
 Q:[a,b]\longrightarrow\Hom(V,W)
\]
be continuous with \(Q_a,Q_b\) invertible.  Choose arbitrary orientations of
\(V\) and \(W\) and define
\[
 \operatorname{par}(Q)
 :=
 \sgn\det Q_b\,\sgn\det Q_a
 \in\{\pm1\}.
\]
The product is independent of the two orientation choices.
\end{definition}

\begin{theorem}[Perturbative parity classification]
\label{thm:determinant-path-parity}
The parity of \Cref{def:determinant-path-parity} has the following properties.
\begin{enumerate}[label=\textup{(\roman*)},leftmargin=2.4em]
\item it is invariant under homotopies through paths with invertible
      endpoints;
\item it is multiplicative under path concatenation;
\item every \(C^1\) path admits an arbitrarily small endpoint-fixed
      perturbation \(\widetilde Q\) with finitely many regular crossings, and
      for every such perturbation
      \[
       \boxed{
       \operatorname{par}(Q)
       =
       (-1)^{\sum_{t_*}\dim\ker\widetilde Q_{t_*}};}
      \]
\item for an isolated infinitely flat crossing, parity gives a canonical
      relation between the signs of the two regular sides even though no
      finite determinant jet does;
\item for a nonisolated singular interval, parity canonically relates the two
      regular endpoint phases but does not canonically decompose the interval
      into individual signed walls.
\end{enumerate}
Thus endpoint-fixed perturbative parity is the maximal path-level
mod-two datum available without additional local regularity.
\end{theorem}

\begin{proof}
Changing the orientation of either \(V\) or \(W\) reverses both endpoint
determinant signs, leaving their product unchanged.  During an
endpoint-invertible homotopy, neither endpoint determinant can cross zero, so
the product is constant.  Under concatenation, the determinant sign at the
common endpoint appears twice and cancels.

For the crossing formula, let \(t_*\) be a regular crossing of a \(C^1\)
path.  Write
\[
 K=\ker Q_{t_*},\qquad C=\operatorname{coker}Q_{t_*},
 \qquad m=\dim K=\dim C.
\]
The derivative induces the crossing map
\[
 \Gamma_{t_*}:K\longrightarrow C,\qquad
 \Gamma_{t_*}=\pi_C\dot Q_{t_*}|_K,
\]
and regularity means that \(\Gamma_{t_*}\) is invertible.  Choose splittings
\(V=K\oplus V_1\) and \(W=C\oplus W_1\) such that
\(Q_{t_*}|_{V_1}:V_1\to W_1\) is an isomorphism.  The Schur-complement
expansion then gives, in oriented coordinates,
\[
 \det Q_t
 =
 c\,(t-t_*)^m+o(|t-t_*|^m),
 \qquad c\ne0.
\]
Hence passage through the crossing changes the determinant sign by
\((-1)^m\).

The singular matrices form a stratified algebraic hypersurface whose
top stratum consists of corank-one matrices.  Finite-dimensional relative
transversality gives an arbitrarily small endpoint-fixed perturbation meeting
the regular strata transversely and in finitely many points; higher-corank
coincidences may be separated by a further arbitrarily small perturbation.
Multiplying the local factors \((-1)^{\dim\ker}\) gives the ratio of endpoint
signs and hence the displayed formula.  Its independence of the chosen
perturbation follows from the already proved homotopy invariance.  This is the
finite-dimensional parity principle; see
\cite{FitzpatrickPejsachowicz1991,DollSchulzBaldesWaterstraat2019}.

For flat or nonisolated singular sets, the endpoint expression remains
defined whenever the endpoints are invertible.  It supplies only their total
relative parity; without a distinguished local crossing or nonzero jet there
is no canonical finer decomposition.
\end{proof}

\begin{remark}[Sharp singular boundary]
At a singular map itself, no nonzero determinant transport exists.  A
pointwise sign can therefore be continued only after adding local jet data,
a regular crossing structure, a path parity, or an independent orientation.
Rank loss is a genuine boundary of the nondegenerate signed phase rather than
a removable coordinate singularity.
\end{remark}

\section{Approximate exchange covariance}
\label{app:approx-exchange}

The exact closure theorem has a quantitative perturbative version.

\begin{theorem}[Mixing-gap stability of exchange closure]
\label{thm:approx-exchange}
Represent probability laws as row vectors.  Let $K$ be a primitive stochastic
matrix on a finite diamond space, let $P$ be the permutation matrix of the
exchange involution, and let $\pi K=\pi$.  Assume
\[
 \norm{PK-KP}_{1\to1}\le\eps
\]
and that every zero-mass row vector $\mu$ satisfies
\[
 \norm{\mu K}_1\le(1-\lambda)\norm{\mu}_1
\]
for some $\lambda>0$.  Then
\[
 \norm{\pi P-\pi}_1\le\frac{\eps}{\lambda}.
\]
For every bounded exchange-odd column vector $d$, meaning $Pd=-d$,
\[
 |\pi d|\le\frac{\eps}{2\lambda}\norm{d}_\infty.
\]
\end{theorem}

\begin{proof}
Set $\mu=\pi P-\pi$, which has total mass zero.  Using $\pi K=\pi$,
\[
 \mu K-\mu
 =\pi PK-\pi K-\pi P+\pi
 =\pi(PK-KP).
\]
Hence $\norm{\mu K-\mu}_1\le\eps$.  The reverse triangle inequality and the
mixing-gap assumption give
\[
 \eps\ge\norm{\mu-\mu K}_1
 \ge\norm{\mu}_1-\norm{\mu K}_1
 \ge\lambda\norm{\mu}_1,
\]
which proves the first estimate.  Since $Pd=-d$,
\[
 \mu d=\pi Pd-\pi d=-2\pi d.
\]
Therefore
\[
 2|\pi d|=|\mu d|
 \le\norm{\mu}_1\norm{d}_\infty
 \le\frac{\eps}{\lambda}\norm{d}_\infty.
\]
\end{proof}

\section{A worked finite signed renewal model}
\label{app:worked-model}

We close with a minimal example displaying the independence of the positive,
retrodictive, signed, and projective structures.

\begin{example}[Three-cycle with a nontrivial signed cover]
\label{ex:three-cycle}
Let $G$ be the directed three-cycle with both orientations on every edge.  Take
$\A=M_2(\C)$ and the faithful state $\rho=\one/2$.  Let the three forward
renewal channels be unitary conjugations
\[
 \Phi_1=\Ad(\sigma_x),
 \qquad
 \Phi_2=\Ad(\sigma_z),
 \qquad
 \Phi_3=\id,
\]
with the corresponding reverse channels equal to themselves.  Then every
channel is its KMS--Petz dual because the Hilbert--Schmidt and symmetric KMS
pairings coincide at the tracial state.

Assign positive branch weights with forward/backward ratios
$e^{a_1},e^{a_2},e^{a_3}$.  The cycle affinity is
$a_1+a_2+a_3$ and may be chosen independently of the operator channels.  Let
the stable record fibre be $R_x=\R^2$ at every vertex, with two orientation
preserving transports and one reflection.  The product determinant around the
cycle is negative, so $[\chi]\ne0$ and the orientation cover is connected,
regardless of whether the cycle affinity vanishes.

The automorphisms $\Ad(\sigma_x)$ and $\Ad(\sigma_z)$ commute even though their
implementers anticommute.  Their projective commutator is $-\one$, so the
nondegenerate spatial revision algebra is $M_2(\C)$.  An odd sign-holonomy
cycle then supplies the canonical temporal exchange relation with the sheet
operator $J$.  A full $3+1$ primitive block requires one additional independent
spatial revision label, as in the rank-four construction of the companion
paper.
\end{example}

\begin{proof}[Verification]
At $\rho=\one/2$, $\GammaR$ is a scalar multiple of the identity map and the
trace adjoint of a unitary conjugation is its inverse.  Pauli conjugations are
involutive, so $\Phi_j^\sharp=\Phi_j$.  The affinity statement is
\Cref{prop:affinity-detailed-balance}.  The reflection contributes one negative
determinant and therefore one nontrivial cycle sign.  Finally
$\sigma_x\sigma_z=-\sigma_z\sigma_x$, while their central sign disappears under
conjugation; this is exactly the projective mechanism of
\Cref{thm:projective-lift}.
\end{proof}

\section{A common-origin local UCP--Clifford phase}
\label{app:common-origin-model}

The preceding bridge is conditional on explicit phase criteria.  This appendix
constructs one finite-range microscopic instrument satisfying them jointly.
The construction is model specific; its purpose is to prove nonemptiness and
compatibility, not universal uniqueness.

Let $\Lambda$ be a finite periodic box in $\Z^2$ and
$\Omega_\Lambda=\{\pm1\}^{\Lambda}$.  For $i\in\Lambda$ and $s\in\{\pm1\}$,
write $\eta^{i,s}$ for the configuration obtained by replacing $\eta_i$ by
$s$.  Fix $\theta\ge0$ and define the heat-bath probability
\[
 q_i(s\mid\eta)
 =\frac{\exp\!\left(\theta s\sum_{j\sim i}\eta_j\right)}
        {2\cosh\!\left(\theta\sum_{j\sim i}\eta_j\right)}.
\]
Choose site weights $\nu_i>0$ with $\sum_i\nu_i=1$.  Let the spatial record
rank be $d=3$, choose $\lambda_{i,a}>0$, $b_a\in[-1,1]$, and
$\varepsilon\in\R$ with
\[
 \delta:=|\varepsilon|\max_a|b_a|<1.
\]
Set
\[
 r_{i,a}(u)
 :=\frac{\lambda_{i,a}(1+\varepsilon b_a u)}
         {\sum_{c=1}^3\lambda_{i,c}(1+\varepsilon b_c u)},
 \qquad u\in\{\pm1\}.
\]
The argument $u=\eta_is$ couples the direction record to whether the local
orientation is preserved or flipped.

Let $S=\C^4$ carry Hermitian Clifford generators
$\Gamma_0,\Gamma_1,\Gamma_2,\Gamma_3$ satisfying
\[
 \Gamma_\mu\Gamma_\nu+\Gamma_\nu\Gamma_\mu
 =2\delta_{\mu\nu}\one.
\]
Put $J=\Gamma_0$, $R_a=\Gamma_a$, and let
$\tau_S=\Tr/4$.  Choose $\kappa_b>0$ for $b\in\{0,1,2\}$ with
$\sum_b\kappa_b=1$, and define UCP maps
\[
 \Psi_{a,0}=\Ad(R_a),\qquad
 \Psi_{a,1}=\Ad(JR_a),\qquad
 \Psi_{a,2}(X)=\tau_S(X)\one.
\]
On
$\widehat\A_\Lambda=C(\Omega_\Lambda)\otimes M_4(\C)$ define the resolved
Heisenberg branch
\begin{equation}
\label{eq:common-origin-branch}
 (\widehat K_{i,s,a,b}F)(\eta)
 =\nu_iq_i(s\mid\eta)r_{i,a}(\eta_is)\kappa_b
   \Psi_{a,b}\!\left(F(\eta^{i,s})\right),
\end{equation}
and set
$\widehat K_\Lambda=\sum_{i,s,a,b}\widehat K_{i,s,a,b}$.

\begin{proposition}[UCP law, marginal, and deck covariance]
\label{prop:common-origin-ucp}
Every resolved branch in \eqref{eq:common-origin-branch} is CP and
$\widehat K_\Lambda$ is UCP.  On scalar observables $F=f\otimes\one$,
\[
 \widehat K_\Lambda(f\otimes\one)
 =(K_\Lambda f)\otimes\one,
\]
where $K_\Lambda$ is the ordinary random-scan Ising heat-bath channel.  Under
deck reversal $(\Theta F)(\eta)=F(-\eta)$, the record transforms as
$(s,a,b)\mapsto(-s,a,b)$ and
\[
 \Theta\widehat K_{i,s,a,b}
 =\widehat K_{i,-s,a,b}\Theta.
\]
\end{proposition}

\begin{proof}
Each $\Psi_{a,b}$ is UCP.  Pullback along
$\eta\mapsto\eta^{i,s}$, multiplication by a nonnegative scalar coefficient,
and composition with a CP map preserve complete positivity.  Since
$\sum_s q_i(s\mid\eta)=1$, $\sum_a r_{i,a}(u)=1$, and
$\sum_b\kappa_b=1$, the sum maps $\one$ to $\one$.  Every internal map fixes
scalar matrices, giving the classical marginal.  Finally
$q_i(-s\mid-\eta)=q_i(s\mid\eta)$ and
$(-\eta_i)(-s)=\eta_is$, so the direction weight is unchanged and deck
covariance follows.
\end{proof}

Let
\[
 \mu_{\Lambda,\theta}(\eta)
 =Z^{-1}\exp\!\left(\theta\sum_{\langle i,j\rangle}\eta_i\eta_j\right)
\]
and define the faithful state
\[
 \widehat\mu_{\Lambda,\theta}(F)
 =\sum_\eta\mu_{\Lambda,\theta}(\eta)\tau_S(F(\eta)).
\]

\begin{theorem}[Resolved transition balance and primitivity]
\label{thm:common-origin-balance}
For every configuration $\eta$, site $i$, redraw value $s$, direction $a$, and
internal mode $b$, put $\xi=\eta^{i,s}$ and define the scalar transition weight
\[
 w_{i,s,a,b}(\eta)
 :=\nu_iq_i(s\mid\eta)r_{i,a}(\eta_is)\kappa_b.
\]
Then the reverse elementary transition $\xi\to\eta$ has fixed redraw value
$\eta_i$ and satisfies
\[
 \boxed{
 \mu_{\Lambda,\theta}(\eta)w_{i,s,a,b}(\eta)
 =\mu_{\Lambda,\theta}(\xi)
   w_{i,\eta_i,a,b}(\xi).}
\]
Each internal map $\Psi_{a,b}$ is self-adjoint for the normalized
Hilbert--Schmidt pairing.  Consequently the complete channel
$\widehat K_\Lambda$ is self-adjoint for the symmetric KMS pairing of
$\widehat\mu_{\Lambda,\theta}$ and hence satisfies KMS detailed balance.  If
$\kappa_2>0$, its Schr\"odinger predual is primitive; therefore
$\widehat\mu_{\Lambda,\theta}$ is the unique faithful finite-volume stationary
state and the dynamics mixes exponentially.
\end{theorem}

\begin{proof}
The heat-bath identity gives
\[
 \mu(\eta)q_i(s\mid\eta)
 =\mu(\xi)q_i(\eta_i\mid\xi).
\]
Moreover $\xi_i\eta_i=s\eta_i=\eta_is$, so the direction factor is identical
on the two orientations of the same elementary transition.  Multiplication by
$\nu_i\kappa_b$ proves the boxed scalar identity.

For $b=0$, $R_a$ is a Hermitian involution.  For $b=1$, $JR_a$ is unitary and
$(JR_a)^{-1}=-(JR_a)$; the central sign disappears under adjoint action, so
$\Ad(JR_a)$ is an involution.  Both inner automorphisms are therefore
self-adjoint on $L^2(M_4(\C),\tau_S)$.  The depolarizing map is the orthogonal
projection onto $\C\one$ and is self-adjoint as well.  Expanding the KMS inner
product of $F$ with $\widehat K_\Lambda G$, changing variables
$\eta\leftrightarrow\xi$ transition by transition, and applying the scalar and
internal self-adjointness identities gives
\[
 \langle F,\widehat K_\Lambda G\rangle_{\widehat\mu}
 =\langle\widehat K_\Lambda F,G\rangle_{\widehat\mu}.
\]
Thus the full channel is its KMS adjoint.

For primitivity, choose a classical block carrying nonzero positive mass.  A
single transition with $b=2$ replaces its internal density by a strictly
positive multiple of $\one$.  All heat-bath, direction, and internal-mode
probabilities are strictly positive.  A finite sequence of single-site redraws
reaches every target classical configuration, while unitary conjugations
preserve strict positivity.  Since the state space is finite, one common power
sends every nonzero positive block state to a strictly positive density in every
block.  The finite-dimensional Perron theorem for primitive CP maps gives
uniqueness and exponential mixing.
\end{proof}

\begin{proposition}[Pressure, sign, and spatial moment]
\label{prop:common-origin-pressure-frame}
Fix an activity $\mathfrak a>1$ and define the depth-one pressure transfer
\[
 \widehat{\mathcal B}_z
 =\mathfrak a e^{-z}\widehat K_\Lambda.
\]
Then its unique pressure zero is
\[
 \beta=\log\mathfrak a,
\]
and its Doob-normalized channel is $\widehat K_\Lambda$.  Let
$\widehat G_\Lambda$ be the positive-support graph of the resolved finite-volume
transitions and quotient it by the free global deck involution
$\eta\mapsto-\eta$.  Primitivity makes $\widehat G_\Lambda$ connected, so
\[
 \widehat G_\Lambda\longrightarrow G_\Lambda:=\widehat G_\Lambda/\langle\Theta\rangle
\]
is a connected, hence nontrivial, principal $\Ztwo$-cover.  If the full resolved
record $(s,a,b)$ is immediately readable, its deck action is the product of
$9$ disjoint transpositions and therefore has determinant $-1$ on the record
contrast module.  The associated determinant character represents the nonzero
classifying this cover.

Attach vectors $v_{i,a}\in\R^3$ and define
\[
 M_i(\eta)
 =\sum_{s,a}q_i(s\mid\eta)r_{i,a}(\eta_is)
   v_{i,a}v_{i,a}^{\mathsf T}.
\]
Assume
\[
 0<\lambda_-\le\lambda_{i,a}\le\lambda_+,
 \qquad
 \sum_{a=1}^3v_{i,a}v_{i,a}^{\mathsf T}\succeq g_*I_3.
\]
Then, with
\[
 r_*:=\frac{\lambda_-(1-\delta)}{3\lambda_+(1+\delta)},
\]
one has
\[
 M_i(\eta)\succeq r_*g_*I_3,
 \qquad M_i(-\eta)=M_i(\eta).
\]
Thus deck-related orientation phases have the same positive rank-three spatial
second moment.
\end{proposition}

\begin{proof}
A unital finite-dimensional channel has spectral radius one, so
$\rho(\widehat{\mathcal B}_z)=\mathfrak a e^{-z}$ and the root is
$\log\mathfrak a$.  Every heat-bath coefficient is positive, hence the support
graph is connected; the global flip has no fixed configuration and acts freely.
A trivial double cover of a connected graph has disconnected total space, so the
connected quotient cover is nontrivial.  Deck reversal fixes $(a,b)$ and
exchanges the two values of $s$; there are $3\cdot3=9$ transpositions.  The
constant line has determinant $+1$, so the contrast module has the same
determinant $-1$, realizing the nontrivial cover character.

For every $a,u$,
\[
 r_{i,a}(u)\ge
 \frac{\lambda_-(1-\delta)}{3\lambda_+(1+\delta)}=r_*.
\]
Hence for $\xi\in\R^3$,
\[
 \xi^{\mathsf T}M_i(\eta)\xi
 \ge r_*\sum_a(\xi\cdot v_{i,a})^2
 \ge r_*g_*\norm{\xi}^2.
\]
Deck reversal sends $(\eta,s)$ to $(-\eta,-s)$ while preserving both
$q_i(s\mid\eta)$ and $r_{i,a}(\eta_is)$, so relabelling the sum proves deck
evenness.
\end{proof}

\begin{theorem}[Clifford factor generated by coherent branches]
\label{thm:common-origin-clifford}
Conditioning on $b=0,1$ leaves the predictive units
\[
 \alpha_a=\Ad(\Gamma_a),
 \qquad
 \beta_a=\Ad(\Gamma_0\Gamma_a).
\]
They generate $\Ad(\Gamma_0)$ because $\beta_a\alpha_a=\Ad(\Gamma_0)$.  Let
$H=\F_2^4$ with basis $e_0,e_1,e_2,e_3$ and choose implementing monomials
$R_x=\Gamma_0^{x_0}\Gamma_1^{x_1}\Gamma_2^{x_2}\Gamma_3^{x_3}$.  Their
commutator form has matrix
\[
 \Omega_{\mu\nu}=
 \begin{cases}0,&\mu=\nu,\\1,&\mu\ne\nu,
 \end{cases}
\]
over $\F_2$.  The form is nondegenerate and
\[
 \C_\sigma[H]\cong\Cl_4(\C)\cong M_4(\C).
\]
\end{theorem}

\begin{proof}
The coherent branches contain $\Ad(\Gamma_a)$ and
$\Ad(\Gamma_0\Gamma_a)$, so their products contain every generator.  Reordering
Clifford monomials produces only a scalar sign, defining a projective cocycle;
distinct generators anticommute and equal generators commute, giving $\Omega$.
If $\Omega x=0$ and $S=\sum_jx_j$, the $i$th equation is $S+x_i=0$, so every
$x_i=S$.  Summing gives $S=4S=0$ over $\F_2$, hence $x=0$.  Thus the alternating
form is nondegenerate.  The $16$ Clifford monomials are linearly independent
and span the $16$-dimensional algebra $M_4(\C)$, proving the isomorphisms.
\end{proof}

\begin{theorem}[Unified common-origin phase]
\label{thm:common-origin-phase}
Assume
\[
 \theta\ge\frac12\log12,
 \qquad \mathfrak a>1,
 \qquad \kappa_b>0,
\]
and the uniform spatial nondegeneracy hypotheses above.  Then the single
resolved branch family \eqref{eq:common-origin-branch} simultaneously has:
\begin{enumerate}[label=\textup{(\roman*)},leftmargin=2.4em]
\item primitive finite-volume UCP dynamics and a unique faithful stationary
      state;
\item exact resolved transition balance and full-channel KMS detailed balance;
\item in the thermodynamic limit, two extremal low-temperature deck-related orientation phases, separated by the deck-odd local spin test;
\item pressure root $\beta=\log\mathfrak a$;
\item a nontrivial connected signed cover and determinant character;
\item a uniformly positive rank-three deck-even spatial second moment;
\item a primitive projective revision factor
      $\Cl_4(\C)\cong M_4(\C)$ generated by the coherent resolved branches.
\end{enumerate}
Therefore the finite algebraic signed modular Clifford criteria---through
the coframe stage---are jointly realizable by one local microscopic edge
instrument.
\end{theorem}

\begin{proof}
Items (i)--(ii) are \Cref{prop:common-origin-ucp,thm:common-origin-balance}.
Summing over direction and internal records gives the two-dimensional Ising
heat-bath marginal, whose plus and minus boundary limits are distinct by the
Peierls estimate in the proof of \Cref{thm:torsion-phase-coexistence}; this gives
(iii).  Item (iv), the sign, and the spatial moment are
\Cref{prop:common-origin-pressure-frame}.  The final item is
\Cref{thm:common-origin-clifford}.  Every property belongs to the same branch
family, rather than to commuting factors appended after the fact.
\end{proof}

\begin{corollary}[Nonemptiness of the finite phase criteria]
\label{cor:full-phase-nonempty}
Choose the common-origin model with coefficients independent of $i$, so its
phase-averaged spatial moment is translation invariant and neighbouring
comparison may be taken to be the identity.  Form its local product with the
two-state diamond kernel
\[
 \Omega_{\diamond}=\{+,-\},\qquad
 K_{\diamond}(\zeta,\zeta')=\frac12,
 \qquad \vartheta_{\diamond}(\zeta)=-\zeta,
\]
and any bounded exchange-odd defect
$d_{\diamond}(-\zeta)=-d_{\diamond}(\zeta)$.  The product kernel is primitive and
exchange covariant; its stationary diamond law is uniform and has zero mean
defect.  The pre-existing determinant line of the common-origin record sector
embeds as an invariant factor and retains its nonzero class.  Consequently all
finite criteria of \Cref{ass:lorentzian-phase}, including exchange-invariant
first-moment closure, are simultaneously satisfied by an explicit local product
model.  A smooth or quantitative continuum conclusion still requires the
separate scaling regularity stipulated in that assumption.
\end{corollary}

\begin{proof}
The common-origin factor supplies sign holonomy, the projective Clifford rank,
pressure supercriticality, broken temporal orientation, and a positive
rank-three moment by \Cref{thm:common-origin-phase}.  Site-independent
coefficients make the averaged moment constant, hence covariant under identity
comparison.  The displayed diamond kernel has strictly positive entries, so it
is primitive, and it commutes with the exchange involution.  Its unique
stationary law is $(1/2,1/2)$, which annihilates every exchange-odd defect.
Tensor products of UCP channels are UCP, tensor products of primitive finite
channels are primitive after a common power, and the determinant local system
from one invariant record factor is unchanged by adjoining an independent
factor.  Thus the combined finite model satisfies all stated criteria.
\end{proof}

\begin{remark}[Scope of the rank-three statement]
The equality between record determinant parity and spatial-rank parity is a
property of this common record alphabet, not a universal proof that every
operational theory has three spatial dimensions.  In this model class,
nontrivial determinant parity requires odd rank and a nonzero spatial bivector
requires rank at least two, so rank three is the first admissible value.
\end{remark}

\section{Logical dependency ledger}
\label{app:ledger}

\begin{longtable}{>{\raggedright\arraybackslash}p{0.25\textwidth}
>{\raggedright\arraybackslash}p{0.20\textwidth}
>{\raggedright\arraybackslash}p{0.46\textwidth}}
\toprule
\textbf{Datum} & \textbf{Logical status} & \textbf{Origin in this paper}\\
\midrule
\endfirsthead
\toprule
\textbf{Datum} & \textbf{Logical status} & \textbf{Origin in this paper}\\
\midrule
\endhead
Rooted future quotient & canonical reduced memory & Future-profile quotient and minimality, \Cref{prop:largest-congruence,thm:minimal-predictive-memory}; no natural raw-history selector is determined without extra structure, \Cref{prop:no-natural-selector}.\\
Approximate predictive compression & quantitatively stable under rank and transversality control & A near-idempotent, near-sufficient map lies near an exact minimal projection; \Cref{thm:predictive-compression-stability,cor:compressed-dynamics-stability}.\\
Complete process quotient & canonical monoidal completion & Largest monoidal congruence below ordinary equivalence, \Cref{thm:parallel-completion,thm:universal-quotient}.\\
Parallel completeness & primitive no-holism axiom & Equality $\Qq_{\parallel}\cong\Qq_1$; \Cref{def:parallel-completeness,rem:no-holism}.\\
Process tomography & interface condition & Identifies operational profiles with full accumulated channels; \Cref{prop:tomographic-channel-equality}.  A single rooted orbit is insufficient by \Cref{ex:rooted-orbit-not-tomographic}.\\
Finite Euclidean Jordan structure & derived from sharp-reversible realisation & \Cref{prop:reversible-realisation-sharp-purification,thm:operational-diagonalization,thm:operational-jordan}.\\
Pure sharpness boundary & irreducible effect-side input & Reversible dilation alone does not determine pure sharpness; \Cref{prop:dilation-does-not-imply-sharpness}.\\
Finite complex $C^*$-algebra & reconstructed from Jordan structure and two-path phases & \Cref{thm:complex-algebra}.\\
Complete positivity & ancillary-stability theorem & \Cref{thm:operational-ucp}.\\
Faithful stationary state & selected by a primitive weighted transfer & \Cref{thm:primitive-stationary-state}.\\
KMS--Petz reverse & canonical for stationary channels and reference-renewing resolved branches & \Cref{cor:common-petz-family,thm:petz-properties,prop:branchwise-petz-reference-renewal}.\\
Scalar path weights & derived on the conditional predictive orbit & Outcome-unifilar prediction, canonical successor refinement, and future tomography force reference renewal, \Cref{lem:conditional-successor,prop:successor-refinement,thm:predictive-reference-renewal,cor:predictive-reference-path-law}.\\
Instrument gauge boundary & Kraus-invariant within a fixed resolved branch & \Cref{prop:kraus-invariance,rem:refinement-not-splitting}.\\
Approximate scalar renewal & quantitative finite-history stability & Tomographic frame control and CP trace-norm contraction, \Cref{thm:approximate-reference-renewal}.\\
Petz path reversal & exact stationary adjointness identity & Reversed-path Petz comparison is identically balanced, forcing the same-orientation affinity convention; \Cref{prop:stationary-petz-path-reversal,rem:same-orientation-affinity}.\\
Real affinity and conductance & resolved reciprocal branch data modulo gauge & \Cref{prop:affinity-detailed-balance,thm:affinity-conductance}.\\
Stable raw instrument & conditional microscopic derivation & A specified interaction-complete monitoring phase selects the nondemolition commutant, its central pointer algebra, and a basis-independent resolved instrument; \Cref{prop:interaction-algebra-minimal,thm:stable-pointer-selection,thm:canonical-pointer-instrument}.\\
Canonical record algebra & conditional predictive quotient of resolved outcomes & \Cref{thm:canonical-record-algebra}.\\
Signed orientation & unique nontrivial one-dimensional record parity, robust under noisy regular transport & \Cref{thm:record-parity-unique,thm:noisy-record-orientation,prop:noisy-record-sign-stability,thm:determinant-cover}.\\
Singular orientation walls & jet and path-level parity classification & Finite-order crossings are controlled by determinant jets and general endpoint-invertible paths by perturbative parity; \Cref{prop:det-jet,thm:determinant-path-parity}.\\
Affinity--parity separation & no canonical coefficient reduction & \Cref{thm:independence,thm:no-affinity-parity-reduction}.\\
$\Zfour$ amplitude lifts & classified torsor, noncanonical & \Cref{prop:z4-lift}.\\
Projective multiplier & automatic on a finite factor & \Cref{thm:projective-lift}.\\
Primitive Clifford factor & canonical radical quotient & \Cref{thm:clifford-factor}.\\
Selected temporal sign & phase plus deck-odd observability & \Cref{thm:no-finite-orientation-selection,thm:deck-observability}.\\
Modular exponent $\beta$ & pressure selected in a supercritical recurrent phase & \Cref{thm:pressure-root}.\\
Entropy--affinity--depth law & exact identity for the stationary pressure process & \Cref{lem:stationary-affinity-cohomology,thm:entropy-affinity-depth}.\\
Affine modular clock & canonical on the graded path cover & \Cref{thm:pressure-path-clock,prop:finite-cycle-clock-obstruction}.\\
Orthogonal frame comparison & derived from second-moment covariance & \Cref{prop:metric-comparison}.\\
Translational closure & exchange-invariant phase or zero-source pressure stationarity & \Cref{thm:exchange-closure,thm:exchange-pressure-stationarity}; torsional counterphases in \Cref{thm:torsion-phase-coexistence}.\\
Nonemptiness of the finite signed phase & explicit local product model & \Cref{cor:full-phase-nonempty}.\\
Physical time unit & metrological calibration & Not fixed by the dimensionless pressure law.\\
Lorentzian continuum & downstream theorem & Companion work \cite{PelissierLorentzian2026}.\\
\bottomrule
\end{longtable}

% =====================================================================

\section{Necessity and no-go ledger}
\label{app:necessity-ledger}

The following table records why the principal assumptions and phase
conditions cannot be deleted wholesale.  It distinguishes a missing proof
from a genuine logical obstruction.

\begin{longtable}{>{\raggedright\arraybackslash}p{0.24\textwidth}
>{\raggedright\arraybackslash}p{0.27\textwidth}
>{\raggedright\arraybackslash}p{0.40\textwidth}}
\toprule
\textbf{Input or condition} & \textbf{Role} & \textbf{Sharp boundary or counterexample}\\
\midrule
\endfirsthead
\toprule
\textbf{Input or condition} & \textbf{Role} & \textbf{Sharp boundary or counterexample}\\
\midrule
\endhead
Arbitrary future continuations &
Makes rooted equivalence a continuation congruence &
Immediate tests alone can identify histories later separated by an allowed update; the future quantifier in \Cref{def:future-equivalence} is load-bearing.\\

Choice of representative or normal form &
Realizes the canonical quotient as a raw-history retract &
The quotient admits no natural invariant selector when a profile-preserving automorphism exchanges representatives; \Cref{prop:no-natural-selector}.\\

Parallel completeness &
Eliminates process distinctions visible only with ancillas &
The two-rebit construction gives processes ordinarily indistinguishable but ancilla-distinct; \Cref{prop:rebit-parallel-counterexample}.\\

Process tomography &
Identifies rooted operational profiles with full accumulated channels &
Two maps may agree on the entire available rooted orbit and differ on an unprepared direction; \Cref{ex:rooted-orbit-not-tomographic}.\\

Minimal spectral rank and transversality &
Turn approximate compression into a nearby minimal exact retract &
An exact sufficient projection may retain arbitrary redundant memory, while a correct-rank projection may erase the quotient; \Cref{rem:predictive-stability-boundaries}.\\

Pure-process reversal &
Supplies pure sharpness in the reconstruction package &
Reversible dilation and essential uniqueness leave the admitted effect theory undetermined; the uniformly noisy effect model in \Cref{prop:dilation-does-not-imply-sharpness} destroys pure sharpness.\\

Minimal continuous two-path interference &
Selects complex rather than real, quaternionic, spin-factor, or exceptional Jordan type &
Without \textup{(O4)}--\textup{(O5)}, the other finite Euclidean Jordan factors remain compatible; \Cref{thm:two-path-complex-face,cor:pairwise-complex-type}.\\

Complete ancillary stability &
Upgrades positive deterministic maps to UCP maps &
Single-system positivity alone permits positive non-CP transformations; complete positivity is precisely positivity at every matrix level, \Cref{thm:operational-ucp}.\\

Resolved instrument data &
Makes branch weights and record outcomes operationally meaningful &
An averaged channel has inequivalent Kraus/instrument decompositions.  A canonical instrument follows only relative to a stable pointer phase, \Cref{thm:canonical-pointer-instrument,rem:pointer-extraction-boundary}.\\

Interaction-complete monitoring &
Selects the maximal nondemolition pointer algebra &
Monitoring a proper subalgebra may leave spurious stable records; pure Hamiltonian motion need not converge to a pointer expectation, \Cref{thm:stable-pointer-selection,rem:pointer-extraction-boundary}.\\

Conditional predictive successor resolution &
Forces branchwise reference renewal and scalar path weights &
A primitive stationary averaged channel can have branches that do not preserve the reference ray; \Cref{ex:average-not-branchwise}.\\

Determinant regularity &
Defines signed transport of future-distinct record contrasts &
The sign is locally constant only off the rank-loss wall; at \(\det Q_e=0\) a contrast is erased and no pointwise orientation transport exists, \Cref{prop:noisy-record-sign-stability}.\\

Nonzero determinant class &
Produces the principal signed cover &
Positive scalar affinity does not produce a mod-two class, and no nonzero homomorphism \((\R,+)\to\Ztwo\) exists; \Cref{thm:independence,thm:no-affinity-parity-reduction}.\\

Sufficient reversible projective rank &
Produces a nontrivial Clifford factor &
A purely dissipative predictive memory may possess no nontrivial units; radical reduction can also collapse the projective form, \Cref{thm:clifford-factor}.\\

Pressure supercriticality &
Selects a positive modular exponent &
If \(r(0)\le1\), no positive pressure root exists; \Cref{thm:pressure-root}.\\

Graded path cover &
Carries a nonzero affine modular clock consistently &
A common nonzero modular exponent on a finite scalar-return cycle is impossible; \Cref{prop:finite-cycle-clock-obstruction}.\\

Orientation-state selection and deck-odd tests &
Produces and operationally retains one temporal sign &
A finite primitive deck-symmetric transfer has a unique deck-even stationary state, and deck-even tests identify the two sheets; \Cref{thm:no-finite-orientation-selection,thm:deck-observability}.\\

Positive definite spatial second moment &
Produces a nondegenerate spatial coframe &
A degenerate direction law has no invertible quadratic frame and therefore no Lorentzian spatial block; \Cref{prop:metric-comparison}.\\

Exchange-invariant or zero-source diamond phase &
Removes translational torsion while leaving rotation free &
One-step detailed balance and stationarity admit torsional extremal phases; \Cref{thm:torsion-phase-coexistence,cor:torsion-selection}.\\

Continuum scale separation and regularity &
Permits the companion strong-resolvent Dirac limit &
Finite signed modular data alone do not imply a smooth continuum; this remains an explicit domain condition in \Cref{thm:realized-curved-sector}.\\

Dimensionful calibration &
Converts the dimensionless modular parameter to laboratory units &
No dimensionless reconstruction can select an absolute second or joule; the pressure root fixes only an intrinsic dimensionless scale.\\
\bottomrule
\end{longtable}


\begin{thebibliography}{99}

\bibitem{PelissierLorentzian2026}
A.~P\'elissier,
\emph{Renewal Spectral Geometry and the Emergence of Lorentzian Spacetime},
companion manuscript (2026).

% ---- Operational reconstructions of quantum theory ----

\bibitem{Hardy2001}
L.~Hardy,
\emph{Quantum theory from five reasonable axioms},
arXiv:quant-ph/0101012 (2001).

\bibitem{ChiribellaDAPS2009}
G.~Chiribella, G.~M.~D'Ariano, and P.~Perinotti,
\emph{Probabilistic theories with purification},
Physical Review A \textbf{81} (2010), 062348.

\bibitem{ChiribellaDArianoPerinotti2011}
G.~Chiribella, G.~M.~D'Ariano, and P.~Perinotti,
\emph{Informational derivation of quantum theory},
Phys.\ Rev.\ A \textbf{84} (2011), 012311.

\bibitem{MasanesMuller2011}
L.~Masanes and M.~P.~M\"uller,
\emph{A derivation of quantum theory from physical requirements},
New J.\ Phys.\ \textbf{13} (2011), 063001.

\bibitem{DakicBrukner2011}
B.~Daki\'c and \v{C}.~Brukner,
\emph{Quantum theory and beyond: is entanglement special?},
in \emph{Deep Beauty: Understanding the Quantum World through
Mathematical Innovation}, ed.\ H.~Halvorson,
Cambridge University Press (2011), 365--392.

\bibitem{DArianoChiribellaPerinotti2017}
G.~M.~D'Ariano, G.~Chiribella, and P.~Perinotti,
\emph{Quantum Theory from First Principles: An Informational Approach},
Cambridge University Press, Cambridge (2017).

\bibitem{SelbyScandoloCoecke2021}
J.~H.~Selby, C.~M.~Scandolo, and B.~Coecke,
\emph{Reconstructing quantum theory from diagrammatic postulates},
Quantum \textbf{5} (2021), 445.

\bibitem{vandeWetering2019}
J.~van~de~Wetering,
\emph{An effect-theoretic reconstruction of quantum theory},
Compositionality \textbf{1} (2019), 1.

\bibitem{BarnumWilce2014}
H.~Barnum and A.~Wilce,
\emph{Local Tomography and the Jordan Structure of Quantum Theory},
Foundations of Physics \textbf{44} (2014), 192--212.

\bibitem{HardyWootters2012}
L.~Hardy and W.~K.~Wootters,
\emph{Limited holism and real-vector-space quantum theory},
Found.\ Phys.\ \textbf{42} (2012), 454--473.

\bibitem{delaTorre2012}
G.~de~la~Torre, L.~Masanes, A.~J.~Short, and M.~P.~M\"uller,
\emph{Deriving quantum theory from its local structure and reversibility},
Phys.\ Rev.\ Lett.\ \textbf{109} (2012), 090403.

\bibitem{ChiribellaScandolo2015}
G.~Chiribella and C.~M.~Scandolo,
\emph{Operational axioms for diagonalizing states},
Electronic Proceedings in Theoretical Computer Science \textbf{195} (2015),
96--115.

\bibitem{BarnumLeeScandoloSelby2017}
H.~Barnum, C.~M.~Lee, C.~M.~Scandolo, and J.~H.~Selby,
\emph{Ruling out Higher-Order Interference from Purity Principles},
Entropy \textbf{19} (2017), article 253.

\bibitem{BarnumUdudecVdW2023}
H.~Barnum, C.~Ududec, and J.~van de Wetering,
\emph{Self-duality and Jordan structure of quantum theory follow from
homogeneity and pure transitivity}, arXiv:2306.00362.

\bibitem{HancheOlsen1985}
H.~Hanche--Olsen,
\emph{JB-algebras with tensor products are $C^*$-algebras},
in \emph{Operator Algebras and their Connections with Topology and Ergodic
Theory}, Lecture Notes in Mathematics 1132, Springer, 1985, 223--229.

% (HancheOlsen1985 is already present in the manuscript; no new entry needed)

% ---- Computational mechanics and predictive states ----

\bibitem{CrutchfieldYoung1989}
J.~P.~Crutchfield and K.~Young,
\emph{Inferring statistical complexity},
Phys.\ Rev.\ Lett.\ \textbf{63} (1989), 105--108.

\bibitem{ShaliziCrutchfield2001}
C.~R.~Shalizi and J.~P.~Crutchfield,
\emph{Computational mechanics: pattern and prediction, structure and
simplicity},
J.~Stat.\ Phys.\ \textbf{104} (2001), 817--879.

\bibitem{GuWiesnerRieperVedral2012}
M.~Gu, K.~Wiesner, E.~Rieper, and V.~Vedral,
\emph{Quantum mechanics can reduce the complexity of classical models},
Nat.\ Commun.\ \textbf{3} (2012), 762.

\bibitem{BinderThompsonGu2018}
F.~C.~Binder, J.~Thompson, and M.~Gu,
\emph{Practical unitary simulator for non-Markovian complex processes},
Phys.\ Rev.\ Lett.\ \textbf{120} (2018), 240502.

% ---- Decoherence, pointer bases, quantum Markov semigroups ----

\bibitem{Zurek1981}
W.~H.~Zurek,
\emph{Pointer basis of quantum apparatus: into what mixture does the wave
packet collapse?},
Phys.\ Rev.\ D \textbf{24} (1981), 1516--1525.

\bibitem{Zurek2003}
W.~H.~Zurek,
\emph{Decoherence, einselection, and the quantum origins of the classical},
Rev.\ Mod.\ Phys.\ \textbf{75} (2003), 715--775.

\bibitem{Schlosshauer2004}
M.~Schlosshauer,
\emph{Decoherence, the measurement problem, and interpretations of quantum
mechanics},
Rev.\ Mod.\ Phys.\ \textbf{76} (2005), 1267--1305.

\bibitem{BlanchardOlkiewicz2003}
Ph.~Blanchard and R.~Olkiewicz,
\emph{Decoherence induced transition from quantum to classical dynamics},
Rev.\ Math.\ Phys.\ \textbf{15} (2003), 217--243.

\bibitem{Evans1977}
D.~E.~Evans,
\emph{Irreducible quantum dynamical semigroups},
Commun.\ Math.\ Phys.\ \textbf{54} (1977), 293--297.

\bibitem{FrigerioVerri1982}
A.~Frigerio and M.~Verri,
\emph{Long-time asymptotic properties of dynamical semigroups on
$W^*$-algebras},
Math.\ Z.\ \textbf{180} (1982), 275--286.

\bibitem{FagnolaSassoUmanita2019}
F.~Fagnola, E.~Sasso, and V.~Umanit\`a,
\emph{The role of the atomic decoherence-free subalgebra in the study of
quantum Markov semigroups},
J.~Math.\ Phys.\ \textbf{60} (2019), 072703.

\bibitem{TicozziViola2008}
F.~Ticozzi and L.~Viola,
\emph{Quantum Markovian subsystems: invariance, attractivity, and control},
IEEE Trans.\ Autom.\ Control \textbf{53} (2008), 2048--2063.

\bibitem{BlumeKohoutZurek2006}
R.~Blume-Kohout and W.~H.~Zurek,
\emph{Quantum Darwinism: entanglement, branches, and the emergent
classicality of redundantly stored quantum information},
Phys.\ Rev.\ A \textbf{73} (2006), 062310.

% ---- Retrodiction, Petz duals, quantum detailed balance ----

\bibitem{AccardiCecchini1982}
L.~Accardi and C.~Cecchini,
\emph{Conditional expectations in von Neumann algebras and a theorem of
Takesaki},
J.~Funct.\ Anal.\ \textbf{45} (1982), 245--273.

\bibitem{Petz1986}
D.~Petz,
\emph{Sufficient subalgebras and the relative entropy of states of a von
Neumann algebra},
Communications in Mathematical Physics \textbf{105} (1986), 123--131.

\bibitem{Petz2003}
D.~Petz,
\emph{Monotonicity of quantum relative entropy revisited},
Reviews in Mathematical Physics \textbf{15} (2003), 79--91.

\bibitem{Alicki1976}
R.~Alicki,
\emph{On the detailed balance condition for non-Hamiltonian systems},
Rep.\ Math.\ Phys.\ \textbf{10} (1976), 249--258.

\bibitem{KFGV1977}
A.~Kossakowski, A.~Frigerio, V.~Gorini, and M.~Verri,
\emph{Quantum detailed balance and KMS condition},
Commun.\ Math.\ Phys.\ \textbf{57} (1977), 97--110.

\bibitem{GoldsteinLindsay1995}
S.~Goldstein and J.~M.~Lindsay,
\emph{KMS-symmetric Markov semigroups},
Math.\ Z.\ \textbf{219} (1995), 591--608.

\bibitem{FagnolaUmanita2010}
F.~Fagnola and V.~Umanit\`a,
\emph{Generators of KMS symmetric Markov semigroups on $\mathcal B(h)$:
symmetry and quantum detailed balance},
Commun.\ Math.\ Phys.\ \textbf{298} (2010), 523--547.

\bibitem{Crooks2008}
G.~E.~Crooks,
\emph{Quantum operation time reversal},
Phys.\ Rev.\ A \textbf{77} (2008), 034101.

\bibitem{BuscemiScarani2021}
F.~Buscemi and V.~Scarani,
\emph{Fluctuation theorems from Bayesian retrodiction},
Phys.\ Rev.\ E \textbf{103} (2021), 052111.

% ---- Nonequilibrium network thermodynamics ----

\bibitem{Schnakenberg1976}
J.~Schnakenberg,
\emph{Network theory of microscopic and macroscopic behavior of master
equation systems},
Rev.\ Mod.\ Phys.\ \textbf{48} (1976), 571--585.

\bibitem{LebowitzSpohn1999}
J.~L.~Lebowitz and H.~Spohn,
\emph{A Gallavotti--Cohen-type symmetry in the large deviation functional
for stochastic dynamics},
J.~Stat.\ Phys.\ \textbf{95} (1999), 333--365.

\bibitem{AndrieuxGaspard2007}
D.~Andrieux and P.~Gaspard,
\emph{Fluctuation theorem for currents and Schnakenberg network theory},
J.~Stat.\ Phys.\ \textbf{127} (2007), 107--131.

\bibitem{ZingerDeterminant2016}
A.~Zinger,
\emph{The determinant line bundle for Fredholm operators: construction,
properties, and classification},
Mathematica Scandinavica \textbf{118} (2016), 203--268.

\bibitem{Bowen1975}
R.~Bowen,
\emph{Equilibrium States and the Ergodic Theory of Anosov Diffeomorphisms},
Lecture Notes in Mathematics \textbf{470}, Springer, Berlin (1975).

% ---- Thermodynamic formalism, modular/thermal time ----

\bibitem{Ruelle2004}
D.~Ruelle,
\emph{Thermodynamic Formalism}, 2nd ed.,
Cambridge University Press, Cambridge (2004).

\bibitem{ParryPollicott1990}
W.~Parry and M.~Pollicott,
\emph{Zeta functions and the periodic orbit structure of hyperbolic
dynamics},
Ast\'erisque \textbf{187--188} (1990).

\bibitem{ConnesRovelli1994}
A.~Connes and C.~Rovelli,
\emph{Von Neumann algebra automorphisms and time--thermodynamics relation in
generally covariant quantum theories},
Classical and Quantum Gravity \textbf{11} (1994), 2899--2917.

% ---- Torsion / Einstein--Cartan ----

\bibitem{Hehl1976}
F.~W.~Hehl, P.~von~der~Heyde, G.~D.~Kerlick, and J.~M.~Nester,
\emph{General relativity with spin and torsion: foundations and prospects},
Rev.\ Mod.\ Phys.\ \textbf{48} (1976), 393--416.

\bibitem{ChiribellaProcessTomography2021}
G.~Chiribella,
\emph{Process tomography in general physical theories},
arXiv:2109.12067 (2021).

\bibitem{JordanVonNeumannWigner1934}
P.~Jordan, J.~von Neumann, and E.~Wigner,
\emph{On an algebraic generalization of the quantum mechanical formalism},
Annals of Mathematics \textbf{35} (1934), 29--64.

\bibitem{BarnumHilgert2020}
H.~Barnum and J.~Hilgert,
\emph{Spectral Properties of Convex Bodies},
Journal of Lie Theory \textbf{30} (2020), 315--344.

\bibitem{AlfsenShultz1998}
E.~M.~Alfsen and F.~W.~Shultz,
\emph{Orientation in operator algebras},
Proceedings of the National Academy of Sciences \textbf{95} (1998), 6596--6601.

\bibitem{AlfsenShultz2003}
E.~M.~Alfsen and F.~W.~Shultz,
\emph{Dynamical Correspondences}, in \emph{Geometry of State Spaces of
Operator Algebras}, Birkh\"auser, 2003, 191--207.

\bibitem{Choi1975}
M.-D.~Choi,
\emph{Completely positive linear maps on complex matrices},
Linear Algebra and its Applications \textbf{10} (1975), 285--290.

\bibitem{Seneta2006}
E.~Seneta,
\emph{Non-negative Matrices and Markov Chains}, revised edition,
Springer, 2006.

\bibitem{Georgii2011}
H.-O.~Georgii,
\emph{Gibbs Measures and Phase Transitions}, second edition,
de Gruyter, 2011.

% ---- Downstream: Gibbs samplers ----

\bibitem{ChenKastoryanoGilyen2023}
C.-F.~Chen, M.~J.~Kastoryano, and A.~Gily\'en,
\emph{An efficient and exact noncommutative quantum Gibbs sampler},
arXiv:2311.09207 (2023).


\bibitem{FitzpatrickPejsachowicz1991}
P.~M.~Fitzpatrick and J.~Pejsachowicz,
\emph{Parity and generalized multiplicity},
Transactions of the American Mathematical Society \textbf{326} (1991),
281--305.

\bibitem{DollSchulzBaldesWaterstraat2019}
N.~Doll, H.~Schulz--Baldes, and N.~Waterstraat,
\emph{Parity as $\mathbb Z_2$-valued spectral flow},
Bulletin of the London Mathematical Society \textbf{51} (2019), 836--852.

\bibitem{JencovaPetz2006}
A.~Jen\v{c}ov\'a and D.~Petz,
\emph{Sufficiency in quantum statistical inference},
Commun.\ Math.\ Phys.\ \textbf{263} (2006), 259--276.

\bibitem{CDPNetworks2009}
G.~Chiribella, G.~M.~D'Ariano, and P.~Perinotti,
\emph{Theoretical framework for quantum networks},
Phys.\ Rev.\ A \textbf{80} (2009), 022339.

% ---- Instruments and operational quantum probability ----

\bibitem{DaviesLewis1970}
E.~B.~Davies and J.~T.~Lewis,
\emph{An operational approach to quantum probability},
Commun.\ Math.\ Phys.\ \textbf{17} (1970), 239--260.

\bibitem{Ozawa1984}
M.~Ozawa,
\emph{Quantum measuring processes of continuous observables},
J.~Math.\ Phys.\ \textbf{25} (1984), 79--87.

\bibitem{Stinespring1955}
W.~F.~Stinespring,
\emph{Positive functions on $C^*$-algebras},
Proc.\ Amer.\ Math.\ Soc.\ \textbf{6} (1955), 211--216.

\bibitem{Kadison1952}
R.~V.~Kadison,
\emph{A generalized Schwarz inequality and algebraic invariants for
operator algebras},
Ann.\ of Math.\ \textbf{56} (1952), 494--503.

\bibitem{Paulsen2002}
V.~Paulsen,
\emph{Completely Bounded Maps and Operator Algebras},
Cambridge University Press, Cambridge (2002).

\bibitem{EvansHoeghKrohn1978}
D.~E.~Evans and R.~H{\o}egh-Krohn,
\emph{Spectral properties of positive maps on $C^*$-algebras},
J.~London Math.\ Soc.\ (2) \textbf{17} (1978), 345--355.

\bibitem{GKS1976}
V.~Gorini, A.~Kossakowski, and E.~C.~G.~Sudarshan,
\emph{Completely positive dynamical semigroups of $N$-level systems},
J.~Math.\ Phys.\ \textbf{17} (1976), 821--825.

\bibitem{Lindblad1976}
G.~Lindblad,
\emph{On the generators of quantum dynamical semigroups},
Commun.\ Math.\ Phys.\ \textbf{48} (1976), 119--130.

\bibitem{BrattelliRobinson1997}
O.~Bratteli and D.~W.~Robinson,
\emph{Operator Algebras and Quantum Statistical Mechanics~2}, 2nd ed.,
Springer, Berlin (1997).

% ---- Statistical mechanics of lattice systems ----

\bibitem{FriedliVelenik2017}
S.~Friedli and Y.~Velenik,
\emph{Statistical Mechanics of Lattice Systems: A Concrete Mathematical
Introduction},
Cambridge University Press, Cambridge (2017).

% ---- Signed graphs, Krein spaces ----

\bibitem{Zaslavsky1982}
T.~Zaslavsky,
\emph{Signed graphs},
Discrete Appl.\ Math.\ \textbf{4} (1982), 47--74.

\bibitem{Bognar1974}
J.~Bogn\'ar,
\emph{Indefinite Inner Product Spaces},
Springer, Berlin (1974).

% ---- Interpretational context ----

\bibitem{Rovelli1996}
C.~Rovelli,
\emph{Relational quantum mechanics},
Int.\ J.\ Theor.\ Phys.\ \textbf{35} (1996), 1637--1678.

\bibitem{Wheeler1990}
J.~A.~Wheeler,
\emph{Information, physics, quantum: the search for links},
in \emph{Complexity, Entropy, and the Physics of Information},
ed.\ W.~H.~Zurek, Addison-Wesley, Redwood City (1990), 3--28.
\end{thebibliography}

\end{document}
